% 本文件由 LaTeX 排版 Agent 根据有效 translation.json 自动生成。
% author=Eusebius of Caesarea (c. 265-c. 340); work_id=2502; title=君士坦丁传

\chapter{\WorkTitle{君士坦丁传}}

% structure: p0001 source=h1 target=omitted reason=page-title-already-rendered-by-parent

\Emph{第一卷}\newline{}\Emph{第二卷}\newline{}\Emph{第三卷}\newline{}\Emph{第四卷}\par

\SourceRecord{Translated by Ernest Cushing Richardson.；Nicene and Post-Nicene Fathers, Second Series；Vol. 1.；Edited by Philip Schaff and Henry Wace.；Buffalo, NY: Christian Literature Publishing Co.,；1890.；<http://www.newadvent.org/fathers/2502.htm>.}

% structure: p0001 source=h1 target=section reason=page-title

\section{君士坦丁传（第一卷）}

% structure: p0002 source=h2 target=subsection reason=relative-heading-rank; base=h2

\subsection{第一章 序言——论君士坦丁之死。}

全体人类\Emph{早已}同心合意，以欢庆的盛宴庆祝这位伟大皇帝的第二和第三个十年统治的完成；我们已经亲自在天主众仆人的集会中迎接他，如同胜利的征服者，并在其统治二十周年时以应得的赞美颂扬他；而更近的，我们仿佛编织了话语的花环，在他自己的宫殿中，于其三十周年时环绕在他神圣的头上。\par

但如今，当我想要表达一些惯常的情感时，我困惑迟疑，不知转向何方，因我完全沉迷于眼前非凡的景象。因为无论我朝哪个方向看去，无论是向东、向西，或遍观世界，或仰望穹苍，处处时时我都看见那位有福者仍在执掌着同一帝国。在地上，我看见他的儿子们，如同他光辉的新反射者，四处散布他们父亲品格的辉光，而他本人依然活着且有力，比以前更完满地治理着人类的一切事务，因他通过其子女的延续而倍增。他们此前已拥有凯撒的尊位；但现在，因披戴着他的自我，并为他的功业所增色，因他们虔诚的卓越，被冠以元首、奥古斯都、可敬者、皇帝等称号。\par

% structure: p0005 source=h2 target=subsection reason=relative-heading-rank; base=h2

\subsection{第二章 序言（续）}

当我想到那位不久前还以肉身与我们同在可见的人，即使在死后，当自然思想摒弃一切多余之物为不合适时，他仍奇妙地拥有与先前同样的帝国居所、尊荣和赞美，确实惊异不已。但更进一步，当我提升思绪甚至到穹苍之拱，并在那里默观他三倍有福的灵魂与天主自己共融，脱去一切必死和尘世的衣袍，闪耀着光芒的荣耀之袍，并领悟到它不再与必死生命的短暂时期和事务相连，而是因一顶永不凋谢的花冠和永无止息、有福存在的永生而受尊崇时，我仿佛丧失了言语和思想的能力，无法说出一个字，反而谴责自己的软弱，强使自己沉默，我将恰如其分地赞美他的任务交给一位更有能力者，即那位不朽的天主和真实的圣言，唯独他能证实他自己的话语。\par

% structure: p0007 source=h2 target=subsection reason=relative-heading-rank; base=h2

\subsection{第三章 天主如何尊荣虔诚的君主，却摧毁暴君}

既然他保证那些荣耀和尊崇他的人将从他手中获得丰厚的赏报，而那些作为敌人和仇敌反对他的人将自取灵魂的灭亡，他已经通过他自己的宣告证实了这些真理，一方面显示了那些否认并反对他的暴君的可怕结局，同时表明他的仆人的死亡以及他的生命都值得赞美和钦佩，并且理应不仅通过可朽坏的纪念碑，而且通过不朽的纪念物来铭记。\par

人类为人类脆弱不定的寿命寻求某种安慰，认为通过建立纪念碑以不朽的尊荣来荣耀他们祖先的记忆。有些人运用绘画的生动描绘和色彩；有些人从无生命的木块雕刻雕像；而另一些人通过在碑石和纪念碑上深刻铭文，试图将他们所尊敬之人的美德传至永久的纪念。所有这些确实都是可朽坏的，被时间消磨，因为它们是可朽坏身体的形象，而非表达不朽灵魂的形像。但是，对于那些对此生终结之后没有稳固幸福盼望的人来说，这些似乎已足够。但是天主，我说那位天主，即万有的共同救主，为那些热爱虔诚的人保存了比人类思想所能领受的更伟大的福乐，甚至在此世赐予未来赏报的预兆和初果，以某种方式将不朽的希望启示给可朽的眼光。古代先知们通过圣经传给我们的神谕宣告了这一点；古代曾以各种美德照耀的虔诚之人的生活向后世见证了同样的事；我们自己的时代也证实了这一点真实，因为在其中君士坦丁，这位在所有曾执掌罗马权力者中唯一是万物主宰天主的友人，已向全人类显现为虔诚生活的如此清晰的榜样。\par

% structure: p0010 source=h2 target=subsection reason=relative-heading-rank; base=h2

\subsection{第四章 天主如何尊荣君士坦丁}

君士坦丁所崇拜的天主自己，通过他的旨意最清晰的显示确认了这个真理，在他统治的开始、过程和结束时都临在扶助他，并把他树立给人类作为虔诚的教训榜样。因此，藉着赐予他的多样祝福，天主使他成为我们所闻的一切统治者中独一无二的，既是巨大的光耀，又是真挚虔诚的最清晰声音的宣告者。\par

% structure: p0012 source=h2 target=subsection reason=relative-heading-rank; base=h2

\subsection{第五章 他统治三十余年，享寿六十余岁}

关于他统治的持续时期，天主以整整三个十年期并外加一些时日尊荣了他，将他整个必死生命的年限延长至这个数字的两倍。并且，天主乐意使他成为自己主权权能的代表，显示他为整个暴君族类的征服者，以及那些疯狂举起不敬之臂反对万有至高之王的地上敌天巨人的毁灭者。他们可以说是出现片刻，然后消失；而独一的真天主，当他使其仆人披上天国的甲胄，能够独自对抗众多敌人，并藉着他使人类摆脱众多不敬神者时，便立他为万国崇拜他的教师，以响亮的声音在众人面前见证他承认真天主，并厌恶那些非神的虚假之神的谬误。\par

% structure: p0014 source=h2 target=subsection reason=relative-heading-rank; base=h2

\subsection{第六章 他是天主的仆人，万国的征服者}

这样，他如同一个忠信善良的仆人行事并作证，公开宣告并承认自己是至高君王的顺从仆役。天主立即赏报他，使他成为统御者与主权者，并且得胜到如此地步，以致在所有统治者中，唯有他始终行进在征服之路上，未被征服、不可战胜；凭借他的战利品，他成为传统所载从未有过的更伟大的统治者。他是如此蒙天主所爱、如此蒙福；如此虔诚，如此在他所行的一切事上蒙福，以至于他极其轻易地获得了比任何先于他者更多的邦国的权柄，并且安然持守其权力，直至他生命的终结。\par

% structure: p0016 source=h2 target=subsection reason=relative-heading-rank; base=h2

\subsection{第七章　与波斯王\OriginalTerm{居鲁士}{Cyrus}及\OriginalTerm{马其顿}{Macedon}的\OriginalTerm{亚历山大}{Alexander}之比较。}

古代历史描述波斯王\OriginalTerm{居鲁士}{Cyrus}为直至当时最为卓越的君王。然而我们若看他一生的终结，便发现其与其昔日繁荣甚不相符，因为他死于妇人之手，死得卑贱而不光彩。\par

再者，\OriginalTerm{希腊}{Greece}诸子颂扬\OriginalTerm{马其顿}{Macedon}的\OriginalTerm{亚历山大}{Alexander}为征服众多不同邦国者；然而我们发现他早年夭折，未及壮年，死于宴饮醉乐之害。他一生不过三十二年，在位期间仅及其三分之一。他如雷霆无情，踏过血河，将整个民族与城市，老幼皆沦为奴隶。但当他刚达壮年，正在惋惜青春欢娱之失时，可怕的死亡击中了他，且为免他再肆毒害人类，将他绝于异国敌土，无嗣无继，无家可归。他的王国也立即分裂，每位将领各取一份据为己有。然而此人却因此等事迹被称颂。\par

% structure: p0019 source=h2 target=subsection reason=relative-heading-rank; base=h2

\subsection{第八章　他几乎征服了整个世界。}

但是我们的皇帝在马其顿人去世的年龄开始执政，却活了双倍的年岁，在位时间更长达三倍。他以温和节制的虔诚训诫教导军队，把战旗带到\OriginalTerm{不列颠人}{Britons}和居住在西大洋深处之邦国。他也征服了整个\OriginalTerm{斯基提亚}{Scythia}，虽然她位于最遥远的北方，分为无数多样而野蛮的部落。他甚至将征服推进到边界最南端的\OriginalTerm{布勒米人}{Blemmyans}和\OriginalTerm{埃塞俄比亚人}{Ethiopians}那里；他也不认为获取东方邦国不值得他关注。简言之，将他的圣光之辉散射到全世界各地，直至最遥远的\OriginalTerm{印度人}{Indians}，即居住于有人居住大地最外缘的邦国，他接受了所有蛮邦统治者、总督和总督的归顺，他们欣然欢迎并致敬，派遣使节和礼物，极其珍视他的相识与友谊；以至于在各自的国家以画像和雕像尊崇他，康斯坦丁在诸皇帝中独被所有人承认和庆祝。尽管如此，即使在那些遥远的邦国中，他也以皇谕放胆宣讲他天主的名号。\par

% structure: p0021 source=h2 target=subsection reason=relative-heading-rank; base=h2

\subsection{第九章　他为虔诚皇帝之子，并将权柄遗赠于诸皇儿。}

他不仅以言语作此见证，同时也身体力行每一条德性之路，丰富地结出虔诚的各类果实。他以慷慨的明证确保朋友的爱戴；因他以人道原则治理，使他的统治轻松可受，为所有阶层的臣民所接受。最终，在悠长岁月之后，当他因神圣的劳苦而疲惫时，他所敬拜的天主以不朽的赏报加冕他，将他从短暂的王国迁至他为圣徒灵魂所预备的永生，并为他兴起三个儿子继承他的权力。既然皇位由他的父亲传给他，同样，按照自然法则，这权柄为他的子女和后代保留，如同父业，延续至无穷世代。事实上，天主自己，这位当我们还在世时就以神圣尊荣分别这位蒙福君王的，并以亲手精选的祝福装饰他离世的天主，应当是他事迹的书写者；因为他已在天上纪念碑上记录了他的劳苦与成功。\par

% structure: p0023 source=h2 target=subsection reason=relative-heading-rank; base=h2

\subsection{第十章　论此史之必要及其训诲价值。}

然而，要恰当地言说这一蒙福品格虽属艰难，沉默虽是较为安稳且风险较小的做法，但我仍必须，若要逃避疏忽与怠惰的指控，犹如以言语描绘一幅肖像，作为对这位虔诚君王的纪念，效仿人类艺术之描摹。因为我若不尽己绵薄之力——尽管微弱且无甚价值——来赞美这位以如此卓越的虔诚尊崇天主者，我将自觉羞愧。我想我的作品在其他方面也必既具教益又属必要，因为它将包含那些取悦万有主宰天主的、王者而高贵的行动的记述。因为，难道不应当感到羞耻吗？\OriginalTerm{尼禄}{Nero}以及其他更恶的不虔不敬的暴君的回忆，都有勤勉的作家用优雅的语言美化他们无价值的事迹，并以卷帙浩繁的历史记录下来，而我却沉默不语，天主亲自赐予我如此一位历代史书未载的皇帝，并允许我来到他面前，享受他的相识与友爱。\par

因此，若说任何人都有责任，那必然是我的责任，要向那些能从崇高事迹的榜样中激发起对天主之爱的人，大力宣扬他的美德。因为有些人写卑鄙人物的生平，写那些无助于道德改善的行为史，出于私心——或爱或恨，有时甚至只是为了炫耀自己的学问——他们用精巧优雅的措辞，过分夸大了对本质上卑劣行为的描述。这样，对那些因天主的眷顾而与邪恶隔绝的人，他们成了邪恶的教师，而不是善的教师，而那些邪恶本应被遗忘和黑暗所掩盖。但我的叙述，尽管与它所必须描述的伟大事迹不相称，仍将因单纯叙述崇高事迹而获得光彩。确实，记载那些令天主喜悦的行为，对心地善良的人而言，将提供一种绝非无益，实则极有教益的研究。\par

% structure: p0026 source=h2 target=subsection reason=relative-heading-rank; base=h2

\subsection{第十一章  他当前的目的仅限于记录\Person{君士坦丁}虔诚的行为。}

因此，我打算略过这位三倍有福的君主的大部分帝王事迹；例如，他在战场上的冲突和交战、他的个人英勇、他对敌人的胜利和成功，以及他获得的许多凯旋；同样，他为个人利益所做的准备、为臣民的社会利益所制定的法律，以及许多其他仍鲜活在众人记忆中的帝国辛劳；我目前这项事业的设计，只在于谈论和书写那些与他的宗教品格相关的情况。\par

既然这些本身几乎无穷无尽，我将从我已知的事实中选择那些最合适、最值得永久记录的事，并尽可能简短地加以叙述。从此以后，确实有充分而自由的机会，可以以各种方式赞美这位真正有福的君主，而此前我们无法这样做，因为根据不确定的人生变迁，我们被禁止在任何人死亡之前断定他为有福。那么，我恳求天主的助佑，愿天上圣言的感召之助与我同在，让我从他生命的最早时期开始叙述我的历史。\par

% structure: p0029 source=h2 target=subsection reason=relative-heading-rank; base=h2

\subsection{第十二章  他如何像\Person{梅瑟}一样，在君王的宫殿中长大。}

古代历史记载说，一个残酷的暴君种族压迫了希伯来民族；天主仁慈地眷顾他们在苦难中的境况，使当时尚是婴孩的\Person{梅瑟}先知在压迫者的宫殿和怀抱中长大，并受教于他们所有的智慧。及至他长大成人，天主的正义要为受压迫的人民报仇的时刻到来，天主的先知服从更强大的主的旨意，离开了王室，在言语和行为上与养育他的暴君疏远，公开承认他的真正弟兄和亲属。于是，天主提拔他成为全民族的领袖，将希伯来人从敌人的奴役中解救出来，并藉着他向暴君种族施行了天主的惩罚。这个古老的故事，虽然被大多数人视为神话，却已传至众人耳中。但现在，同一位天主使我们亲眼目睹了比神话更奇妙的奇迹，且因其新近出现而比任何传说都更确实。因为当世的暴君竟敢与至高天主作战，并严重迫害了他的教会。在这其中，\Person{君士坦丁}——他即将成为这些暴君的毁灭者，但当时尚在幼年，脸上刚长出青春的绒毛——像那位天主的仆人一样，住在暴君的家中；但他虽然年轻，却没有分享不敬神者的生活方式：因为从那时起，他高贵的本性在圣神的引导下，便倾向于虔诚和令天主喜悦的生活。此外，效法他父亲榜样的愿望也激励着儿子走向德行。他的父亲是\Person{君士坦提乌斯}（我们此时应当重提他的事迹），是我们这时代最杰出的皇帝；有必要简要叙述他生平中几件关乎他儿子荣耀的事。\par

% structure: p0031 source=h2 target=subsection reason=relative-heading-rank; base=h2

\subsection{第十三章  论他的父亲\Person{君士坦提乌斯}，他拒绝效法\Person{戴克里先}、\Person{马克西米安}和\Person{马克森提乌斯}迫害基督徒的行为。}

当四位皇帝共同管理\Place{罗马帝国}时，唯有\Person{君士坦提乌斯}采取了一条与其同僚不同的行事方针，与至高天主建立了友谊。\par

因为他们围攻并摧毁天主的教会，将之夷为平地，甚至破坏了祈祷之所的根基，而他却使自己的双手免受他们可憎的不敬神行为的玷污，从未在任何方面与他们相似。他们以不分青红皂白地屠杀虔诚的男女来玷污自己的省份，而他却使自己的灵魂免于这种罪行的污点。他们陷入不敬神的偶像崇拜的迷宫，先使自己、然后使所有在他们权下的人成为邪恶魔鬼之错误的奴隶，而他同时在自己的整个统辖区内开创了最深的和平，并确保他的臣民不受阻碍地庆祝对天主的崇拜的权利。简言之，当他的同僚用最沉重的勒索压迫所有人，使他们的生活不堪忍受，甚至比死亡更糟时，唯有\Person{君士坦提乌斯}以温和宁静的统治治理他的人民，并对他的人民表现出真正父母般的关怀和养育。此人的其他美德数不胜数，为众人所称颂；我将记录其中一两个事例，作为那些我必须略过的美德品质的样本，然后继续按预定顺序叙述。\par

% structure: p0034 source=h2 target=subsection reason=relative-heading-rank; base=h2

\subsection{第十四章  他的父亲\Person{君士坦提乌斯}如何因贫穷被\Person{戴克里先}责备，之后装满国库，并将钱归还给捐献者。}

由于关于这位君王的种种传闻，描述他性格仁慈温和，虔敬程度极高，且说他因对臣民极度宽容，国库中甚至没有存钱；当时占据最高权力的皇帝派人申斥他忽视公益，同时指责他贫穷，并以国库空虚作为证据。对此，他请皇帝的使者稍待片刻，召集其统治下各族最富有的臣民，告知他们他需要金钱，并说现在是他们自愿表达对君王爱戴的时候。\par

他们一听到这话（似乎早已渴望有机会表明诚心善意），便热切踊跃地用金银和其他财富填满了国库；每个人都争相贡献最多。他们面带喜悦欢欣。这时\Person{君士坦提乌斯}请伟大皇帝的使者亲自查看他的宝藏，并指示他们忠实报告所见；并补充说，此刻他把这笔钱收归己用，但这些钱早已由物主保管，如同在忠实司库手中一样安全。使者对所目睹之事惊叹不已；据说，在他们离开后，这位真正慷慨的君王派人召来物主，一一称赞他们的顺从和忠诚，然后将所有财物归还，并命他们回家。\par

因此，这一件事证明了我们试图描绘其品格的这位君王的慷慨；另一件事则清楚地证明他的虔敬。\par

% structure: p0038 source=h2 target=subsection reason=relative-heading-rank; base=h2

\subsection{第十五章　论其同僚发起的迫害}

根据帝国最高当局的命令，各省总督发起了对虔敬者的普遍迫害。实际上，最早的虔诚殉道者正是来自朝廷本身，他们经历了为信仰而战的冲突，最欣然地忍受了火、剑和深海；总之，各种死亡，以至于在短时间内，所有皇宫都被剥夺了虔诚之人。结果是，这些邪恶的发动者完全失去了天主的保护，因为他们迫害敬拜者，同时也沉默了通常为他们自己所做的祈祷。\par

% structure: p0040 source=h2 target=subsection reason=relative-heading-rank; base=h2

\subsection{第十六章　\Person{君士坦提乌斯}如何假装拜偶像，驱逐那些同意献祭的人，却留住在宫中愿意宣认基督的人}

另一方面，\Person{君士坦提乌斯}想出了一个充满睿智的策略，做了一件听起来矛盾但实际上最值得钦佩的事。\par

他向他宫廷中的所有官员，甚至包括最高权位者，提出以下选择：要么向魔鬼献祭，从而获准留在他身边，享受平日的尊荣；要么拒绝，则被禁止接近他本人，完全失去与他相识和交往的资格。因此，当他们各自做出选择，有人这样，有人那样；每个人的选择都已确定后，这位令人钦佩的君王揭示了他策略的隐秘含义，谴责一方的怯懦和自私，而高度赞扬另一方为天主而忠贞不渝。他还宣布，那些背叛他们天主的人必不配得到君王的信任；因为如果他们对更高权力不忠，他们怎么可能对他保持忠诚呢？因此，他决定这些人应全部被逐出宫廷；同时，宣布那些为真理作证、证明自己是天主忠仆的人，也会对君王表现同样的忠诚，于是将自身和帝国的守护托付给他们，说他对这样的人应予以特别尊重，视其为最亲近、最珍贵的朋友，且远比最宝贵的财富更珍视他们。\par

% structure: p0043 source=h2 target=subsection reason=relative-heading-rank; base=h2

\subsection{第十七章　论他的基督徒生活方式}

据说，\Person{君士坦丁}的父亲具有我们简要描述的那种品格。因这样的敬虔，天主赐给他何种死亡，并使他与帝国同僚的结局有何不同，任何人若留心此事的情况即可知晓。因为在长时间展示众多皇家美德之后，他承认惟独至高天主，谴责不敬虔者的多神论，并以圣人们的祈祷坚固其家室，他余生在卓越的安详与宁静中度过，享受被称为有福的状态——既不侵扰别人，自身也不受侵扰。\par

因此，在他平静安宁的整个统治期间，他将整个家庭、他的子女、妻子和家仆都奉献给独一至高天主：以致在他宫墙内聚集的团体与天主的教会毫无区别；其中也有他的神职人员，为他们的君王不断献上祈祷，而此时对大多数人而言，甚至与敬拜天主的人交谈一词也是不允许的。\par

% structure: p0046 source=h2 target=subsection reason=relative-heading-rank; base=h2

\subsection{第十八章　\Person{戴克里先}和\Person{马克西米安}逊位后，\Person{君士坦提乌斯}成为首席奥古斯都，并蒙福得众多子嗣}

这种行为立即得到天主的赏报，使他获得了帝国的最高权力。因为较年长的皇帝们不知何故退位；这一突然的变化发生在他们迫害教会的头一年之后。\par

從那時起，君士坦提烏獨自獲得了首席奧古斯都的尊榮，此前他已因帝國凱撒的冠冕而顯赫，在他們中居首位；但在他這份功績得到證明之後，他被授予了羅馬帝國的至高尊嚴，被命名為後來當選該榮譽的四位中的首席奧古斯都。此外，他在子女數量上超越了大多數皇帝，身邊聚集了非常多的兒女。最後，當他達到了幸福的晚年，即將償還自然的共同債務，並以此生換取另一生命時，天主再次以特殊的方式為他彰顯了祂的大能，使他的長子君士坦丁在他最後的時刻在場，並準備好從他手中接受帝國權力。\par

% structure: p0049 source=h2 target=subsection reason=relative-heading-rank; base=h2

\subsection{第十九章 論其子君士坦丁，他青年時曾陪同戴克里先前往巴勒斯坦。}

後者曾與他父親的帝國同僚在一起，並在他們中間度過了他的一生，正如我們所說，如同天主的古先知。甚至在他最早的青年時期，他們就斷定他配得最高的榮譽。我們自己曾見過一個實例，當時他陪同年長的皇帝經過巴勒斯坦，站在其右側，他當時展現出的帝王氣質令所有見到的人都讚歎不已。因為在儀表與身體的優雅及身高上，無人能與他相比；他在體力上如此超越同輩，以至於使他們畏懼。然而，他在心智品質的卓越上甚至比其優越的身體天賦更為突出；他首先擁有健全的判斷力，並且也受益於自由教育的薰陶。他因天生的才智和天主賦予的智慧而傑出，程度非比尋常。\par

% structure: p0051 source=h2 target=subsection reason=relative-heading-rank; base=h2

\subsection{第二十章 君士坦丁因戴克里先的陰謀逃往其父處。}

當時在位的皇帝們觀察到他剛毅而充滿活力的體格和卓越的心智，心中產生了嫉妒和恐懼，此後便小心伺機在他的品格上施加一些恥辱的標記。但這位年輕人察覺了他們的計謀，其細節因天主的眷顧不止一次傳到他耳中，他便尋求安全的逃脫；在這方面他又一次保持了他與偉大先知梅瑟的相似。事實上，在各方面天主都是他的幫助者；祂早已預定他應在場，準備好繼承他的父親。\par

% structure: p0053 source=h2 target=subsection reason=relative-heading-rank; base=h2

\subsection{第二十一章 君士坦提烏之死，他立其子君士坦丁為皇帝。}

\Emph{立即}，因此，他一脫逃了為他陰險設置的陰謀，就全速趕往他父親處，最終在他躺臥垂死之際到達。君士坦提烏一見他的兒子如此意外地在面前，便從臥榻上躍起，溫柔地擁抱了他，並宣告說，他在死亡面前唯一困擾他的焦慮——即因兒子的缺席而引起的焦慮——現已消除，他感謝天主，說他現在認為死亡勝於最長的生命，並立即完成了他的私人事務的安排。然後，在他自己的宮殿裡、在皇帝的臥榻上，向環繞他的兒女們作了最後的告別，按自然法則將帝國遺贈給了他的長子，便嚥下了最後一口氣。\par

% structure: p0055 source=h2 target=subsection reason=relative-heading-rank; base=h2

\subsection{第二十二章 君士坦提烏的葬禮後，軍隊如何宣告君士坦丁為奧古斯都。}

皇帝的寶座並未空置很久：因為君士坦丁穿上他父親的紫袍，從他父親的宮殿出發，向眾人呈現了他父親的生命與統治在他身上的復興。然後，他與他父親的朋友們一起舉行了葬禮行列，有些人先行，有些人隨後，並以非凡的華麗為這位虔敬的逝者舉行了最後的儀式，所有人都同心同聲以歡呼和讚美尊敬這位三重蒙福的君主，他們以一心一意讚美兒子的統治是逝者的復活，並立即以歡樂的呼喊頌揚他們的新君主為至尊與可敬的奧古斯都。因此，逝去的皇帝因對他兒子的讚美而受榮耀，後者則被稱為有福，因繼承了這樣一位父親。他治下的所有民族也充滿了喜樂和無法言喻的歡欣，因為一刻也沒有被剝奪良好秩序政府的益處。\par

在皇帝君士坦提烏的實例中，天主已向我們的時代顯明了那些在生活中尊敬和愛慕祂的人的結局。\par

% structure: p0058 source=h2 target=subsection reason=relative-heading-rank; base=h2

\subsection{第二十三章 暴君覆滅的簡短說明。}

至於其他那些攻擊天主的众教會的君主，我認為在本書中不宜記述他們的垮臺，也不宜因提及他們與相反品格的人相關聯而玷污對善人的記憶。事實本身的知識，對於那些親眼目睹或聽說過他們各自遭受的災禍的人來說，就足以作為健全的告誡。\par

% structure: p0060 source=h2 target=subsection reason=relative-heading-rank; base=h2

\subsection{第二十四章 君士坦丁獲得帝國是出於天主的旨意。}

因此，萬有的天主，整個宇宙的最高治理者，按祂自己的旨意任命了君士坦丁——這樣一位著名父母的後裔——為君主和統治者：如此，當其他人是因同類的選舉而被提升到這個尊位時，唯獨他，無人能誇口說對他的高陞作出了貢獻。\par

% structure: p0062 source=h2 target=subsection reason=relative-heading-rank; base=h2

\subsection{第二十五章 君士坦丁對野蠻人和不列顛人的勝利。}

他一登上王位，便开始关心父亲的疆业，以极大的仁慈访问那些先前受他父亲管辖的所有省份。一些居住在莱茵河畔和西海岸的蛮族部落胆敢反叛，他使他们全部归顺，并将他们从野蛮状态引入温顺。他满足于阻止其他部落的入侵，并像驱赶未驯服的野兽一般，将那些他看出完全无法适应文明生活稳定秩序的人驱逐出他的领土。在满意地处理完这些事务后，他将注意力转向世界的其他地区，首先渡海前往不列颠诸族，它们位于海洋的怀抱之中。他使他们臣服，然后着手考虑帝国其余地区的状况，以便随时准备在情况需要时提供援助。\par

% structure: p0064 source=h2 target=subsection reason=relative-heading-rank; base=h2

\subsection{第二十六章 他如何决心从\Person{马克森提乌斯}手中解救罗马}

因此，当他将整个世界视为一个巨大的身体，并看出这一切的首脑、罗马帝国的帝都已被暴政压迫所压垮时，起初他将解放的任务交给了统治帝国其他部分的人，视他们为年长者。但当他们中无人能够提供援助，而那些尝试过的人遭遇了惨败时，他说，只要他看见帝都如此受苦，生活便毫无乐趣，于是准备推翻暴政。\par

% structure: p0066 source=h2 target=subsection reason=relative-heading-rank; base=h2

\subsection{第二十七章 他反思了那些崇拜偶像者的败亡之后，选择了基督信仰}

然而，他确信自己需要比军队更强大的援助，因为暴君如此勤勉地施行邪恶的魔法和巫术，他便寻求天主的助佑，认为拥有武器和众多军队是次要的，却相信天主的合作力量是不可战胜、不可动摇的。因此，他考虑该依靠哪一位天主以获得保护和助佑。在此探寻中，他想到：在他之前的许多皇帝中，那些将希望寄托于众神并以祭献和供物侍奉他们的人，起初被谄媚的预言和承诺他们一切繁荣的神谕所欺骗，最终遭遇不幸结局，而他们的神祇没有一个站在他们身边警告他们即将到来的上天之怒；而有一位独独采取了截然相反的道路，谴责他们的错误，终其一生尊崇唯一至高的天主，发现祂是他帝国的救主和保护者，以及一切美善的赐予者。反思这些，并仔细权衡：那些信赖多神的人也以多种形式的死亡而倒下，没有留下家庭、后代、血统、名字或纪念；而另一方面，他父亲的天主却赐予他能力和很多标记的彰显。而且进一步考虑到，那些已经拿起武器对抗暴君、并在众神保护下进军战场的人，遭遇了可耻的结局（因为其中一人未战而可耻地退却；另一人则在自己军队中被杀，仿佛成了死亡的玩物）。我说，回顾所有这些考量，他判断，参与那些不是神的虚妄崇拜、在如此令人信服的证据之后仍偏离真理，实属愚妄；因此他感到有责任唯独尊崇他父亲的天主。\par

% structure: p0068 source=h2 target=subsection reason=relative-heading-rank; base=h2

\subsection{第二十八章 当他祈祷时，天主在正午的天空中向他显现一个光十字架的异象，并带有铭文劝告他以此得胜}

\Emph{于是}他以恳切的祈祷和哀求呼求祂，求祂启示祂是谁，并伸出右手帮助他克服当前的困难。当他这样热切祈祷时，一个最神妙的标记从天上显现给他，若是别人叙述此事，或许难以相信。但既然得胜的皇帝后来在很久以后——当这位历史的作者有幸与他相识交往时——亲自向作者宣告此事，并用誓言证实他的陈述，谁又能怀疑这个记述呢？尤其是后来的见证证实了它的真实性。他说，大约正午时分，当日已开始偏斜，他亲眼看见天上有一个光十字架的标记，在太阳之上，并带有铭文\Emph{以此得胜}。见此景象，他自己惊愕不已，他整个军队，即跟随他此次出征并目睹了这奇迹的军队，也同样惊愕。\par

% structure: p0070 source=h2 target=subsection reason=relative-heading-rank; base=h2

\subsection{第二十九章 天主的基督如何在他睡眠中显现，并命令他在战争中使用以十字架形状制成的军旗}

他还说，他心中疑惑这异象的含义是什么。当他继续沉思和推理其意义时，黑夜突然降临；然后在他睡眠中，天主的基督带着他曾在天上看到的同一个标记显现给他，命令他制作一个他在天上所见的标记的复制品，并在所有与敌人的交战中将其用作护符。\par

% structure: p0072 source=h2 target=subsection reason=relative-heading-rank; base=h2

\subsection{第三十章 十字架军旗的制作}

黎明时分，他起身，将这奇事告诉了他的朋友：然后，召集金匠和宝石匠，他坐在他们中间，向他们描述他所见标记的形状，命令他们用黄金和宝石制作出来。这复制品我本人曾有机会亲眼目睹。\par

% structure: p0074 source=h2 target=subsection reason=relative-heading-rank; base=h2

\subsection{第三十一章 十字架军旗的描述，罗马人现今称为\OriginalTerm{拉布兰旗}{Labarum}}

此旗幟的製作方式如下：一根長矛，外鍍黃金，上橫一桿，構成十字之形。全旗頂端固定一個由黃金和寶石編成的花冠；花冠之內，有救主之名的象徵，即兩個表示基督之名的起始字母，字母 P 被 X 從中心穿過；皇帝在後期習慣將此字母戴在頭盔上。從矛的橫桿懸掛一塊布料，一塊皇家織物，上面綴滿璀璨寶石的刺繡；並以黃金交織，呈現出難以形容的美麗。這面旗幟是方形的，其垂直旗桿下部甚長，上部在十字勝利的標記下方、繡旗正上方，承載著虔誠皇帝及其子女的黃金半身肖像。\par

皇帝經常使用這個救恩的標記，作為對抗一切敵對和邪惡勢力的保障，並命令將其他類似的旗幟在他所有軍隊前頭行進。\par

% structure: p0077 source=h2 target=subsection reason=relative-heading-rank; base=h2

\subsection{\Person{君士坦丁}如何接受教導並閱讀聖經}

這些事不久之後便發生了。但在上述時刻，他因那非凡的異象而驚愕不已，決意除了向他顯現的那位之外不敬拜別的天主，於是派人召來通曉祂教導之奧祕的人，詢問那位天主是誰，以及他所見異象的標記有何含義。他們斷言，祂是天主，獨一無二天主的獨生子；所顯現的標記是不朽的象徵，是祂從前在地上居住時所獲得的戰勝死亡的勝利紀念。他們也教導祂降臨的原因，並向他講解祂降生成人的真實事蹟。於是他接受了這些教導，對眼前所見的神性顯現感到驚嘆。因此，將天上的異象與解釋相比較，他發現自己的判斷得以證實；並確信這些知識是藉由神聖的教導傳授給他的，從此決定專心研讀聖經。\par

此外，他以天主的司鐸為顧問，並視之為己任，以全心敬意尊崇向他顯現的天主。此後，因著對祂滿懷堅定的希望，他急忙去撲滅暴政的威脅之火。\par

% structure: p0080 source=h2 target=subsection reason=relative-heading-rank; base=h2

\subsection{\Person{馬克森提烏斯}在\Place{羅馬}的淫亂行為}

因為那位暴虐地佔據了帝國都城的人，已在不敬神和邪惡的道路上走得很遠，以至於毫無顧忌地從事各種卑鄙污穢的行徑。\par

例如：他常將婦女與其丈夫分開，過一段時間再送回他們身邊，這些侮辱不是針對卑賤或默默無聞的人，而是針對\Place{羅馬}元老院中地位最高的人。此外，他雖可恥地玷污了無數的自由婦女，卻仍無法滿足他那不受約束、放縱的慾望。但當他試圖敗壞基督徒婦女時，卻再也無法達成其陰謀，因為她們寧願獻出生命而死，也不願將自己的身體交給他玷污。\par

% structure: p0083 source=h2 target=subsection reason=relative-heading-rank; base=h2

\subsection{一位總督之妻如何為貞潔自殺}

有一位婦人，是掌握總督權力的元老之妻（她是個基督徒），當她得知那些為暴君執行此類任務的人站在她家門前，並知道她的丈夫因恐懼而吩咐他們帶走她時，她請求給予片刻時間換上日常服飾，便進了內室。獨自一人留在那裡後，她將一把劍刺入自己的胸膛，隨即死去。她將屍體留給了那些拉皮條的人，卻以一個勝過任何言語的行為，向全人類，無論是當代還是後代，宣告：基督徒所著稱的貞潔，是唯一不可戰勝、不可摧毀的東西。這就是這位婦人所表現出的行為。\par

% structure: p0085 source=h2 target=subsection reason=relative-heading-rank; base=h2

\subsection{\Person{馬克森提烏斯}屠殺\Place{羅馬}人民}

因此，所有人，無論平民還是官員，無論地位高低，都因畏懼他那如我所描述的膽大妄為的邪惡而戰慄，並受其殘酷暴政的壓迫。儘管他們默然忍受，承受著這痛苦的奴役，但終究無法逃脫暴君嗜血的殘忍。例如，有一次，他以微不足道的藉口，將民眾交給自己的侍衛屠殺；無數的\Place{羅馬}人民，就在城中心，死於長矛和武器之下——不是死於斯基泰人或野蠻人之手，而是死於自己的同胞。此外，為了奪取他們的財產，被殺的元老數量也無法計算；在不同時期，以各種捏造的罪名，無數元老被處死。\par

% structure: p0087 source=h2 target=subsection reason=relative-heading-rank; base=h2

\subsection{\Person{馬克森提烏斯}對抗\Person{君士坦丁}的巫術與\Place{羅馬}的饑荒}

但暴君邪惡的頂點，是他訴諸巫術：有時為了巫術目的剖開孕婦的子宮，有時搜查新生嬰兒的內臟。他還殺死獅子，並施展某些可怕的伎倆來召喚惡魔，以阻止即將到來的戰爭，希望藉此取得勝利。總之，無法描述這個\Place{羅馬}暴君奴役其臣民的種種壓迫行徑；以至於此時他們已陷入極度的貧困和必需食物的匱乏，這種匱乏是我們同時代人從未記得曾在\Place{羅馬}發生過的。\par

% structure: p0089 source=h2 target=subsection reason=relative-heading-rank; base=h2

\subsection{\Person{馬克森提烏斯}軍隊在\Place{義大利}戰敗}

然而，\Emph{君士坦丁}充满对所有悲惨的怜悯，便开始用一切战争准备来武装自己以对抗专制。因此，他以至高天主为护佑，呼求祂的基督作他的保护者和援助，将胜利的凯旋纪念，即救恩标记，置于士兵和卫队之前，率领全军前进，试图为罗马人重新获得他们祖先传承的自由。\par

而马克森提乌斯，更信赖他的法术而非臣民的爱戴，甚至不敢出城门外，却用大批军队守卫了他专制统治下的每个地方、区域和城市。皇帝因信靠天主的助佑，向前对抗暴君的第一、第二和第三梯队，在第一次攻击中轻易击溃了他们全体，并进入了意大利的腹地。\par

% structure: p0092 source=h2 target=subsection reason=relative-heading-rank; base=h2

\subsection{第三十八章 马克森提乌斯死于台伯河桥。}

现在他已经非常接近罗马城本身了，为了使他免于因暴君之故而与全体罗马人作战，天主自己仿佛用隐秘的绳索将暴君远远地引出城门。此时，圣经记载的那些奇迹——天主昔日对抗不虔敬者所行的（大多数人不信为神话，但信者相信）——他事实上向所有人，无论信者还是不信者，这些奇迹的目击者，都予以证实。正如昔日摩西和敬拜天主的希伯来人的时期，\BibleQuote{法老的战车及其军队，他都投入海中；他精选的战车长，沉于红海。}——此时，马克森提乌斯及其士兵和卫队，\BibleQuote{如石头沉入深渊，}\BibleRef{出 15:5}，当他逃遁于君士坦丁的天助军队之前时，试图渡过挡在路上的河流，河上他用船只搭造了一座坚固的桥，却为自己制造了毁灭的机关，实际上是对付自己，而原希望藉此捕捉那位天主所爱者。因为他的天主站在一旁保护他，而另一人，无神，却证明是这些秘密装置的可怜设计者，自取灭亡。因此，人可正当地说：\BibleQuote{他挖了坑，又挖深了，却掉进自己挖的陷阱里。他的恶行必归到他头上，他的暴力必降到他的头顶。}如此，在目前的情况下，在天主指引下，桥上的机关和隐藏的伏兵，在预定时间之前意外地坍塌，桥开始下沉，船与人整体沉入水底。先是那恶人自身，然后他的武装侍从和卫队，正如圣言之前所描述的，\BibleQuote{如铅沉入浩瀚的水中。}\BibleRef{出 15:10} 因此，这些从天主获得胜利的人，即或不用同样的话语，却用与祂的伟大仆人摩西的百姓同样的精神，歌颂并讲说如同他们昔日针对那不虔敬的暴君所说的：\BibleQuote{我们歌颂主，因他受尊荣至极；马和骑手，他都投入海中。他成了我的救助和护盾，以致救恩。}又唱道：\BibleQuote{主啊，众神之中谁能像你？谁能像你，在圣洁中荣耀，在赞美中奇美，施行奇事？}\par

% structure: p0094 source=h2 target=subsection reason=relative-heading-rank; base=h2

\subsection{第三十九章 君士坦丁进入罗马。}

此时，唱过这些及类似的对天主的赞美诗，效法祂的伟大仆人摩西，君士坦丁进入帝国首都，凯旋而归。全体元老院和城中其他显贵，仿佛解除了监狱的束缚，与全体罗马民众一起，面容流露出内心的喜悦，以欢呼和洋溢的喜悦迎接他；男、女、孩童，与无数仆人，不停地呼喊称他为解救者、保存者和恩人。但他，内心对天主怀有虔诚，不因这些喝彩而傲慢，也不因听到的赞美而自高；反而，意识到自己从天主得到了帮助，立即向祂献上感恩，作为他胜利的创始者。\par

% structure: p0096 source=h2 target=subsection reason=relative-heading-rank; base=h2

\subsection{第四十章 君士坦丁手持十字架的雕像及其铭文。}

此外，他以公开宣告和不朽的碑刻向所有人昭示了救恩标记，在帝国首都中央设立了这对抗敌人的伟大胜利凯旋纪念，并且明确地以不可磨灭的文字刻上：救恩标记是罗马政府和整个帝国的保障。因此，他立即下令将一个十字架形状的高高举起的矛置于代表他自身的雕像的手下，在罗马最繁华的地方，并在其上用拉丁语刻上以下铭文：\Emph{因这救恩标记，即真正的勇武考验，我保全并解放了你们的城市于暴政的轭下。我也已使罗马元老院和人民获得自由，并恢复他们古代的荣耀与辉煌。}\par

% structure: p0098 source=h2 target=subsection reason=relative-heading-rank; base=h2

\subsection{第四十一章 各行省的欢庆与君士坦丁的仁政。}

如此，这位虔诚的皇帝以得胜十字架的宣认为荣，以极大的见证胆量向罗马人宣示了天主圣子。全城居民，元老院与民众，仿佛从暴虐统治的压迫中复苏，似乎享受更纯净的光辉，重生进入崭新生命。所有民族，远至\Place{西海}之滨，摆脱了此前困扰他们的灾难，因欢乐的节日而喜悦，不停赞颂他为得胜者、虔诚者、共同恩主；确实，所有人口同心声，宣称\Person{君士坦丁}是藉天主的恩宠为人类出现的普遍祝福。皇帝的敕令也遍传各处，使那些被不义剥夺地产的人得以重新享有自己的财产，而被不义流放者得以返回家园。此外，他释放了因暴君残忍而遭受这些苦难的人，使他们脱离监禁及各种危险与恐惧。\par

% structure: p0100 source=h2 target=subsection reason=relative-heading-rank; base=h2

\subsection{赐予主教的荣誉与建造教堂}

皇帝也亲自邀请天主的仆役们为伴，以最高的尊敬与荣誉对待他们，在言行上恩待他们，视他们为奉献于其天主服务的人。因此，他们虽衣着简朴、外貌卑微，却被接纳到他的餐桌；但在他的评价中并非如此，因为他认为自己看到的不是凡俗眼目所见的人，而是其中的天主。他还使他们成为他的旅途伴侣，相信他们事奉的那位天主会因此帮助他。此外，他从自己的私人资源中给予天主的教会慷慨的捐赠，既扩大又增高神圣的建筑，并以丰厚的奉献装饰教会的庄严圣所。\par

% structure: p0102 source=h2 target=subsection reason=relative-heading-rank; base=h2

\subsection{\Person{君士坦丁}对穷人的慷慨}

他还大量分发现金给那些有需要的人，此外，甚至对那些与他无涉的外教人也显出仁爱和恩惠；甚至对广场上可怜无助的乞丐，他不仅提供金钱或必需的食物，还提供体面的衣服。但对于那些曾经富裕、后又遭遇境况逆转的人，他的援助更为慷慨。他真正以帝王气度赐予他们丰厚的恩惠；有的赐予土地，有的以各种职位尊荣。他如父亲般照顾不幸者的孤儿，赈济寡妇的贫困，并特别关怀她们。他甚至将因父母去世而无依无靠的贞女嫁给他个人相识的富人。但这样做之前，他先赐予新娘适当的嫁妆，使她们能与婚姻结合。总之，正如太阳升起于大地时，慷慨地将光线普照众人，\Person{君士坦丁}也是如此，清晨从皇宫出发，仿佛与天上的光体一同升起，将自身恩惠的光线赐予所有来到他面前的人。几乎不可能接近他而不得到一些好处，也从未发生过任何期望得到他帮助的人希望落空的事。\par

% structure: p0104 source=h2 target=subsection reason=relative-heading-rank; base=h2

\subsection{他如何出席主教会议}

他对众人的一般态度便是如此。但他对天主的教会特别关心：而在几个省份中，有些人在判断上彼此分歧，他如同由天主设立的总主教一般，召集了他的仆役们的会议。他也不耻于出席并与他们一同坐在会议中，而是参与他们的讨论，为一切有关天主和平的事服务。他坐在他们中间，如同众人中的一员，遣散了侍卫、士兵和所有负责保卫他的人员；但受天主的敬畏所保护，并被忠实朋友的守护所环绕。他高度赞许那些他视为倾向于正确判断、表现出平静与温和态度的人，因为他显然乐于看见普遍的意见和谐；而对于那些固执己见的人，他则怀着厌恶看待。\par

% structure: p0106 source=h2 target=subsection reason=relative-heading-rank; base=h2

\subsection{他对无理之人的宽容}

此外，他耐心忍受了一些对他本身发怒的人，用温和柔顺的言辞指导他们自制，不要骚乱。其中一些人尊重他的劝告而停止了；但对于那些不能正确判断的人，他完全将他们交托给天主处置，自己从未想过对任何人采取严厉措施。因此，自然发生了\Place{非洲}的叛乱者竟至于暴力到敢于采取公开的鲁莽行为；似乎有一些恶灵，嫉妒眼前的大繁荣，驱使他们做出残暴的事，以激起皇帝对他们的愤怒。然而，这恶意的行为并未给他带来任何好处；因为皇帝嘲笑这些事，宣称它们源于恶者；这些不是清醒之人的行为，而是疯子或附魔者的行为；应当怜悯他们而非惩罚他们，因为惩罚疯人如同同情他们的状况一样是至高的仁爱。\par

% structure: p0108 source=h2 target=subsection reason=relative-heading-rank; base=h2

\subsection{战胜蛮族}

因此，皇帝在一切行为中尊崇天主，万有的主宰，并对他自己的教会行不懈的监督。天主酬报他，使所有蛮族臣服于他脚下，以致他能在各处树立对敌人的胜利纪念碑；并宣告他为全人类的征服者，使他成为对手的恐惧；但这并非他的天性，因为他其实是人类中最温良、最柔和、最仁慈的人。\par

% structure: p0110 source=h2 target=subsection reason=relative-heading-rank; base=h2

\subsection{\Person{马克西米努斯}之死，他曾试图谋反，以及\Person{君士坦丁}藉天主启示所发现的其他人}

正当他忙于此时，第二位逊位者因被查出参与叛国阴谋，遭受了极为可耻的死亡。他是第一个因其罪行和不敬，其画像、雕像及所有类似荣誉和尊荣的标志到处被销毁的人。此后，同一家族的其他成员也被发现暗中策划反对皇帝的阴谋；所有这些意图都通过天主在异象中向祂的仆人以奇妙的方式启示了出来。\par

因为祂经常赐予他自身的显现，神圣的临在以一种最奇妙的方式向他显现，并赐予他未来事件的许多预示。事实上，无法用言语表达天主乐意赐予祂仆人的神圣恩宠那不可言喻的奇事。被这些包围着他，他平安地度过了余生，喜乐于臣民的爱戴，也喜乐看见他治理下的所有人都过着满足的生活；但最令他欣喜的是天主的教会兴旺发达。\par

% structure: p0113 source=h2 target=subsection reason=relative-heading-rank; base=h2

\subsection{第四十八章 庆祝君士坦丁的十周年纪念}

就在他处于这种境况时，他完成了执政的第十年。此时，他下令举行普遍的庆典，并向天主、万有的君王献上感恩的祈祷，作为无火无烟的祭献。他从这一工作中获得了极大的愉悦；然而，他听到东方各省遭受蹂躏的消息时，却并不如此。\par

% structure: p0115 source=h2 target=subsection reason=relative-heading-rank; base=h2

\subsection{第四十九章 李锡尼如何压迫东方}

因为他得知在那地区，有一只凶猛的野兽正困扰着天主的教会和各省的其他居民，好像是邪灵努力产生与虔诚皇帝的行为完全相反的效果；以致罗马帝国分成两部分，在众人看来如同黑夜与白昼；因为黑暗笼罩了东方各省，而最明亮的白昼照亮了另一部分的居民。后者从天主手中领受了许多祝福，但这些祝福的景象却令那嫉恨一切善的嫉妒以及压迫帝国另一部分的暴君无法忍受；尽管他的统治顺利，且因与如此伟大的皇帝君士坦丁联姻而受尊荣，却不愿追随这位虔诚君主的脚步，反而竭力模仿不敬者的邪恶意图和行径；他宁愿选择亲眼见过那些可耻结局之人的道路，也不愿与他的上司保持友好关系。\par

% structure: p0117 source=h2 target=subsection reason=relative-heading-rank; base=h2

\subsection{第五十章 李锡尼如何企图谋反君士坦丁}

\Emph{因此}他对他的恩人发动了一场不可和解的战争，完全不顾友谊的法则、誓言的约束、亲族的纽带以及已有的条约。因为最仁慈的皇帝已通过将妹妹许配给他，给予了他真诚爱意的证明，从而赐予他家族关系及他自己古老皇室血统的特权，并授予他帝国同僚的级别和尊严。但对方却采取了完全相反的做法，密谋反对他的上司，用各种手段以伤害回报恩人。起初，他假装友谊，用诡计和欺骗行事，期望以此隐藏他的计谋；但天主使祂的仆人能够识破这些在黑暗中策划的阴谋。然而，他在初次尝试中被发现后，又诉诸新的欺诈；时而假装友谊，时而声称受到庄严条约的保护。然后突然违背每一项约定，又通过使节恳求宽恕，但最终可耻地背弃诺言，他宣布公开战争，并带着绝望的愚妄从那时起决心武装起来反对天主本身，因为他知道皇帝是崇拜天主的。\par

% structure: p0119 source=h2 target=subsection reason=relative-heading-rank; base=h2

\subsection{第五十一章 李锡尼对主教的阴谋及其禁止主教会议}

起初，他秘密查询他统治下的天主的仆人们——这些人从未在任何方面冒犯过他的政府——以便对他们提出虚假指控。当他找不到任何指控理由、也没有真正反对他们的理由时，他接着颁布了一项法律，规定主教们绝不可彼此交流，也不可有人离开本堂往访邻近教会；最后，也不得举行讨论共同事务的主教会议或议会。这显然是显示他对我们恶意的一个借口。因为我们要么被迫违反法律，从而应受惩罚，要么遵守其禁令，从而废止教会的法规；因为除非通过主教会议，否则不可能令人满意地解决重大问题。在其他情况下，这个憎恨天主的人，决心与爱天主的君主作对，也制定了这样的法律。因为前者召集天主的司祭是为了尊荣他们，促进和平与意见一致；而后者，其目的是摧毁一切善事，竭尽全力破坏整体的和谐。\par

% structure: p0121 source=h2 target=subsection reason=relative-heading-rank; base=h2

\subsection{第五十二章 基督徒的流放及其财产的没收}

\Person{君士坦丁}，这位天主之友，曾将自由出入皇宫的权利赐予祂的敬拜者；而这位天主的仇敌，却以完全相反的精神，将他权下所有的基督徒都逐出皇宫。他放逐了那些曾证明自己是他最忠诚、最忠心的臣仆，又强迫那些他本人曾因其卓越功勋而授予荣誉和地位的人，像奴隶一样为他人做卑贱的差役。最终，他图谋将所有人的财产据为己有，甚至以死亡威胁那些承认救主名号的人。此外，他本人因纵欲而堕落无救，又因不断行奸淫和其他无耻恶行而卑劣不堪，竟以自己的卑劣品性作为人类本性的标本，否认人间存在贞洁和节制的德行。\par

% structure: p0123 source=h2 target=subsection reason=relative-heading-rank; base=h2

\subsection{第五十三章 命令妇女不得与男子在教堂聚集}

\Emph{因此}他通过了第二条法律，命男子不得与女子一同出现在祈祷之所，并禁止女子参加神圣的德行学校，也不得接受主教的教导，而指定女子担任同性的教师。这些规定遭到普遍嘲笑，于是他另谋手段以摧毁教会。他命令民众惯常的集会应在城门外空旷的乡间举行，声称城外的开阔空气远比城墙内的祈祷之所更适合群众集会。\par

% structure: p0125 source=h2 target=subsection reason=relative-heading-rank; base=h2

\subsection{第五十四章 拒绝献祭者应被革除军职，监禁者不得给食}

然而，他在这一点上也未能得到顺服。最终，他撕下伪装，下令：凡在帝国各城市担任军职者若拒绝向魔鬼献祭，即应被剥夺各自职务。于是，各省长官的队伍都失去了敬拜天主的人；而颁布此令的他本人也遭受损失，因为他剥夺了自己获得虔诚者祈祷的权利。我又何必再提他如何命令，无人应遵守普通的人情，将食物分给在监狱中受折磨的人，甚至不得怜悯那些饿死的囚徒？简而言之，无人可行善事，那些本性驱使他们同情同胞的人，竟不得向他们施以丝毫仁慈。这实在是所有法律中最无耻、最可耻的，超越了人类本性最恶劣的堕落：一条法律，竟对施以怜悯者施以与受怜悯者相同的惩罚，将纯粹的人道行为处以最严厉的刑罚。\par

% structure: p0127 source=h2 target=subsection reason=relative-heading-rank; base=h2

\subsection{第五十五章 \Person{里基尼乌斯}的无法行径与贪婪}

此乃\Person{里基尼乌斯}的法令。然而，我又何必列举他在婚姻方面、以及在临终之事上的新规？他竟然妄图废除罗马人古老而明智的法律，代之以野蛮残忍的制度，编造千百种借口压迫臣民。正因如此，他发明了一种新的丈量土地方法，将最小的一块地算得大于实际面积，出自贪得无厌的欲望。同样，他登记早已不在人世、久已逝去的乡村居民姓名，以此手段获取可耻的收益。他的吝啬毫无限度，他的贪婪永无止境。因此，当他用黄金、白银和无尽的财富装满所有国库时，他却痛呼自己的贫穷，仿佛遭受坦塔罗斯的折磨。但我又何必提及他流放了多少无辜者，没收了多少财产，囚禁了多少出身高贵、品德可敬的人，将他们的妻子交给他的放荡奴隶卑劣地侮辱，以及他自己——虽已感到年迈体弱——对多少已婚妇女和贞女施加暴力？我不必详述这些，因为他最后的行为之恶劣，使得前事显得微不足道。\par

% structure: p0129 source=h2 target=subsection reason=relative-heading-rank; base=h2

\subsection{第五十六章 他终于发起迫害}

他愤怒的最后表现，便是在公开敌视教会中，将攻击矛头直指主教们，视他们为最大的仇敌，尤其仇恨那些伟大的虔诚皇帝视为朋友的人。因此，他向我们将怒火倾泻到极致，丧失理智，既不顾念此前迫害基督徒者的榜样，也不顾及那些他本人被兴起去惩罚和毁灭因其不敬行为之人：他也不理会自己亲眼所见的事实，尽管他亲眼看到我们这些灾难的主要发起者（无论他是谁）被天主的鞭笞所击打。\par

% structure: p0131 source=h2 target=subsection reason=relative-heading-rank; base=h2

\subsection{第五十七章 \Person{马克西米安}因瘘管溃疡生蛆而衰弱，颁布谕令支持基督徒}

因为此人发起了对教会的攻击，并首先以义人和虔诚之人的血玷污了自己的灵魂，天主就对他施行审判，起初影响他的身体，后来延及他的灵魂。突然，他私处出现脓肿，随后形成深层的瘘管；这些疾病以无法医治的毒性侵袭肠子，肠内滋生大量蛔虫，发出恶臭。此外，他整个人因暴食过度而堆积了大量脂肪，如今正在腐烂，据说所有接近他的人都感到难以忍受的可怕景象。因此，他在与这些痛苦斗争之后，最终，尽管为时已晚，他认识到自己过去对教会犯下的罪行；在天主前承认自己的罪过后，他停止了迫害基督徒，并急忙颁布重建其教会的帝国诏令和敕书，同时命令他们进行惯常的敬礼，并为他祈祷。\par

% structure: p0133 source=h2 target=subsection reason=relative-heading-rank; base=h2

\subsection{第五十八章 迫害基督徒的\Person{马克西米努斯}被迫逃亡，并伪装成奴隶藏匿。}

然而，我们现在所说的这位，他曾亲眼目睹这些事，并凭自己的亲身经历知道这些事，却突然从脑海中抹去了对它们的记忆，既没有思考第一个迫害者的惩罚，也没有思考第二个迫害者身上所执行的神圣审判。后者确实试图在罪恶生涯中超越其前任，并以发明针对我们的新酷刑而自豪。火、剑、钉刺、野兽或深海都不足以满足他。除了所有这些，他还发现了一种新的惩罚方式，并颁布法令规定应毁掉他们的视力。于是，不仅许多男子，还有妇女和儿童，在失去视力和双脚关节（通过切割或烧灼）后，被这样送到矿场从事痛苦的劳役。因此，这个暴君不久之后也被天主的正义审判所赶上，当时他正信赖他所崇拜为神祇的魔鬼的帮助，并依靠他无数军队，冒险投入战斗。因为，在那次场合，他感到自己失去对天主的一切希望，就脱下了与他不相称的帝国服装，怯懦地躲在周围的人群中，企图通过逃跑获得安全。\par

此后，他穿着奴隶的服装在田野和村庄中潜藏，希望以此有效地隐藏自己。然而，他未能逃脱天主大而鉴察万有的眼睛：因为就在他期待安然度过余生时，他被天主燃烧的箭击中扑倒在地，全身被神圣复仇的打击所吞噬；以致他原本面貌的痕迹全部消失，只留下枯骨和骷髅般的外观。\par

% structure: p0136 source=h2 target=subsection reason=relative-heading-rank; base=h2

\subsection{第五十九章 因病失明的\Person{马克西米努斯}颁布了一道有利于基督徒的法令。}

而天主的打击仍沉重地加于他，以致他的眼睛突出并脱离眼眶，使他完全失明：这样，他以最公义的报应遭受了他最初为天主的殉道者所设计的同样惩罚。然而，甚至经历这些痛苦之后，他也恳求基督徒的天主宽恕，并承认他敌对天主的亵渎行为；他也像之前的迫害者一样撤回前言；并通过法律和法令明确承认他在崇拜他视为神祇的那些对象上的错误，宣告他现在通过切身体验知道基督徒的天主是唯一真神。这些事实\Person{李锡尼}不仅从他人的见证中得知，而且本人亲身知道；然而，好像他的理智被某种错误的乌云遮蔽一般，他仍然坚持同样的邪恶道路。\par

\SourceRecord{Translated by Ernest Cushing Richardson.；Nicene and Post-Nicene Fathers, Second Series；Vol. 1.；Edited by Philip Schaff and Henry Wace.；Buffalo, NY: Christian Literature Publishing Co.,；1890.；<http://www.newadvent.org/fathers/25021.htm>.}

% structure: p0001 source=h1 target=section reason=page-title

\section{\WorkTitle{君士坦丁传（卷二）}}

% structure: p0002 source=h2 target=subsection reason=relative-heading-rank; base=h2

\subsection{第一章　李锡尼乌斯的秘密迫害：在本都的阿马西亚处死几位主教}

我们所说之人，就这样继续冲向那等待天主的仇敌的毁灭；他再一次以致命的模仿，效法那些他亲眼目睹因其不虔敬行为而自取灭亡之人的榜样，重新点燃了对基督徒的迫害，像熄灭已久之火，并把这亵渎的火焰煽得比任何前人都更猛烈。\par

起初，他虽如凶猛的野兽或弯曲扭动的蛇一般，向天主喷吐愤怒与威吓，但因畏惧君士坦丁，不敢公然攻击其治下属于天主教会；而是掩饰其恶毒的敌意，试图通过秘密与有限的手段致主教们于死地，他找到方法，借各省总督提出的指控，除掉了其中最杰出者。他们受苦的方式带有某种奇异且前所未闻的特征。无论如何，在本都的阿马西亚所施行的残暴，超过了任何已知的残酷行径。\par

% structure: p0005 source=h2 target=subsection reason=relative-heading-rank; base=h2

\subsection{第二章　拆毁教堂，屠杀主教}

因为在那城中，自迫害开始以来，有些教堂再次被夷为平地，另一些则由各区总督关闭，以防止常去的人聚集，或向天主献上应有的敬拜。发出这些暴行命令的人，太清楚自己的罪行，不致期望这些服务是为他的益处而进行的；他确信我们所做的一切，以及我们为求得天主恩宠的一切努力，都是为了君士坦丁的缘故。\par

这些谄媚的总督们确信这种行为会取悦不虔敬的暴君，便将教会中最杰出的主教们处以死刑。于是，无罪的人被带走，毫无理由地像谋杀犯一样受罚；有些人遭遇了一种新的死亡，身体被一块块割开；经过这比任何悲剧更可怕的残酷刑罚后，被扔进深海作为鱼食。这些恐怖的结果再次如先前一般，虔诚的人们逃亡，田野和荒漠再次接纳了天主的敬拜者。暴君既已在这目标上得逞，更进一步决定发动对基督徒的全面迫害；他本会实现其目的，没有什么能阻止他执行决议，若不是那保护自己子民的天主预见了即将来临的灾祸，并以他特别的引导，将他的仆人君士坦丁带到帝国这一部分，使他如明亮的光辉照耀在黑暗和幽暗之夜中。\par

% structure: p0008 source=h2 target=subsection reason=relative-heading-rank; base=h2

\subsection{第三章　君士坦丁如何为面临迫害危险的基督徒所激动}

他察觉到所听闻的祸患已无法忍受，便采取明智的谋划，将天性的仁慈与某种程度的严厉调和，急忙去救助那些深陷压迫的人。因为他判定，通过除掉一个人来确保人类大部分的安全，理当被视为一项虔诚圣善的任务。他也判定，若他只听凭仁慈的发号施令，将他的怜悯施予一个完全不配的人，则一方面不会给一个无论如何都不肯放弃恶行的人带来真正益处，而且他对其臣民的暴怒只会增加；另一方面，那些受他压迫的人将永远失去得救的希望。\par

受这些思考的影响，皇帝决定不再拖延，向那些陷入如此极端苦难的人伸出援助之手。他因此进行了惯常的备战准备，调集了全部兵马。但在这一切之前，高举着我之前描述过的旗帜，那是他对天主完全信靠的象征。\par

% structure: p0011 source=h2 target=subsection reason=relative-heading-rank; base=h2

\subsection{第四章　君士坦丁以祈祷备战，李锡尼乌斯以占卜备战}

他也带上了天主的司祭们，深知此时比任何时候都更需要祈祷的功效，并认为他们应常在他身边陪伴，作为灵魂最忠实的守护者。\par

当暴君得知君士坦丁对敌人的胜利，并无其他原因，唯因天主的协助，而且上述之人常与他同在并陪伴左右；此外，那救恩苦难的标记还走在皇帝本人及其全军之前；他对这些防备嗤之以鼻（正如所料），同时以亵渎的话嘲弄辱骂皇帝。\par

另一方面，他聚集了埃及的占卜师、算命者、巫师和术士，以及那些他以为是神祇的祭司和先知。然后，他献上他认为时机所需的祭品，询问他对战争的成功结局能抱多大指望。他们异口同声地回答，他无疑会战胜敌人，在战争中得胜；各地的神谕也以华丽的诗句向他预示同样的前景。占卜师以鸟飞兆象保证这是吉兆；祭司则宣告祭品内脏的颤动也显明同样的事。于是，在这些虚假的保证下，他得意洋洋地率领军队前进，准备战斗。\par

% structure: p0015 source=h2 target=subsection reason=relative-heading-rank; base=h2

\subsection{第五章　李锡尼乌斯在树林中献祭时关于偶像和基督所说的话}

当他准备交战时，他请他最信任的侍卫和最珍贵的朋友到他视为神圣之处与他相会。那是一个水源充足、树荫茂密的园林，园内有多座他所认为的神祇的大理石雕像。他点燃蜡烛，按惯例举行祭献以尊崇这些神祇后，据说发表了如下讲话：\par

\Quote{朋友们，战友们！这些是我们祖国的神祇，我们以源自远古祖先的敬礼尊崇他们。但如今统率与我们敌对之军的人，已背弃了祖先的宗教，采纳了无神论的观念，在愚昧中敬拜某种陌生而前所未闻的天主，以其可鄙的旗帜玷污了他的军队，仰仗该天主的援助而拿起武器，如今进军，与其说是针对\Emph{我们}，不如说是针对他所离弃的那些神祇。然而，此刻将要证明我们中谁判断有误，并在我们的神祇与我方敌军所声称尊崇的神祇之间作出裁决。因为要么它宣告胜利属于我们，从而最公正地表明我们的神祇是真正的拯救者和扶助者；要么，如果这位君士坦丁的天主——我们不知其从何而来——确实胜过我们的神祇（他们数量众多，至少从数量上占优势），那么从此以后任何人都不应再怀疑应该敬拜哪位天主，而应立即依附于那更强大的力量，并将胜利的光荣归于他。那么，假设这位我们如今视为笑柄的陌生的天主真的得胜，那时我们确实必须承认并尊崇他，并向那些我们为之徒然点燃蜡烛的神祇永别。但如果我们的神祇得胜（他们无疑会胜利），那么，一旦我们赢得眼前的胜利，就应立即对那些蔑视神祇的人开战，毫不迟延。}\par

这就是他对当时在场者所说的话，不久之后，一些亲耳听闻的人将此报告给本历史作者。他讲完话后，立即下令部队开始进攻。\par

% structure: p0019 source=h2 target=subsection reason=relative-heading-rank; base=h2

\subsection{第六章 在里基尼乌斯统治下的诸城中所见的异象，如同君士坦丁的军队穿城而过}

正当这些事情发生时，据说在暴君统治下的诸城中观察到超自然的现象。君士坦丁军队的不同分队似乎在正午时分穿过这些城市显现在人们眼前，仿佛他们已获得胜利。实际上，当时任何地方都没有一个士兵在场，但这景象是通过一种神圣而更高的力量显现的，预示着即将发生的事。因为一旦双方军队准备交战，那破坏友好盟约的人首先发起战斗；于是，君士坦丁呼号\Quote{至高的救主天主}之名，并将此作为口令传给士兵，在第一次交锋中就战胜了他；不久之后，在第二次战役中，他又获得一次更重要、更具决定性的胜利，救恩的旗帜行于他的军队前列。\par

% structure: p0021 source=h2 target=subsection reason=relative-heading-rank; base=h2

\subsection{第七章 胜利随处紧随十字架旗帜在战场出现}

确实，无论这旗帜出现在何处，敌军很快就在他得胜的军队面前溃逃。皇帝察觉此事，每当他看到任何部分军队处境艰难，就下令将那救恩的旗帜移往该方向，如同某种抵御灾难的胜利护符；于是，战斗人员仿佛受到神圣激励，获得新的力量和勇气，胜利随即而来。\par

% structure: p0023 source=h2 target=subsection reason=relative-heading-rank; base=h2

\subsection{第八章 选出五十人背负十字架}

于是，他从侍卫中挑选了那些在体力、勇气和虔诚方面最杰出的人，委托他们专门保管和护卫这旗帜。因此，不少于五十人，其唯一职责就是包围并警惕地护卫旗帜，他们轮流将旗帜扛在肩上。这些情景是在事件发生很久以后，皇帝本人在闲暇时向本叙述作者讲述的；他还补充了另一件值得记录的事情。\par

% structure: p0025 source=h2 target=subsection reason=relative-heading-rank; base=h2

\subsection{第九章 一名背十字架者擅离职守逃跑被杀，另一名忠于职守者得以保全}

因为他说，有一次在战斗最激烈时，突然的骚动和恐慌袭击了他的军队，使当时扛旗的士兵陷入极度恐惧，以至于他将旗交给另一个人，以便自己逃离战斗。然而，当他的同伴接过旗，他退出并放弃了对旗帜的一切责任时，他被一支标枪击中腹部而致命。这样，他为自己的怯懦和不忠付出了代价，当场倒地而死；但那位接替他作为救恩旗帜背负者的人，却发现这旗帜成了他生命的保障。因为尽管他遭到连续的标枪射击，但这人安然无恙，旗帜的杆承受了所有武器。这确实是极其奇妙的情形：敌人的标枪都落在这支细长的矛范围内并留在其中，从而拯救了旗手免于死亡；因此，任何从事这一职守的人都从未受过伤。\par

这个故事并非我杜撰，此事我也同样归功于皇帝本人的权威，他当着我的面连同其他事情一起讲述。如今，借着天主的能力取得这些初步胜利后，他调动军队继续前进。\par

% structure: p0028 source=h2 target=subsection reason=relative-heading-rank; base=h2

\subsection{第十章 诸场战役与君士坦丁的胜利}

然而，敌军的先锋无法抵挡皇帝的第一波进攻，便放下武器，俯伏在他脚前。他全都饶恕了，喜悦于保全人命。但仍有其他人继续持有武器，投入战斗。皇帝试图以友好的提议安抚他们，但当这些不被接受时，他便命令军队发动攻击。于是他们立刻转身逃跑；有些人被追上并按战争法则杀死，另一些人在混乱的逃亡中互相践踏，死在自己同伴的剑下。\par

% structure: p0030 source=h2 target=subsection reason=relative-heading-rank; base=h2

\subsection{第十一章 里基尼乌斯的逃亡与巫术。}

在这种情况下，他们的统帅发现自己失去了追随者的援助，既丧失了他不久前还颇为庞大的常备军与盟军，也亲身证明了他对那些他以为是神明的依靠是多么徒然，于是可耻地逃跑了；他确实现身逃脱，保全了性命，因为虔诚的皇帝禁止士兵过于紧追，从而给了他逃跑的机会。皇帝这样做是希望他日后在认清自己绝境时，能放弃那疯狂而狂妄的野心，回归更健全的理智。君士坦丁以其过度的仁爱如此思想，并愿意耐心忍受过去的伤害，将这宽恕赐予一个如此不配的人；但里基尼乌斯远未放弃他的恶行，反而罪上加罪，敢行比以往更胆大包天的暴行。不仅如此，他再次妄行可憎的巫术，又一次狂妄自大；因此，人们完全可以像论及昔日埃及暴君那样说他：\BibleQuote{天主使他的心硬了}。\BibleRef{出 9:12}\par

% structure: p0032 source=h2 target=subsection reason=relative-heading-rank; base=h2

\subsection{第十二章 君士坦丁如何在帐幕中祈祷并获得胜利。}

但当里基尼乌斯放纵于这些不敬，盲目地冲向毁灭的深渊时，皇帝在另一方面，当他看到自己将不得不在第二次战役中迎战敌人时，便将中间的时间奉献给他的救主。他在营外远处支搭了十字架的帐幕，在那里以纯洁而神圣的方式度过时光，向天主献上祈祷；效法他那位古先知的榜样，圣经曾见证说，他在营外支搭了帐幕。他仅由几位他极为看重其信仰与虔诚的人陪伴。每逢他计划与敌人交战，他都继续遵守这一习惯。他行事审慎，以求稳妥，并渴望在一切事上受神圣指引。他热切地向天主恳求，总是稍后便蒙受他临在的显现。然后，仿佛为神圣的冲动所驱使，他从帐幕中冲出，突然命令他的军队立即行动，刻不容缓，并立刻拔剑。他们随即发动攻击，勇猛作战，以致不可思议地迅速取得胜利，并竖立起战胜敌人的胜利纪念碑。\par

% structure: p0034 source=h2 target=subsection reason=relative-heading-rank; base=h2

\subsection{第十三章 他对待俘虏的仁慈。}

因此，皇帝和他的军队长期以来每当有交战的可能时，便习惯于如此行动；因为他的天主始终在他心中，他渴望按他的旨意行事，并凭良心避免任何草菅人命。他不仅关心自己臣民的生命，甚至也关心敌人的生命。因此，他命令他的胜利之师饶恕俘虏的性命，告诫他们作为人类，不要忘记共同本性的要求。每当他看到士兵的激情失控时，他便以金钱赏赐来平息他们的愤怒，奖赏每一个救了敌人性命的人一定重量的黄金。而皇帝自己的智慧使他发现了这种激励保全人命的办法，因此，甚至大量蛮族也被这样救下，他们的性命全仗皇帝的黄金。\par

% structure: p0036 source=h2 target=subsection reason=relative-heading-rank; base=h2

\subsection{第十四章 再次提及他在帐幕中的祈祷。}

现在，这些以及成千上万类似的行为，是皇帝所熟悉并习惯性去做的。而在当时，他照例在战前退入他帐幕的隐密处，在那里以祈祷向天主祈祷。同时，他严格戒除一切安逸或奢侈的生活，以禁食和身体克苦来管束自己，以恳求和祈祷祈求天主的恩宠，以便得着他的赞同与帮助，并随时准备执行他所愿进入他思想的一切。简言之，他对所有人一视同仁地施以警醒的关怀，并为敌人的安全与为自己子民的安全，一同向天主代求。\par

% structure: p0038 source=h2 target=subsection reason=relative-heading-rank; base=h2

\subsection{第十五章 里基尼乌斯奸诈的友谊与偶像崇拜行为。}

而那个不久前在他面前逃跑的人，如今掩饰自己的真实想法，再次请求恢复友谊与同盟，皇帝认为适宜在某些条件下允准他的请求，期望这样的措施可能有益，并普遍有利于大众。然而里基尼乌斯虽然假装顺从所规定的条件，并以誓言证明他的诚意，但就在此时，他却在暗中集结军事力量，再次密谋战争与冲突，甚至邀请蛮族加入他的旗帜，并且开始为他寻求其他神明，因为他在迄今所信靠的那些神明上受了欺骗。他不曾思考自己曾公开谈论假神的事，或选择承认那曾为君士坦丁而战的天主，反而因寻求众多新神而显得可笑。\par

% structure: p0040 source=h2 target=subsection reason=relative-heading-rank; base=h2

\subsection{第十六章 里基尼乌斯如何劝告他的士兵不要攻击十字架的旗帜。}

如今，\Person{君士坦丁}已藉经验认识到那居于救恩之旗中的神圣而奥秘的力量，他的军队藉此力量惯于得胜；他便告诫士兵们，绝不可将攻击指向这旗帜，甚至不可不慎让目光停留其上；他向他们保证，这旗帜拥有可怕的力量，尤其与他为敌；因此他们最好小心避免与它相撞。既已如此吩咐，他便准备与那仍因人性而犹豫、推迟他预见到将临于其对手之命运的人决战。\par

% structure: p0042 source=h2 target=subsection reason=relative-heading-rank; base=h2

\subsection{第十七章　\Person{君士坦丁}的胜利}

但他一察觉敌人坚持其决定，且已拔剑在手，便尽情发泄义怒，仅一次冲锋便瞬间击溃了全部敌军，如此一举战胜了他们和他们的神祇。\par

% structure: p0044 source=h2 target=subsection reason=relative-heading-rank; base=h2

\subsection{第十八章　\Person{李锡尼}之死及庆典}

他继而依照战争之法处置这天主的敌人及其追随者，使之受到适当的惩罚。因此，那暴君本人及其在邪恶中支持他计谋的人，一同被处以正义的死刑。此后，那些不久前还被对假神的虚妄信心所欺骗的人，真诚地承认了\Person{君士坦丁}的天主，公开宣认他们相信他是真实而唯一的天主。\par

% structure: p0046 source=h2 target=subsection reason=relative-heading-rank; base=h2

\subsection{第十九章　欢庆与节日}

如今，恶人既被除去，太阳在暴虐政权的阴沉乌云之后再次灿烂照耀。罗马统治的每个部分彼此融合；东方各民族与西方联合，整个罗马帝国仿佛由其首脑即一位单一的最高统治者所妆饰，他的唯一权威遍及全体。如今，虔诚之光的明亮光辉使那些先前坐在黑暗和死影中的人日子欢欣。往日的痛苦不再被记念，因为众人都联合赞美得胜的君王，并承认他的护守者乃是唯一真神。这样，那位以一切虔诚美德闪耀其品格的皇帝——\OriginalTerm{胜利者}{Victor}（他自己采用了这个最合适的称号，以表达天主赐予他战胜一切恨他或反对他之人的胜利）——取得了东方的统治权，从而单独治理罗马帝国，使之如前所述重新统一在一个首脑之下。因此，正如他首先向所有人宣告天主的唯一主权，他自己作为罗马世界的唯一君主，将其权柄延及全人类。每个人曾受压迫之苦的恐惧现已消除；曾低头悲伤的人如今以笑容和表达内心喜悦的眼神彼此相视。他们首先按所受的教导，以游行和赞美诗将最高主权归于天主，实为万王之王；然后以不断的欢呼向得胜的皇帝及其最明智而虔诚的儿子们——诸凯撒——致敬。先前的苦难被遗忘，所有过去的邪恶被宽恕；在享受现世幸福的同时，也混合着对未来持续祝福的期盼。\par

% structure: p0048 source=h2 target=subsection reason=relative-heading-rank; base=h2

\subsection{第二十章　\Person{君士坦丁}为宣信者颁布的法令}

此外，皇帝充满仁慈精神的谕令也颁布到我们中间，正如在帝国另一部分的居民中一样；他的法律，散发出对天主的虔诚精神，预兆着多重的祝福，因为它们为各省的臣民在各民族中确保了许多利益，同时规定了适合天主教会紧急需要的措施。首先，它们召回了那些因拒绝参与偶像崇拜而被各自省份总督放逐或赶出家园的人。其次，它们解除了那些因同样原因被判决在民事法庭服役之人的负担，并命令归还任何被剥夺财产之人的财产。那些在考验时期因天主的事业而以灵魂刚毅著称的人，因此被判在矿场劳役、或被放逐到孤岛、或被强迫在公共工程中劳作的人，都立即解除了这些负担；另一些人，因宗教上的恒心而失去了军阶，皇帝的慷慨洗刷了这耻辱：他准许他们要么恢复军阶并享受先前的特权，要么如果他们更喜欢稳定的生活，则永久免除一切兵役。最后，所有因耻辱和侮辱被迫从事妇人劳役的人，他也同样与其他获释者一同解放了。\par

% structure: p0050 source=h2 target=subsection reason=relative-heading-rank; base=h2

\subsection{第二十一章　关于殉道者及教会财产的法律}

这些就是皇帝书面命令为那些为信仰受难者的人身所确保的恩惠，他的法律也为他们的财产作了充分的规定。\par

至于那些为宣认主名而舍生的天主圣殉道者，他指示其产业应由至亲继承；若至亲缺失，则继承权归于教会。此外，凡由国库转给他人的财产，不论是以售卖或赠与方式，连同仍保留在国库中的，皇帝的慷慨谕令均指示应归还原主。他如此广泛施予的恩惠，为天主的教会带来了这些益处。\par

% structure: p0053 source=h2 target=subsection reason=relative-heading-rank; base=h2

\subsection{第二十二章　他如何赢得民众的欢心。}

但他的慷慨还给予帝国中的异教民族和其他国家更广泛、更众多的恩惠。因此，我们[东方]地区的居民，听说帝国另一部分所享有的特惠，曾祝福那些幸运的受益者，并渴望自己能享有同样命运，如今他们亲眼目睹自己拥有这一切福泽，便异口同声地宣告自己的幸福，并承认这样一位君主出现在人间，实属奇妙之事，是世界历史从未记载过的。这便是他们的心声。\par

% structure: p0055 source=h2 target=subsection reason=relative-heading-rank; base=h2

\subsection{第二十三章　他宣告天主是其繁荣的根源：关于他的谕旨。}

如今，因着天主救主的强大援助，万族都臣服于皇帝的权下，他便公开向众人宣告那赐予他一切恩惠之天主的圣名，宣称是祂，而非他自己，是他过往胜利的根源。他下令将此宣言以拉丁文和希腊文写就，传谕帝国各省。其文辞之精妙，可从其两封书信中得知：一封致天主的教会，另一封致帝国各城市的异教民众。我思量将后者插入此处，因其与我当前主题相关，一则使此文件的历史副本得以记录并传之后世，二则有助于证实我当前叙述的真实性。此件取自皇帝法令的权威副本，由我本人持有，皇帝亲笔签名，宛若为我的陈述盖上了真实的印记。\par

% structure: p0057 source=h2 target=subsection reason=relative-heading-rank; base=h2

\subsection{第二十四章　君士坦丁关于对天主的虔诚及基督宗教的法律。}

\OriginalTerm{胜利者君士坦丁，最伟大的奥古斯都}{Victor Constantinus, Maximus Augustus}致巴勒斯坦省居民。\par

凡对至高者的属性抱有公正且健全见解之人，早已最清楚地看出，而且毋庸置疑，那些谨守基督宗教神圣义务者与那些敌视或藐视此宗教者之间，自古存在何等巨大的差别。然而在当今，我们可通过更显明的证据和更决定性的实例，看到质疑此真理是何等不合理，以及至高天主的能力何其强大：因为显然，凡忠实遵守祂圣洁律法、畏惧触犯其诫命者，得蒙丰盛祝福，并拥有稳固的希望和实现其事业的充足能力。反之，凡怀有邪恶思想者，则遭遇到与它们邪恶选择相应的结果。因为一个既不愿承认也不愿适当敬拜那作为一切祝福根源之天主的人，怎能期待获得任何祝福？事实本身便是我所说的印证。\par

% structure: p0060 source=h2 target=subsection reason=relative-heading-rank; base=h2

\subsection{第二十五章　古代实例。}

因为确实，任何人若在心中回溯从最早时期至今的事件进程，并反思过去所发生的事，便会发现，凡以正义和正直为行为基础者，不仅使其事业得以成功，而且犹如从这甘美的根上收获了甜蜜的果实。再者，凡观察那些大胆施行压迫或不义的人之生涯者——他们或将其无知狂怒指向天主自身，或对同胞毫无善意，竟敢以放逐、羞辱、没收财产、屠杀或其他类似苦难加害于他们，而毫无悔意或无意转向更善之路——他将发现那些人已得到与其罪行相称的报应。这些结果是自然且合理地应发生的。\par

% structure: p0062 source=h2 target=subsection reason=relative-heading-rank; base=h2

\subsection{第二十六章　论受迫害者与迫害者。}

因为凡以正直心意致力于任何行动，常将敬畏天主置于思虑之前，并保持对祂坚定不移的信德，不让眼前的恐惧或危险超过对未来赏报之望德者——这样的人，即使暂时经历痛苦考验，也因相信将来更大赏报而轻看其苦难；他们的品格也随过去苦难的严重程度而愈益明亮。另一方面，对于那些或无耻地轻忽正义原则、或拒绝承认至高天主，并且竟敢对忠实奉行祂敬礼的人施加最残酷的侮辱和惩罚者，他们同样未能认识到自己压迫他人的可悲，以及那些即使在苦难中仍保持对天主热忱者所拥有的幸福与祝福——我再说，对于这样的人，他们的军队屡遭屠戮，屡次败逃，他们的备战终归于彻底毁灭与溃败。\par

% structure: p0064 source=h2 target=subsection reason=relative-heading-rank; base=h2

\subsection{第二十七章　迫害如何成为侵略者灾难的缘由。}

从我所描述的这些原因，兴起了惨烈的战争和破坏性的荒废。由此导致生活必需品的匮乏，以及一连串随之而来的苦难：因此，这些不敬之事的作者们，要么遭遇了极度痛苦的惨死，要么苟且于可耻的存在，并承认这比死亡本身更糟，从而仿佛承受了与其罪行严重程度相称的惩罚。因为每个人都依据他被盲目狂怒所驱使、试图（如他所想）击败天主圣意的程度，经历了相应程度的灾祸：以至于他们不仅感受到今世苦难的压力，还被对未来世界惩罚的极其鲜明的恐惧所折磨。\par

% structure: p0066 source=h2 target=subsection reason=relative-heading-rank; base=h2

\subsection{第二十八章 天主拣选\Person{君士坦丁}为降福的仆人}

如今，如此众多的不敬虔压迫着人类，共和政体濒于彻底毁灭，仿佛被某种瘟疫侵蚀，因此需要强而有效的拯救；天主为这些邪恶设计了怎样的解脱和补救呢？（所谓天主，就是指那唯一而真实的天主，拥有全能永恒权能者：一个从天主领受恩惠的人，以最崇高的赞美言辞承认这些恩惠，绝不能视为傲慢。）因此，我本人就是他拣选并认为适于完成其旨意的工具。于是，从遥远的\Place{不列颠}洋，以及按自然法则太阳沉入地平线以下的区域开始，藉着天主权能的帮助，我驱逐并彻底消除了所有盛行的邪恶形式，希望人类通过我的工具受到启迪，得以回归对天主神圣律法的正当遵守，同时我们至为蒙福的信仰藉其全能之手的引导而兴盛。\par

% structure: p0068 source=h2 target=subsection reason=relative-heading-rank; base=h2

\subsection{第二十九章 \Person{君士坦丁}对天主的虔诚表达；以及宣信者的赞美}

我说，在他的手引导之下；因为我愿永远不忘对他恩宠应有的感激。因此，相信这至为卓越的服务已作为特殊恩赐托付于我，我继续前行直至东方地区，这些地区因遭受更严重的灾难，似乎更需要我的更有效的补救。同时我深信，我自己的生命，我所呼吸的每一口气，简言之，我内心深处最隐秘的思想，完全归于至高天主的眷顾。如今我深知，那些真诚追求天上希望、并将这希望视为他们生活的特殊主导原则、定睛于天堂本身的人，无需依赖人的恩惠，反而因他们远离了这尘世生命中低劣邪恶的事物，而享有更高的尊荣。然而，我认为我有责任立刻且彻底地免除所有这样的人一度承受的严酷境遇，以及他们所受的不公惩罚，尽管他们没有任何罪过或应得责任。因为，如果这些人在那些以迫害他们对天主的虔诚为首要目的的统治期间，充分彰显了灵魂的坚忍和恒毅，而在一位身为祂仆人的君王治理下，他们品德的荣耀反而不能更加光明和蒙福，那确实奇怪了。\par

% structure: p0070 source=h2 target=subsection reason=relative-heading-rank; base=h2

\subsection{第三十章 一项法律，免除流放、法庭服役及没收财产}

因此，所有那些因不愿放弃他们全心全意献身的对天主的敬畏和信仰，而背井离乡、移居异地，并因而在不同时期遭受法庭严酷判决的人；以及所有那些被登记在公共法庭的册簿上，尽管过去免于此种职务的人；我命令，这些人现在应向万有的解放者天主献上感恩，因为他们得以恢复祖传的产业和往日的安宁。那些被剥夺了财物，迄今过着悲惨生活，悲叹失掉所有一切的人，也得以重新回到他们先前的家园、家庭和地产，并欢喜接受天主丰厚的恩慈。\par

% structure: p0072 source=h2 target=subsection reason=relative-heading-rank; base=h2

\subsection{第三十一章 同样赐予岛上流放者释放}

此外，朕命令，所有那些违背自己意愿被拘留在岛上的人，应得享本现行法令的恩惠；好使那些至今被崎岖山峦和环绕的海洋屏障包围的人，如今从这阴郁荒凉的孤寂中得释放，能通过重访他们最亲爱的朋友，实现他们最深的愿望。那些在卑贱凄惨的污秽中延长痛苦生命的人，也将他们的复原视作意外之获，从此抛弃一切焦虑，能在我们治下自由、无所畏惧地生活。因为在我们夸耀并相信自己是天主的仆人时，任何人在我们的统治下还能生活在恐惧中，那肯定是极不寻常、简直不可思议的事；我们的使命是纠正其他人的错误。\par

% structure: p0074 source=h2 target=subsection reason=relative-heading-rank; base=h2

\subsection{第三十二章 亦赐予那些被耻役于矿场和公共工程的人}

再者，对于那些被判罚从事矿场苦役或公共工程劳役的人，让他们享受悠闲的甜蜜，以代替这些长期持续的辛劳，从此过一种远为轻松、更符合他们心愿的生活，以安静的休憩换取他们任务的无尽艰辛。若有人丧失了共同的自由特权，或不幸遭受耻辱，让他们各人赶快回到自己的出生地，并以合宜的喜悦恢复他们以前的社会地位，他们仿佛因长期居留异地而与之分离。\par

% structure: p0076 source=h2 target=subsection reason=relative-heading-rank; base=h2

\subsection{第三十三章 论那些从事军事服务的宣信者}

再者，关于那些先前被授予某种军职荣誉、而后因残酷不公的理由被剥夺的人——因为他们宁愿承认对天主的忠诚，也不愿保留所拥有的军衔——我们给予他们完全的自由选择：要么恢复他们先前的职位，倘若他们愿意再次从军；要么在光荣退役后，安享太平。因为那些在危险面前展现出如此宽宏与坚毅的人，理当有权选择享受平静的闲暇，或恢复他们先前的军衔。\par

% structure: p0078 source=h2 target=subsection reason=relative-heading-rank; base=h2

\subsection{第三十四章 妇女居所中被判劳役或奴役的自由人的解放。}

最后，若有人被不公正地剥夺了贵族血统的特权，并被判刑发配到妇女居所从事织麻布等残酷悲惨的劳役，或为公库利益而沦为奴役，不因其高贵出身而豁免；那么，这样的人应恢复他们先前享有的尊荣和应有的品级，从此因自由的祝福而欢欣，过喜乐的生活。同样，若自由人因某种不义、不人道甚至疯狂而沦为奴隶，感受从自由到奴役的突然转变，时常哀叹这不寻常的苦役，应因我们这道法令而重获自由，回到家人身边，追求与自由状态相宜的职业；并让他从记忆中抹去那些他曾感到压迫、与其身份极不相称的劳役。\par

% structure: p0080 source=h2 target=subsection reason=relative-heading-rank; base=h2

\subsection{第三十五章 关于殉道者和宣信者的财产继承，以及遭受流放或财产没收者的财产问题。}

我们也不可忽略那些因各种借口而被剥夺的产业。若有人曾以无畏和坚定的决心参与那崇高神圣的\OriginalTerm{殉道}{martyrdom}争战，并且也被剥夺了财富；或若\OriginalTerm{宣信者}{confessor}也有同样遭遇，他们已为自己赢得了永恒宝藏的希望；或若那些因不肯屈服于迫害者、不肯背弃信仰而被逐出故乡的人遭受了财产损失；最后，若有人虽逃脱了死刑，却被剥夺了世俗财物；我们规定，所有这些人的遗产应转移给其最近的亲属。既然法律明确将此权利赋予最亲近的亲属，那么不难确定这些遗产各自归属何人。显然合理的是，若逝者自然死亡，继承权应属于那些本有最近亲缘关系的人。\par

% structure: p0082 source=h2 target=subsection reason=relative-heading-rank; base=h2

\subsection{第三十六章 教会被宣布为无亲属者的继承人，并确认此类人的无偿赠与。}

但若上述之人——即殉道者、宣信者，或因类似原因被逐出故乡者——没有尚存的亲属按顺序继承其财产；在此情形下，我们规定当地的教会应在每种情况下继承遗产。逝者因其缘故忍受了所有极端苦难的教会成为其继承人，这绝不会对逝者不公。我们认为还需补充一点：若上述任何人以无偿赠与方式捐出部分财产，这些财产的占有权应合理地归于接受者。\par

% structure: p0084 source=h2 target=subsection reason=relative-heading-rank; base=h2

\subsection{第三十七章 应归还土地、园圃或房屋，而非实际收成。}

为使这道法令没有模糊之处，人人都能轻易明白其要求，兹令所有人知晓：若他们现正持有一块土地、房屋、园圃或任何属于上述之人的其他物品，那么如实承认并尽快归还将是件好事且有益。另一方面，即使有人似乎已从此不义占有中获取丰厚利润，我们认为公正并不要求归还这些利润。但他们必须明确声明从中获得了多少收益、来自何处，并恳求我们宽恕此过；以便他们过去的贪欲在某种程度上得以补赎，且愿至高的天主悦纳此补偿作为悔罪的标志，欣然地宽恕其罪。\par

% structure: p0086 source=h2 target=subsection reason=relative-heading-rank; base=h2

\subsection{第三十八章 应如何就此提出请求。}

但那些已成为此类财产之主的人（若允许他们拥有此头衔是恰当或可能的），可能会为自己的行为辩解，声称在那种悲惨景象随处可见的时代——人们被残酷逐出家园、无情杀戮、毫无怜悯地被驱逐；无辜者财产被充公司空见惯，迫害与财产没收不断——他们无法避免占有这些财产。若有人以此类理由辩护并仍坚持其贪婪心态，他们应认识到这种行为将给自身带来惩罚，尤其因为纠正邪恶正是我们事奉至高的天主的特质。因此，将来继续保留那些过去因可怕必要而被迫夺取之物将是危险的；尤其因为我们有责任既通过劝诫也通过警戒事例来抑制贪婪的欲望。\par

% structure: p0088 source=h2 target=subsection reason=relative-heading-rank; base=h2

\subsection{第三十九章 国库必须将土地、园圃和房屋归还给教会。}

国库本身，即使拥有我们所提及的任何物品，也不得保留；但好似不敢对圣教会提出异议似的，它应当公正地将其暂时不义地占有的东西，为了教会的利益而予以归还。因此，我们规定：凡按正义应属于教会的一切，无论财产是房屋、田地、园囿，或任何性质之物，均须按其全部价值与完整性，连同不受减损的所有权，一并归还。\par

% structure: p0090 source=h2 target=subsection reason=relative-heading-rank; base=h2

\subsection{第四十章　殉道者的坟墓与墓地应归教会所有。}

再者，关于那些因保存殉道者遗骸而受尊崇、并继续作为他们光荣离世纪念的地方，我们怎能怀疑它们理应属于教会，或克制自己不发出有关此事的命令呢？因为确实没有比这更好的慷慨，没有比这更令人喜悦或有裨益的劳苦了：即在圣神的引导下如此行事，使那些被不义和邪恶之徒以虚伪的借口所占据的东西，按正义的要求得以归还，并再次确保归于圣教会。\par

% structure: p0092 source=h2 target=subsection reason=relative-heading-rank; base=h2

\subsection{第四十一章　凡购买或接受教会产业作为礼物者，必须归还。}

由于在涵盖所有情况的条款中，忽略那些曾通过购买从国库获得任何此类财产、或将其作为礼物保留的人，是不妥当的；因此，所有如此鲁莽地放纵其贪得无厌之欲的人，当明白：尽管他们敢于进行此类购买，已尽其所能使我们对其失去宽仁，但他们仍不会得不到此宽仁，只要在各情况下可能且合宜。至此已作决定。\par

% structure: p0094 source=h2 target=subsection reason=relative-heading-rank; base=h2

\subsection{第四十二章　恳切劝勉敬拜天主。}

如今，由于最清晰、最令人信服的证据表明，先前压迫全人类的苦难，凭借全能天主的德能，以及祂屡次藉我们施行的谋略与援助，已从世界各地被驱逐；因此，众人无论个人或团体，都应当留意并认真思量，这力量何其伟大，这恩宠何其有效：它已消灭并彻底摧毁了那些最邪恶之人（我如此称呼他们）的一代；为善人恢复了喜乐，并将其遍洒各国；如今更充分保证，能以应有的敬畏尊崇神圣法律，并对那些献身于侍奉该法律的人给予应有的尊重。这些人仿佛从黑暗深渊中升起，以对当前事件进程的启明认识，从此将对其训诫，给予符合其虔诚品格的应有敬畏与尊崇。\par

兹令此法令在朕之东方各省颁布。\par

% structure: p0097 source=h2 target=subsection reason=relative-heading-rank; base=h2

\subsection{第四十三章　君士坦丁的律令如何得以施行。}

皇帝致我的第一封信中即包含这些训令。此律令的规定迅速得以施行，一切进行方式与不久前在暴君残酷统治下所胆大妄为之暴行截然不同。那些依法应得之人，也领受了皇帝的慷慨恩惠。\par

% structure: p0099 source=h2 target=subsection reason=relative-heading-rank; base=h2

\subsection{第四十四章　皇帝擢升基督徒担任官职，并禁止外邦人在此类职位上献祭。}

此后，皇帝继续处理重要事务。首先，他向各省派遣总督，大多为信奉救恩信仰之人；若有任何人显得倾向于外邦崇拜，皇帝即禁止他们献祭。此法律亦适用于那些品级与尊严超越各省总督者，甚至包括占据最高职位、拥有近卫军长官权力者。若他们是基督徒，则可自由按信仰行事；否则，法律要求他们禁绝偶像崇拜之祭祀。\par

% structure: p0101 source=h2 target=subsection reason=relative-heading-rank; base=h2

\subsection{第四十五章　禁止献祭与命令修建教堂的法规。}

此后不久，大约同时颁布了两项法律：其一旨在限制过去各城各国所行之偶像崇拜的可憎之事，规定不得设立偶像、行占卜及其他虚假愚昧之术，或以任何方式献祭。另一法令命令增高原有的祈祷所，并加阔天主的圣殿，仿佛预料到，既然多神论的癫狂已被完全消除，几乎全人类从此都将归附于天主的事奉。皇帝出于自身的虔诚，亲自拟定并撰写这些指示给各省总督；法律还告诫他们不可吝惜金钱，可从皇室国库支取。类似的指示亦写给了各教会的主教；皇帝亦欣愿将其转达于我，此乃他亲笔致我的第一封信。\par

% structure: p0103 source=h2 target=subsection reason=relative-heading-rank; base=h2

\subsection{第四十六章　君士坦丁致尤西比乌斯及其他主教的书信，关于修建教堂，并指示在各省总督协助下，修缮旧堂并扩建新堂。}

得胜者君士坦丁·最大奥古斯都致尤西比乌斯。\par

鉴于那不虔敬的暴虐统治直到现今一直在迫害我们救主的仆人，最亲爱的兄弟，我相信并完全确信：所有教会的建筑物，或因实际疏忽而沦为废墟，或因恐惧时代暴虐之风而未得到充分关注。\par

但如今，自由既已恢复，那毒蛇也因至高天主的眷顾与我们的中介被逐出公共事务的管理，我们确信人人都能看见神圣大能的效力，那些因畏惧迫害或不信而陷入任何错误的人，如今将承认真天主，并在今后采纳那合乎真理与正直的生活。因此，关于你本人所管辖的教会，以及你所熟识的其他教会的主教、司铎和执事，请你劝勉所有人热心关注教堂的建筑，或修缮、扩建现有的教堂，或在必要时建造新的教堂。\par

\Quote{我们也授权你，并通过你授权其他人，向行省总督和近卫军长官索取工程所需之物。因为他们已接到命令，要最勤勉地服从你圣座的命令。愿天主保佑你，亲爱的弟兄。}此训令的副本被发送到各省各教会的司铎；行省总督也接到了相应指示，这道帝国法令迅速得以执行。\par

% structure: p0108 source=h2 target=subsection reason=relative-heading-rank; base=h2

\subsection{第四十七章 他撰写了谴责偶像崇拜的谕令}

此外，皇帝在敬拜天主方面不断进步，向各省居民发出了一封劝诫信，论及他前任掌权者所陷入的偶像崇拜之错误；他在信中雄辩地劝勉臣民承认至高天主，并公开宣称他们忠诚于他的基督作为救主。这封信也是他亲笔所写，我认为有必要为当前这部著作从拉丁文译出，以便我们仿佛亲耳听见皇帝本人在全人类面前发出这些感想。\par

% structure: p0110 source=h2 target=subsection reason=relative-heading-rank; base=h2

\subsection{第四十八章 君士坦丁关于众神崇拜之错误的敕令给各省人民，以善恶之通论开篇}

\Signature{凯旋者君士坦丁，最伟大的奥古斯都，致东方各省人民。}\par

凡为自然之主权法律所包含者，似乎都向所有人传达了关于神圣秩序的预见与智慧的足够观念。任何以知识真道为导向、以达致那目标为念的人，都不会怀疑，正确理性的公正理解，连同自然视觉本身，藉着真正德行的惟一影响，引导人认识天主。因此，任何智者见大众受相反情感支配时，都不会感到惊讶。因为倘若恶行不与它对照，显明乖谬与愚昧的道路，德行之美便无用处，不为人见。因此，其一获赏报，而至高天主亲为另一者施行审判。\par

现在，我将尽可能明确地向你们陈述我自己对未来福乐之盼望的性质。\par

% structure: p0114 source=h2 target=subsection reason=relative-heading-rank; base=h2

\subsection{第四十九章 论君士坦丁虔诚的父亲，以及迫害者戴克里先和马克西米安}

我素来将先前诸皇帝视为我无法同情之人，缘于其品性之野蛮残暴。事实上，我父亲是惟一一贯实行人道义务者，并以可赞的虔敬在一切行动上呼求天主圣父的祝福；而其余诸人，心神不健全，热衷于残暴手段胜过温和措施，他们不加约束地放纵此性情，在其统治整个时期迫害真道。不，他们的恶毒狂怒如此猛烈，以至于在宗教与世俗事务皆处于深厚和平时，他们仿佛点燃了内战的火焰。\par

% structure: p0116 source=h2 target=subsection reason=relative-heading-rank; base=h2

\subsection{第五十章 迫害的起因源于阿波罗的神谕，据说他因\Quote{正义者}而无法发出神谕}

大约在那时，据说阿波罗从一个深邃阴暗的洞穴中，不藉任何人的声音宣告说：世上的\Emph{正义者}妨碍他说出真理，因此从三脚架发出的神谕是虚假的。为此他垂头丧气以示悲伤，哀叹失去神谕之灵将给人类带来的灾祸。但让我们看看此事的后果。\par

% structure: p0118 source=h2 target=subsection reason=relative-heading-rank; base=h2

\subsection{第五十一章 君士坦丁年轻时，从撰写迫害敕令者那里听说\Quote{正义者}就是基督徒}

现在我呼求你，至高天主，为我作证：我年轻时，听见当时罗马皇帝之首——他实在是不幸的、真正不幸的，且为精神错觉所苦——急切地询问侍从，这些世上的正义者是谁；当时一位异教司祭回答说，他们无疑是基督徒。他如同饮下一口蜜酒般急切地接受了这回答，并拔出了原为惩治罪恶而设立的刀剑，指向那神圣无瑕之人。因此，他立即发布了那些血腥的敕令，若可如此说，是以蘸血的剑尖写就；同时命令他的法官们绞尽脑汁发明更新更可怕的刑罚。\par

% structure: p0120 source=h2 target=subsection reason=relative-heading-rank; base=h2

\subsection{第五十二章 施加于基督徒的多种刑罚与折磨形式}

那时，人们真可看到那些可敬的天主崇拜者日复一日，在持续无情的残暴中，遭受最严重的凌辱；那连仇敌也从未轻慢的谦逊品格，竟成为其狂怒同胞的玩物。有哪种火刑、哪种拷打或折磨形式，未曾不加区分年龄或性别地施于所有人？那时，真可说是大地流泪，包罗万物的苍穹因血污而哀悼；白昼之光本身也因这景象而悲伤昏暗。\par

% structure: p0122 source=h2 target=subsection reason=relative-heading-rank; base=h2

\subsection{第五十三章 蛮族和善地接待基督徒}

但这一切的后果是什么呢？怎么，连野蛮人如今也可以夸口，说他们的行为与这些残酷的罪行形成了对比；因为他们接待并仁慈地囚禁了那时从我们中间逃走的人，不仅确保他们免于危险，还让他们自由地实践神圣的宗教。现在，罗马人民背负着那个永久的污点，那是当时被逐出罗马世界、投靠野蛮人的基督徒们烙印在他们身上的。\par

% structure: p0124 source=h2 target=subsection reason=relative-heading-rank; base=h2

\subsection{第五十四章 因神谕而发动迫害的人遭受了怎样的惩罚}

但我何须再详述这些可悲的事件，以及由此弥漫世界的普遍忧伤？犯下这可怕罪行的人如今已不复存在：他们遭遇了悲惨的结局，被投入下界深处永无止境的惩罚。他们自相残杀，没有留下名字或后代。确实，若不是那邪恶的\OriginalTerm{皮提亚神谕}{Pythian oracle}对他们施加了迷惑的力量，这场灾难绝不会降临到他们身上。\par

% structure: p0126 source=h2 target=subsection reason=relative-heading-rank; base=h2

\subsection{第五十五章 \Person{君士坦丁}归光荣于天主，感恩十字圣号，并为教会及人民祈祷}

如今我恳求你，全能的天主，求你怜悯并恩待你东方的万民，即这些省份中因长期苦难而疲惫的子民，并通过你的仆人赐予他们医治。圣洁的天主，我向你，万有的主，献上这祈祷，并非毫无缘由。在你的指引下，我谋划并完成了充满福祉的措施；在你神圣的记号前导下，我率领你的军队走向胜利；每逢公共危险之际，我仍沿用你完善性的同一象征去迎敌。因此，我已将一颗因爱和敬畏而妥善调和的灵魂奉献于你的服务。我真心爱你的名，并敬畏你那所赐予丰富证据、以确认和增强我信德的能力。我急于倾尽所有力量去修复你最神圣的居所，那是那些亵渎和不敬的人以暴力的污染所玷污的。\par

% structure: p0128 source=h2 target=subsection reason=relative-heading-rank; base=h2

\subsection{第五十六章 他祈祷众人皆成为基督徒，但不强迫任何人}

我自己的愿望是，为世界的共同福祉和全人类的利益，愿你的子民享受和平与无扰的和睦生活。因此，让那些仍以谬误为乐的人，也享有与相信者同等的和平与安宁。因为或许这种平等权利的恢复会将他们引向正道。谁也不可骚扰他人，让各人按自己灵魂所愿而行。但愿心智健全的人确信这一点：只有那些被你召唤去信赖你神圣律法的人，才能过圣洁和纯洁的生活。至于那些将远离我们的人，如果他们愿意，就让他们拥有他们的谎言庙宇吧；\Emph{我们}却有你真理的光辉殿宇，那是你赐给我们作故乡的。然而，我们祈祷他们也能领受同样的祝福，从而体验那由情感合一所带来的由衷喜乐。\par

% structure: p0130 source=h2 target=subsection reason=relative-heading-rank; base=h2

\subsection{第五十七章 他归光荣于天主，天主藉其圣子赐光明给迷误之人}

确实，我们的敬拜并非新奇或近期的事，而是你为你的适当尊荣所定的，从我们相信这宇宙体系最初建立之时起。虽然人类深深堕落，被多种谬误所迷惑，但你藉你圣子的位格彰显了纯净的光明，使邪恶的势力不能完全得胜，并藉此向所有人见证了你。\par

% structure: p0132 source=h2 target=subsection reason=relative-heading-rank; base=h2

\subsection{第五十八章 他再次为他治理宇宙而光荣他}

你的作为向我们保证了此事的真实性。是你的能力除去我们的罪过，使我们成为忠信者。太阳和月亮有它们固定的轨道。众星围绕这地球运行无误。季节的更替遵循无误的规律。大地的坚固结构由你的圣言奠定；风在指定的时间受到推动；水流不息，海洋被不可移动的屏障所限制，凡在天地海之内的一切，都为了奇妙而重要的目的而设计。\par

若不是这样，若不是一切都由你意志的决定所规范，如此巨大的多样性、如此多种的权力划分，无疑会给整个人类和它的事务带来毁灭。因为那些彼此争斗的力量，它们对人类所怀的敌意，哪怕现在肉眼不见，也会发展到更致命的地步。\par

% structure: p0135 source=h2 target=subsection reason=relative-heading-rank; base=h2

\subsection{第五十九章 他归光荣于天主，作为善的永恒导师}

愿丰盛的感谢归于你，全能的天主、万有的主！因我们的本性从人不同的追求中越被认识，你神圣教义的训导就越得到那些思想正直、真诚献身于真德行之人的确认。至于那些不允许自己从谬误中被治愈的人，愿他们不将此归咎于别人，只归咎于自己。因为那具有至高无上疗效的良药，公开地摆在所有人面前。但愿没有人去伤害那经验本身证明为纯洁无瑕的宗教。因此，从今以后，让我们共同享有我们所能得到的特权，即和平的恩赐，努力保持我们的良心远离一切相反之事。\par

% structure: p0137 source=h2 target=subsection reason=relative-heading-rank; base=h2

\subsection{第六十章 诏令结尾的劝诫：不应有人骚扰邻人}

再次，愿没有人用他自己因确信真理而领受的东西，去损害他人；但愿每个人，如果可能，将他所理解并认识到的用于邻人的益处；否则，就放弃这种尝试。因为自愿为不朽而奋斗是一回事，因惧怕惩罚而强迫他人这样做是另一回事。\par

\Quote{这些话是我们的；我们之所以在这些话题上说得比我们惯常的宽仁所会指示的更多，是因为我们不愿掩饰或背叛真信仰；尤其因为我们知道有些人说异教神殿的仪式和黑暗的权势已被完全移除。我们本当诚挚地向所有人推荐这种移除，若不是那些邪恶谬误的反叛精神仍在某些人的心中顽固地扎根，以致使人对全人类普遍回归真理之路不抱希望。}\par

% structure: p0140 source=h2 target=subsection reason=relative-heading-rank; base=h2

\subsection{第六十一章 争论如何因与\Person{阿里乌}有关之事在\Place{亚历山大里亚}发起}

皇帝就这样，如同一位强大的天主使者，通过他自己的信函向所有省份讲话，同时警告他的臣民远离迷信的谬误，并鼓励他们追求真正的虔敬。但正当他为此举的成功而欢欣期待时，他听到了一个极为严重的骚乱侵犯了教会和平的消息。他深感忧虑地听到这个消息，并立刻试图为这邪恶寻找补救之法。这个骚乱的起源可以这样描述：天主的子民确实处于繁荣状态，善工丰盛。没有外来的恐怖袭击他们，但因天主的恩宠，一片明亮而极其深沉的和平环绕着他的教会。然而，与此同时，嫉妒之灵正伺机摧毁我们的福乐；它起初悄悄潜入，但不久就在圣徒的集会中横行。最终它波及了主教们本身，以嫉妒地维护神圣真理教义为借口，使他们彼此愤怒敌对。由此，仿佛从一点火星燃起了一场大火，最初起源于亚历山大里亚教会，蔓延至整个埃及、利比亚和上底比斯。最终它的蹂躏扩展到了帝国的其他省份和城市；不仅可见教会的主教们在言辞的争斗中互相冲突，而且民众本身也完全分裂，一些人依附一派，另一些人依附另一派。甚至，这些事的丑闻变得如此臭名昭著，以至于默感教义的神圣之事在不信者的剧场中也遭到最可耻的嘲笑。\par

% structure: p0142 source=h2 target=subsection reason=relative-heading-rank; base=h2

\subsection{第六十二章 关于同一\Person{阿里乌}和\OriginalTerm{梅利提安人}{Melitians}}

在亚历山大里亚，有些人就最高问题顽固地争执。在整个埃及和上底比斯，其他人因更早的争论而不和：因此教会到处被分裂所扰乱。既然身体如此患病，整个利比亚也感染了这传染病；其余更远的省份也受到同样混乱的影响。因为亚历山大里亚的争论者向各省的主教派出使者，他们分别站在一方成为党徒，并参与了同样的分裂精神。\par

% structure: p0144 source=h2 target=subsection reason=relative-heading-rank; base=h2

\subsection{第六十三章 \Person{君士坦丁}如何派遣使者和信件论和平}

皇帝一得知这些事实——他听到时心中十分悲伤，视之为亲自加于自身的灾祸——他立即从随行的基督徒中选出一位他深知其信仰的清醒和纯正、并且在此之前已因其宗教宣告的勇敢而著称的人，派他去在亚历山大里亚的争论各方之间谈判和平。他还让他带去一封非常必要且合宜的信，给这场争端的始作俑者；这封信作为他看顾天主子民之警觉的样本，或许适合引入我们这部关于他生平记述之中。其内容如下。\par

% structure: p0146 source=h2 target=subsection reason=relative-heading-rank; base=h2

\subsection{第六十四章 \Person{君士坦丁}致\Person{亚历山大}主教和\Person{阿里乌}司铎的信}

\Signature{胜利者\Person{君士坦丁}，最大奥古斯都，致\Person{亚历山大}和\Person{阿里乌}}\par

我呼求那位天主作证——我也理应如此——祂是我努力的帮助者，也是万民的保全者，说我承担那我现在已经履行的职分有双重原因。\par

% structure: p0149 source=h2 target=subsection reason=relative-heading-rank; base=h2

\subsection{第六十五章 他不断为和平忧虑}

我的计划首先是，将所有民族对神祇所形成的不同判断带入一种可以说是稳固的一致状态；其次，恢复当时在一种严重疾病的邪恶力量下受苦的世界体系。抱着这些目标，我试图通过思想的隐秘眼睛完成一个，而通过军事权威的力量矫正另一个。因我意识到，如果我能成功地在所有天主仆人中建立我所希望的共同和谐，一般事务的进程也将按照他们所有人虔诚的愿望而发生相应的改变。\par

% structure: p0151 source=h2 target=subsection reason=relative-heading-rank; base=h2

\subsection{第六十六章 他也调整了在\Place{非洲}兴起的争论}

发现整个非洲因那些以轻率无聊精神妄图将人民的宗教分裂为不同派别之人的影响，而被一种无法容忍的疯狂愚蠢之灵渗透，我急于制止这种混乱，但找不到与时机相称的其他补救方法，除了派遣你们中的一些人协助恢复争论者之间的相互和谐，在我移除了那个阻拦你们神圣会议的不法判决的共同人类敌人之后。\par

% structure: p0153 source=h2 target=subsection reason=relative-heading-rank; base=h2

\subsection{第六十七章 宗教始于\Place{东方}}

因为既然神圣之光的力量和神圣敬拜的法律——最初藉着天主的恩宠，从东方的怀抱（可以说是）出发，以它们神圣的光辉照亮了世界——我自然相信你们会是首先促进其他民族得救的人，并决心以全部思想精力和勤勉探询来寻求你们的帮助。因此，一旦我取得了对敌人的决定性胜利和无庸置疑的凯旋，我的首要探询就是关于那件事，我感觉那事具有最高利益和重要性。\par

% structure: p0155 source=h2 target=subsection reason=relative-heading-rank; base=h2

\subsection{第六十八章 因纷争而忧伤，他劝告和平}

但是，啊，天主的荣耀眷顾！当我听说你们之间存在着比该国中继续存在的更为严重的分裂时，我的耳朵，不，我的内心受到了何等深的创伤！以致我本希望借助你们的帮助为别人的错误寻求补救，而你们却处于比他们更需要医治的状态。然而，在仔细调查这些分歧的起源和基础之后，我发现其原因确实微不足道，完全不值得如此激烈的争论。因此，我感到不得不以这封信向你们说话，同时呼吁你们的团结和明智，我呼求天主眷顾在此任务中协助我，以和平使者的身份打断你们的分歧。这是有理由的：因为如果我期望在更高力量的帮助下，能够通过明智地诉诸听众的虔诚情感，毫不困难地将他们唤回到更好的精神，即使分歧的原因更大，我怎能不承诺自己更容易、更迅速地解决这个分歧，当妨碍普遍情感和谐的原因本身是微不足道且无关紧要时？\par

% structure: p0157 source=h2 target=subsection reason=relative-heading-rank; base=h2

\subsection{第69章 \Person{亚历山大}与\Person{亚略}之间争论的起源，以及这些问题不应被讨论。}

\Emph{我明白}，那么，当前争论的起源是这样的。当你，\Person{亚历山大}，要求司铎们各自对神圣法律中的某段经文持何种意见时，或者更确切地说，我问你关于一个无益的问题，那时，你，\Person{亚略}，轻率地坚持那些根本不应该被想到的事情，或者即使想到，也应该深埋于沉默之中。因此，你们之间产生了分歧，共融被撤回，神圣的子民分裂成不同的派别，不再保持一个身体的合一。现在，因此，你们俩都要表现出同等的克制，接受你们的同仆正义地给出的建议。那么这个建议是什么？首先提出这样的问题，或对它们作出回答，是错误的。因为那些讨论点并非由任何法律权威所规定，而是由错误的闲暇所培养的好争论精神所暗示，即使它们仅旨在作为智力练习，也应该被限制在我们自己的思想领域，不应草率地在民众集会上提出，也不应轻率地托付给大众的耳朵。因为有多少人能够准确理解或充分解释如此崇高和深奥性质的题目？即使有人完全胜任这一点，他能说服多少人？或者，在处理如此微妙的问题时，谁能确保自己不会危险地偏离真理？因此，我们在这些情况下应该慎言，以免，因为我们自己由于自然能力的薄弱而无法清楚解释面前的主题，或者另一方面，因为听众理解的迟钝而无法准确理解我们所说的话，由于这两种原因之一，人民就会陷入亵渎或分裂的境地。\par

% structure: p0159 source=h2 target=subsection reason=relative-heading-rank; base=h2

\subsection{第70章 劝勉合一。}

因此，让那轻率的提问和轻率的回答都得到你们彼此的宽恕。因为你们分歧的原因并非神圣法律中的任何主要教义或诫命，也没有在你们中间兴起任何关于敬拜天主的新异端。你们实在具有同一判断；因此你们大可以联合在共融与交往之中。\par

% structure: p0161 source=h2 target=subsection reason=relative-heading-rank; base=h2

\subsection{第71章 在本身微不足道的事情上不应有争论。}

因为你们继续为这些微小且极不重要的问题争辩，让天主这么大一部分子民受你们判断的指引，而你们自己却彼此分裂，这实在不相宜。我相信，这不仅是不得体，而且是恶事。然而，我要用一个小比喻来提醒你们。你们知道，尽管哲学家们都遵循一个体系，但他们也常在某些问题上产生争执，或许在知识程度上也有差异；然而，他们因共同教义的联合力量而被召回思想的和谐。若果真如此，那么你们这些至高天主的仆役，岂不更应当在同一宗教的宣认上同心一意吗？但让我们更加深思熟虑、更加密切地检视我所说的，看看是否应当因为我们自己之间一些琐碎而愚蠢的言辞差异，就使弟兄彼此以敌人相待，并因你们在如此细微且完全无关紧要之点上争吵不休的人，而使这庄严的主教会议被亵渎的分裂所撕裂？这是粗俗的，更带有孩童般无知的特点，而不是与司铎和明智者的智慧相符。让我们心甘情愿地远离这些魔鬼的诱惑。我们伟大的天主和共同的救主将同一光明赐给了我们众人。请允许我，祂的仆人，在祂眷顾的引导下，成功完成我的任务，使我能够通过我的劝勉、勤勉和恳切的告诫，将祂的子民召回共融与友爱之中。因为正如我所说的，你们对于我们的宗教有同一的信德和同一的见解，并且神圣诫命在各方面都命令我们全体保持和谐的精神，那么，既然那导致你们之间轻微分歧的事并不影响整体的有效性，就不该在你们中间引起分裂或分离。我这样说，绝非要强迫你们对这本属闲散的议题（无论其真正性质如何）完全统一意见。因为你们的主教会议的尊严可以保持，你们全体团体的共融也能维持不破裂，尽管你们在无关紧要的事上可能存在很大分歧。因为我们并非所有人在每件事上都意见一致，也没有一种性情和判断是人人共有的。因此，关于天主眷顾，愿你们有一样的信德，一样的见解，对天主的判断一致。至于你们在无关紧要或毫无意义的问题上的微妙辩论，即使无法在情感上取得一致，也当将这些分歧托付给自己的思想和念头的隐秘保管。如今，愿共同友情的珍贵、对真理的信德、对天主的尊敬以及遵守祂法律的义务在你们中间坚定不移。那么，恢复你们彼此之间的友谊、友爱和关怀吧；把惯常的拥抱还给民众；你们自己也当净化灵魂，好比重新彼此承认。因为常常发生这样的事：当和解因消除敌意的原因而实现时，友谊甚至比先前更加甜美。\par

% structure: p0163 source=h2 target=subsection reason=relative-heading-rank; base=h2

\subsection{虔诚关怀的急切之情使他流泪；因这些事，他原定前往东方的行程被推迟。}

\Quote{那么，还我宁静的白天和不受搅扰的夜晚，使不暗淡之光的喜乐、宁静生活的欢愉从此成为我的份。否则我必然要不断流泪哀悼，无法在平安中度过余生。因为天主的子民，我同为仆人者，竟这样被一股不合理且有害的争辩精神彼此分裂，我又如何能保持心灵的安宁呢？我且向你们证明，我为此事有多么悲伤。不久以前，我到了\Place{尼科美底亚}，本想立即从该城前往东方。正当我赶赴你们那里，已走完大半路程时，这事的消息改变了我计划，使我不必被迫亲眼看见那几乎连听到都几乎无法忍受的事。因此，从今以后，愿你们统一的见解为我开启通往东方地区的道路——你们的纷争已对我关闭了这条路——并让我尽快看见你们和所有其他人民一同欢庆，向天主献上赞美和感谢的言词，为普世和平与自由的恢复而感恩。}\par

% structure: p0165 source=h2 target=subsection reason=relative-heading-rank; base=h2

\subsection{收到这封信后，争执仍未平息。}

虔诚的皇帝就这样通过上述信件努力促进天主教会和平。而受托送信的那位卓越之人不仅传达了信件本身，也支持了送信者的观点；因为正如我所说，他在各方面都是个虔诚的人。然而，邪恶过于深重，非一封信所能补救，以致争论双方的激烈程度不断加剧，祸害的影响扩展到所有东方省份。这些都是嫉妒和某个恶灵造成的，它们以嫉妒的眼光看待教会的兴旺。\par

\SourceRecord{Translated by Ernest Cushing Richardson.；Nicene and Post-Nicene Fathers, Second Series；Vol. 1.；Edited by Philip Schaff and Henry Wace.；Buffalo, NY: Christian Literature Publishing Co.,；1890.；<http://www.newadvent.org/fathers/25022.htm>.}

% structure: p0001 source=h1 target=section reason=page-title

\section{\WorkTitle{君士坦丁传（第三卷）}}

% structure: p0002 source=h2 target=subsection reason=relative-heading-rank; base=h2

\subsection{第一章 \Person{君士坦丁}的虔诚与迫害者的邪恶之比较}

如此，那憎恶良善的恶神，因嫉妒教会所享的福乐，在和平与喜乐期间，仍在她内部激起纷争的风暴。与此同时，蒙天主眷顾的皇帝并未忽略应尽之责，反而在一切行为上，与那些残酷暴君近来的暴行形成鲜明对比，从而战胜了所有反对他的仇敌。首先，暴君们自身背离了真天主，用各种强制手段迫使他人崇拜假神；而\Person{君士坦丁}不仅以言词，更以行动使世人确信这些假神仅为虚妄，劝勉他们承认独一的真天主。他们曾以亵渎之言嘲弄基督，而\Person{君士坦丁}却以他们嘲弄的对象作为自己的护佑，并以救主苦难的标记为荣。他们迫害、驱逐基督的仆人，使其无家可归；而\Person{君士坦丁}却将他们一一召回，使他们重返故乡。他们使这些仆人蒙受羞辱；\Person{君士坦丁}却使他们的处境在众人眼中变得荣光可羡。他们无耻地抢劫并出卖虔诚之人的财物；\Person{君士坦丁}不仅补偿了这些损失，还以丰厚的赏赐使他们更加富足。他们通过书面法令散布对教会主教的中伤诽谤；\Person{君士坦丁}反而以个人荣誉的标记尊崇这些人，并通过谕令和法规使他们比先前更受尊敬。他们彻底拆毁并夷平圣堂；\Person{君士坦丁}则下令扩大现存者，并以恢弘的规模、动用国库资金建造新的圣堂。他们下令焚烧并彻底销毁神圣的经卷；\Person{君士坦丁}则命人大量抄写这些经卷，并以宫廷的经费加以华丽装饰。他们严格禁止各地主教在任何场合召开会议；而\Person{君士坦丁}却从各省将他们召至宫廷，迎接他们进入皇宫，甚至自己的内室，认为他们配得分享他的家宴。他们以祭品尊崇邪魔；\Person{君士坦丁}揭示了他们的谬误，并将那些已无用的祭品材料分给能善加利用的人。他们下令豪华装饰异教神庙；\Person{君士坦丁}却将那些曾是迷信崇拜主要对象的庙宇彻底夷平。他们使天主的仆人遭受最耻辱的刑罚；\Person{君士坦丁}则惩罚迫害者，奉天主的名义给予他们公正的惩治，同时恒常敬礼圣殉道者的纪念。他们将天主的敬拜者逐出皇宫；\Person{君士坦丁}却始终完全信任他们，深知他们比任何人都更忠信良善。那些被贪婪所害的暴君，仿佛自愿忍受坦塔罗斯的折磨；而\Person{君士坦丁}以帝王的慷慨打开所有宝库，以丰厚而高贵的胸怀分施恩赐。暴君们为掠夺或没收受害者的财物而杀人无数；而在\Person{君士坦丁}统治期间，正义之剑处处闲置，各省的人民和地方官员更多地受父权式的治理，而非强制。凡适当思考这些事实者，必定认为一个崭新而清新的生存时代开始显现，一种前所未见的光明忽然从黑暗中照耀人类；人人必须承认，这些完全是天主的作为，祂兴起这位虔诚的皇帝以抵抗众多不敬之徒。\par

% structure: p0004 source=h2 target=subsection reason=relative-heading-rank; base=h2

\subsection{第二章 再论\Person{君士坦丁}的虔诚，以及他对十字圣号的公开见证}

当我们思量那些不义之事史无前例，他们胆敢对教会施行的暴行是世间任何时代所未闻的，天主便亲自为我们带来全新的景象，并由此成就了迄今未曾记载或观察到的事迹。还有什么奇迹比我们这位皇帝的美德更奇妙呢？天主的智慧将他作为恩赐赐予人类。的确，他在众人面前始终大胆地、公开地为天主的基督作证；他非但毫不退缩于公开宣认基督圣名，反而渴望向所有人表明，他视此为最高荣誉，时而将这十字圣号印在脸上，时而以它作为引领他走向胜利的旗帜而夸耀。\par

% structure: p0006 source=h2 target=subsection reason=relative-heading-rank; base=h2

\subsection{第三章 关于他的画像：上方有十字架，下方有龙}

除此以外，他命人将十字圣号绘于一块高大的画板上，置于其宫殿门廊前，使众人可见，该标记置于他头部上方，而下方的则是那可憎而凶残的人类仇敌——他借着不敬之徒的暴政摧残了天主的教会——以龙的形态坠入毁灭的深渊。因为神圣的经卷在先知书中描述他为龙和弯曲的蛇；因此，皇帝公开描绘了这条龙的形象，在其自身和子女脚下，被箭刺穿，坠入深海。\par

他用这种方式意图描绘人类隐秘的仇敌，并指示他因放置在头上的救恩性战利品而被交付于灭亡的深渊。这幅寓意画就这样通过绘画的色彩传达出来；我对皇帝的伟大智慧充满惊奇，他仿佛受到神圣的启示，表达了先知们关于这个怪兽所预言的，说：\Quote{天主必用他大而强、可畏的刀剑攻击那龙，就是那飞行的蛇；并要毁灭海中的龙。}这正是皇帝在以上描述的图画中真实而忠实地呈现的。\par

% structure: p0009 source=h2 target=subsection reason=relative-heading-rank; base=h2

\subsection{第四章　关于\Person{阿里乌}在\Place{埃及}所挑起争端的进一步说明}

君王乐于从事此类事务；但是那扰乱\Place{亚历山大}诸天主教会平安的嫉妒之灵的影响，连同\Place{底比斯}和\Place{埃及}的分裂，继续使他心中颇不安宁。因为事实上，在每个城市中，主教与主教陷入顽固的冲突，民众与民众互相敌对；几乎如同传说中的\OriginalTerm{碰撞}{Symplegades}，彼此剧烈碰撞。甚至有些人完全失去理性，做出鲁莽而粗暴的行为，甚至侮辱皇帝的雕像。这种情况几乎无法激起他的愤怒，反而使他心中忧伤；因为他深深哀叹这些心神错乱者所表现的愚妄。\par

% structure: p0011 source=h2 target=subsection reason=relative-heading-rank; base=h2

\subsection{第五章　关于庆祝复活节之分歧}

但在此之前，另一种最恶性的混乱早已存在，并长期折磨着教会；我指的是关于救恩性的复活节庆日的分歧。因为一方坚持应遵守犹太人的习俗，另一方则主张应遵守确切的周期重复，而不服从那些错误且与福音恩宠无关之人的权威。\par

因此，民众在各地由此产生分歧，神圣宗教的礼仪长期混乱（以致在庆祝同一复活节庆节的时间上的不同见解，使守节者之间产生极大分歧，一些人以斋戒和刻苦折磨自己，另一些人则致力于节日的轻松），没有人能想出补救这一邪恶的方法，因为争论在双方之间均衡持续。治愈这些分歧，对全能的天主独一神而言是容易的；而君士坦丁似乎是在世上唯一能够为这一善果充当祂工具的人。因为当他得知我所描述的事实，并察觉他给\Place{亚历山大}基督徒的信函未能产生应有的效果时，他立刻振奋精神，宣称他必须全力进行这一场反对那扰乱教会和平的隐秘仇敌的战争。\par

% structure: p0014 source=h2 target=subsection reason=relative-heading-rank; base=h2

\subsection{第六章　他如何下令在\Place{尼西亚}召开公会议}

于是，仿佛要调动一支神圣的军队对抗这个仇敌，他召集了一次公会议，并以表达他对主教们尊敬的信函邀请他们迅速出席。这不仅仅是下达一道简单的命令，而且皇帝的善意对执行大有帮助：因为他允许一些人使用公共交通工具，另一些人则获得充足的马匹供应。为这次会议选定的地点，即是\Place{比提尼亚}的\Place{尼西亚}城（其名称源自\Quote{\Emph{\OriginalTerm{胜利}{Victory}}}），也恰到好处。因此，皇帝的谕令一经普遍传达，所有人都极其乐意地赶往那里，仿佛要在赛跑中争先恐后；因为他们受对会议幸福结局的期待、对当下和平的希望，以及目睹如此可钦佩的皇帝这一新奇之事的渴望所激励。当他们都聚集在一起时，显然可见这次行动是天主的工作，因为那些不仅在观点上而且在地域、地点和民族上原本相隔最远的人，如今被聚集在一起，被容纳在一座城墙之内，仿佛形成了一个由各种精选花朵组成的巨大司铎花环。\par

% structure: p0016 source=h2 target=subsection reason=relative-heading-rank; base=h2

\subsection{第七章　论有各国主教出席的公会议}

果然，来自\Place{欧罗巴}、\Place{利比亚}和\Place{亚细亚}的所有教会中最杰出的天主的司铎们齐集于此。一座祈祷之所，仿佛是神圣地扩展了，足以同时容纳\Place{叙利亚}人和\Place{基里基雅}人、\Place{腓尼基}人和\Place{阿拉伯}人、来自\Place{巴勒斯坦}的代表、以及来自\Place{埃及}的人；\Place{底比斯}人和\Place{利比亚}人，连同来自\Place{美索不达米亚}地区的人。一位\Place{波斯}的主教也出席了这次会议，甚至一位\Place{斯基泰}人也未在人数中缺席。\Place{本都}、\Place{迦拉达}、\Place{旁非利亚}、\Place{卡帕多细亚}、\Place{亚细亚}和\Place{夫黎基雅}提供了他们最杰出的主教；而居住在\Place{色雷斯}和\Place{马其顿}、\Place{亚该亚}和\Place{埃皮鲁斯}最偏远地区的人也仍然出席。甚至来自\Place{西班牙}本身，一位名声广传的人也作为个人坐在这个伟大的集会中。帝都的主教因极度年老未能参加，但他的司铎们出席并代替了他的位置。君士坦丁是历代第一位用和平之纽带捆绑如此花环，并作为对所有仇敌取得胜利的感恩祭献给其救主的君王，从而在我们这个时代展示了宗徒团体的相似物。\par

% structure: p0018 source=h2 target=subsection reason=relative-heading-rank; base=h2

\subsection{第八章　会议如同\WorkTitle{宗徒大事录}中那样由各国人士组成}

因为据说\BibleRef{宗 2:5}，在宗徒时代，有\Quote{来自天下各国的虔诚犹太人}聚集；其中有帕提亚人、玛待人、厄兰人和住在美索不达米亚、犹太、卡帕多细亚、本都、亚细亚、夫黎基雅、旁非利亚、埃及及靠近基勒乃的利比亚一带的人，以及侨居的罗马人、犹太人和归依犹太教的人、克里特人、阿拉伯人。但那次集会较小，因为并非所有参与者都是天主的仆人；而此次聚会中，主教的人数超过二百五十位，随行的司铎和执事，以及辅祭和其他仆从的群众则完全无法计数。\par

% structure: p0020 source=h2 target=subsection reason=relative-heading-rank; base=h2

\subsection{第九章 论二百五十位主教的德行与年龄}

在这些天主的仆人中，有的以智慧和口才著称，有的以生活的庄重和性格的坚韧著称，还有的集所有这些恩宠于一身。他们中有年高德劭者，也有年轻且心智成熟者，还有刚刚开始牧职生涯者。皇帝下令每日为他们提供充足的供应。\par

% structure: p0022 source=h2 target=subsection reason=relative-heading-rank; base=h2

\subsection{第十章 皇宫中的会议。\Person{君士坦丁}进入并坐于会场。}

当指定的日子到来，会议为最终解决争议问题而召开时，每位成员都聚集在皇宫的中央大殿，该殿似乎比其他建筑更为宏伟。殿内两侧按秩序排列着许多座位，与会者按等级就座。当全体与会者井然有序地坐定后，一片肃静，期待皇帝的到来。首先，他的三位近亲依次进入，随后另有一些人也在他之前进来，并非通常伴随他的士兵或卫兵，而只是信仰中的朋友。此刻，当信号表明皇帝进入时，全体起立，最后他亲自穿过会场中央，如同来自天主的某种天上使者，身穿仿佛闪烁光芒的衣服，反射着紫袍的炽烈光辉，并以黄金和宝石的璀璨光彩为饰。这便是他外在的容貌；至于他的内心，显然他以虔诚和敬畏天主为特征。这从他低垂的目光、脸上的红晕和步态中可以看出。至于其他个人优点，他在身高和体态之美以及威严的气度和不可战胜的力量与活力上超越了所有在场者。所有这些优点，加上温和的举止和符合皇位的宁静，表明他心智品质的卓越非言语所能赞美。当他走到座位上方时，起初他站着，待一把精雕的金制矮椅为他放好后，他等待主教们向他示意，然后坐下，之后全体与会者也同样坐下。\par

% structure: p0024 source=h2 target=subsection reason=relative-heading-rank; base=h2

\subsection{第十一章 主教\Person{欧瑟伯}发言后的会场寂静}

坐在会场右侧首位的主教随即起身，向皇帝致辞，以感谢全能天主的口吻发表了一篇简短的演说。他坐下后，全场寂静，所有人都注视着皇帝；皇帝欣然环顾会场，面带愉悦，收敛心神，以平静温和的语气说出了以下话语。\par

% structure: p0026 source=h2 target=subsection reason=relative-heading-rank; base=h2

\subsection{第十二章 \Person{君士坦丁}论和平的会议致辞}

\Quote{最亲爱的朋友们，我最大的愿望曾经是享受你们共同在场的景象；如今这个愿望实现了，我感到必须感谢万有的天主，因为在他所有的恩惠之外，他赐予了我高于一切的祝福，允许我看到你们不仅全部聚集在一起，而且都怀着共同的和谐情意。因此我祈求，今后没有恶意的仇敌来干扰我们的幸福；我祈求，既然暴君们不虔敬的敌意已藉着天主我们的救主的力量永久消除，那喜好邪恶的灵不再发明其他手段使神圣法律遭受亵渎的诽谤；因为据我判断，天主教会内部的争斗比任何一种战争或冲突都更邪恶、更危险；在我看来，这些分歧比任何外在的麻烦都更令人痛苦。因此，当我凭天主的旨意和合作战胜了我的敌人后，我以为除了感谢他，并与那些藉我之手被他恢复自由的人同乐之外，再没有别的事可做了；但当我听到我始料未及的消息——即你们纷争的消息——时，我判断这不是次要的事，而是热切希望藉我之手为这一邪恶也找到补救办法，于是立即派人召你们前来。现在我因看到你们的集会而欢喜；但我觉得，只有当我看到你们全体意见一致，并且你们所有人都怀着和平与和谐的共同精神时——这本是你们作为献身于天主事奉的人所应当向他人推荐的——我的愿望才算完全实现。因此，亲爱的朋友们，不要迟延；天主的仆人们，我们共同的主和救主的忠实仆人们，不要迟延：从此刻起，抛弃你们之间存在的分裂的原因，通过拥抱和平的原则来消除争论的困惑。因为如此行事，你们同时将最悦乐至高天主，并将赐予我——你们的同仆——极大的恩惠。}\par

% structure: p0028 source=h2 target=subsection reason=relative-heading-rank; base=h2

\subsection{第十三章 他如何引导意见分歧的主教们走向和谐一致}

皇帝以拉丁语说完这些话后，由另一人翻译，便允许与会的主教们发表意见。于是有些人开始指责邻座者，后者则为自己辩护并反唇相讥。如此，各方提出了无数主张，从一开始便爆发了激烈的争论。尽管如此，皇帝仍耐心地听取所有人的意见，专注地接纳每一条提案，并时而依次协助各方的论证，逐渐使得最激烈的争辩者也倾向于和解。同时，他以其和蔼可亲的态度和对希腊语（他并非全然不懂）的运用，显得格外迷人可亲，说服一些人，以理由说服另一些人，赞扬发言得体者，并敦促所有人达成共识，最终成功使他们就所有争议问题达成一致的意见和判断。\par

% structure: p0030 source=h2 target=subsection reason=relative-heading-rank; base=h2

\subsection{第十四章　关于信德与复活节庆期的全体一致宣言}

结果是，他们不仅在信德上团结一致，而且就庆祝这带来救恩的复活节的日期也达成了全体共识。那些经全体决议批准的事项也被记录成文，并由每位成员签字。于是，皇帝相信他已借此战胜了教会的仇敌，获得了第二次胜利，便着手为天主举行一次凯旋的庆典。\par

% structure: p0032 source=h2 target=subsection reason=relative-heading-rank; base=h2

\subsection{第十五章　君士坦丁在其即位二十周年庆典之际如何款待主教们}

大约此时，他完成了执政的第二十个年头。为此，各省人民普遍举行了公共庆典，而皇帝本人则邀请并宴请了他所促成和解的这些天主的仆人，并藉着他们，如同将合宜的祭献献给天主。没有一位主教缺席这场皇家宴会，其场面之壮观非言语所能描述。禁卫军的分遣队及其他部队手持出鞘的刀剑环绕着宫殿入口，而天主的仆人们无惧地穿过这些士兵，进入皇宫的最内层，其中一些人是皇帝同桌的宾客，另一些人则斜倚在两侧排列的长榻上。有人可能会认为，这是一个预表基督国度景象的画面，如梦似幻而非现实。\par

% structure: p0034 source=h2 target=subsection reason=relative-heading-rank; base=h2

\subsection{第十六章　赠予主教的礼物及致全体人民的信函}

在举行完这场辉煌的庆典之后，皇帝亲切地接待了所有客人，并在已施恩惠之外，又慷慨地根据各人的等级亲自赠送礼物。他还通过一封亲笔信，向未与会者通报了主教会议的各项决定。我也要将这封信如同刻在纪念碑上一般，插入我这篇记述他生平的叙述中。其内容如下：\par

% structure: p0036 source=h2 target=subsection reason=relative-heading-rank; base=h2

\subsection{第十七章　君士坦丁就尼西亚大公会议致各教会的信函}

\Signature{君士坦丁·奥古斯都，致各教会。}\par

鉴于国运昌盛充分证明了天主对我们的莫大恩宠，我已判定，我的首要努力目标应当是：使组成天主教会的有福民众中，能保持信德的统一、真诚的爱德以及对全能天主敬拜的共同情感。然而，除非所有，或至少大部分主教能聚集一堂，并讨论一切有关我们至圣宗教的细节，这一目标才能有效且确定地实现；为此，我们召开了尽可能多的人参加的会议，我本人也以你们中一员的身份出席（我绝不会否认我最引以为乐的事，即我是你们的同仆），每个问题都得到了适当而充分的审查，直至那洞悉万物的天主所赞许、并趋向合一与和谐的判断被显明出来，从而在信德上不再留下任何进一步讨论或争议的余地。\par

% structure: p0039 source=h2 target=subsection reason=relative-heading-rank; base=h2

\subsection{第十八章　他论述他们关于复活节庆期的共识，并反对犹太人的做法}

在这次会议中，讨论了关于至圣逾越节的问题，全体出席者一致决议，这一庆节应在所有地方由所有人同一天庆祝。因为还有什么比这更合宜、更光荣的事：我们赖以寄托永生希望的这庆节，应按照同一确定的规定和安排，由所有人毫无例外地遵守？首先，在庆祝这至圣庆节时，若追随犹太人的做法，实属不当；他们以极大的罪恶不虔地玷污了自己的双手，因此理当遭受灵魂的盲目。因为我们有能力，若放弃他们的习俗，按更真实的规定——从受难之日直到现在所保存的——将这一礼仪的适当遵守延续到后世。因此，让我们与可憎的犹太群体毫无共同之处；因为我们从救主领受了另一条道路。一条既合法又光荣的道路已为我们的至圣宗教敞开。亲爱的弟兄们，让我们同心同意采纳这条道路，并远离他们的一切卑劣行为。因为他们的夸耀实属荒谬，竟说没有他们的指教，我们就不可能遵守这些。因为他们既然犯下了弑主的杀父之罪，此后便不受理性支配，而受无法驾驭的情欲左右，并被他们内在的狂乱精神所驱使，又怎能作出正确的判断？因此，在这一问题上，如同在其他问题上一样，他们对真理毫无知觉，以致全然不知这一问题的真正调整，有时一年内庆祝两次逾越节。那么，我们为何要追随那些公开陷于严重错误的人？我们绝不会同意在同一年内第二次庆祝这庆节。但假若这些理由不够充分，你们的明智仍应努力并不断祈祷，以免你们的灵魂纯洁似乎因与这些最邪恶之人的习俗交往而受任何玷污。我们还必须考虑，在如此重要的问题上，且关于如此神圣的庆节，意见分歧是错误的。因为我们的救主留给我们一个庆节，以纪念我们得救的日子——我是指他至圣受难之日；他愿意他的公教会是一个，其成员虽分散在许多不同地方，却由同一渗透的精神，即天主的旨意所滋养。愿你们的圣德的明智反思：在同一日子里，一些人守斋，另一些人欢宴，这是多么可悲可耻；同样，在逾越节的日子之后，一些人出席宴会和娱乐，而另一些人则履行规定的斋戒。那么，天主的眷顾显然愿意（我想你们全都清楚看到），这一习俗应得到适当纠正，并归结为一条统一的规则。\par

% structure: p0041 source=h2 target=subsection reason=relative-heading-rank; base=h2

\subsection{第19章　劝勉追随世界大部分地区的榜样}

因此，既然需要纠正此事，使我们与那个弑杀主的凶手民族毫无共同之处；既然西部、南部和北部世界所有教会，以及东部一些教会所遵守的安排是合乎体统的；因此所有人在此场合一致认为值得采纳。我自己也已承担使这一决定得到你们明智的认同，并希望你们的智慧乐意接纳那同时为罗马城、非洲、全意大利、埃及、西班牙、高卢、不列颠、利比亚和全希腊，以及亚细亚、本都和奇里乞亚各教区所遵守的做法，它们完全同心合意。并且你们不仅要考虑到我所列举的地区教会数目远比其他地区多，而且考虑到所有人联合起来渴望合理的要求，并避免参与犹太人的背信行为，是最为合宜的。总之，用尽可能简短的话说，全体一致决定，至圣的逾越节应在同一天庆祝。因为一方面，在如此神圣的问题上意见分歧是不合适的；另一方面，最好采取一项免于奇怪愚蠢和错误的决定。\par

% structure: p0043 source=h2 target=subsection reason=relative-heading-rank; base=h2

\subsection{第20章　劝勉服从大公会议的法令}

\Quote{那么，请以完全甘愿的心接受这真正神圣的训令，并视其为天主的恩赐。因为凡在主教们的神圣集会中所决定的，都应被视为天主旨意的显示。因此，一旦你们将这些事项传达给我们所有亲爱的弟兄，你们就有义务从那时起为自己采纳并吩咐他人遵守上述安排和对这至圣之日的适当庆祝；以便当我来到你们的爱德面前——这是我长久渴望的——我能够与你们在同一天庆祝圣节，并在一切事上与你们一同喜乐，当我看到撒殚的残酷权势因天主的助佑，通过我们的努力而被除去，而你们的信德、平安与和谐处处兴盛。愿天主保佑你们，亲爱的弟兄们！}\par

皇帝将这封信的一份真实副本传送至每一省份，阅读者可在其中如镜中般看见他思想及其对天主虔诚的纯真诚意。\par

% structure: p0046 source=h2 target=subsection reason=relative-heading-rank; base=h2

\subsection{第21章　主教们离别时劝勉保持和谐}

当大公会议即将结束时，他召集所有主教在指定之日来见他，他们到达后，他向他们发表告别演说，在演说中他劝勉他们勤于维护和平，避免彼此之间的争辩和嫉妒，若有人显得智慧和口才出众，应将一人的卓越视为众人共有的恩赐。另一方面，他提醒他们，更有天赋的人应克制自己，不要因谦卑弟兄的不足而自高，因为判断真正优越的是天主的特权。相反，他们应体恤弱者，记着绝对完美在任何情况下都是罕见的品质。因此，每个人都应愿意对别人的小过予以宽恕，以爱德看待并忽略人性的软弱；将彼此和睦视为最高荣誉，以免他们的分歧给那些随时亵渎天主的话语的人留下笑柄：我们应尽一切力量拯救这些人，而除非我们的行为在他们看来有吸引力，否则无法做到。但你们很清楚，见证不会给所有人带来祝福，因为有些听众只满足于获得肉身的需要，另一些人则讨好上司；有些人钟情于好客待他们的人，还有些人因接受礼物而爱戴施予者；但真正渴望见证之言的人很少，真理之友更是稀有。因此有必要努力迎合所有人，如同医师一般，为每人提供有益于灵魂健康的东西，以使救恩之道被所有人完全尊重。他劝勉的前部分便是如此；最后他命令他们为他向天主献上恳切的祈祷。这样告别后，他允许他们返回各自的国家；他们满心欢喜地回去，从此他们在皇帝面前达成的一致意见继续盛行，长期分裂的人如同同一身体上的肢体般团结在一起。\par

% structure: p0048 source=h2 target=subsection reason=relative-heading-rank; base=h2

\subsection{第二十二章 他如何遣散一些人，致信其他人，以及他的赠礼}

皇帝满心欢喜于这样的成功，似以信件的形式向未出席大公会议的人献上甘美的果实。他还下令向城乡全体民众赐予丰厚的金钱礼物，以此庆祝他统治二十周年的节日。\par

% structure: p0050 source=h2 target=subsection reason=relative-heading-rank; base=h2

\subsection{第二十三章 他如何致信埃及人，劝勉他们和平}

此时，当其他人都已和平，唯独埃及人之间仍激荡着无法和解的争执，再次扰乱皇帝的安宁，但未激起他的愤怒。因为他以极大的尊敬对待争辩双方，视他们如父辈，甚至如天主的先知；他再次召他们前来，再次耐心地居中调解，以礼物尊荣他们，并将其仲裁的结果通过信件传达。他确认并批准了大公会议的法令，呼吁他们努力追求和谐，不要分裂和撕裂教会，而要铭记天主的审判。皇帝亲手书写了一封信传达这些指示。\par

% structure: p0052 source=h2 target=subsection reason=relative-heading-rank; base=h2

\subsection{第二十四章 他如何经常写信给主教和人民，内容具有宗教性质}

除此之外，他还有大量关于同类主题的著作，他写了无数信件，有些给主教，在其中他下达了有益于天主的教会的指示；有时这位三倍有福者会向全体教会会众讲话，称他们为自己的弟兄和同仆。但我们或许日后会抽时间将这些信函单独收集整理，以免插入它们损害本史书的完整性。\par

% structure: p0054 source=h2 target=subsection reason=relative-heading-rank; base=h2

\subsection{第二十五章 他如何下令在耶路撒冷、我们救主复活圣地建造一座教堂}

此后，虔诚的皇帝着手另一项真正值得记载的工作，在巴勒斯坦省。那么这是什么工作？他认为自己有义务使我们救主复活的有福之地成为众人敬仰和朝圣的对象。他因此立即下令在该处建造一所祈祷之所；他这样做，不仅是出于自己头脑的自然冲动，更是由救主自己所感动。\par

% structure: p0056 source=h2 target=subsection reason=relative-heading-rank; base=h2

\subsection{第二十六章 圣墓曾如何被不敬之人用垃圾和偶像覆盖}

因为在过去，不敬之人（或者不如说，整个邪魔族类藉着他们）曾努力将那不朽的神圣纪念碑——光辉的天使曾从天降下，为那些仍存铁石心肠、以为永生之主仍躺在死人中的人滚开石头，并向妇女们宣告喜讯，藉着他们寻求的那位是活着的这一确信除去她们硬心的不信——交与遗忘的黑暗。这不敬和不信神的人曾想将这圣穴完全从人的眼中除去，愚昧地以为这样就能有效掩盖真理。于是，他们费力从远方运来大量泥土，覆盖了整个地方；然后堆到一定高度，铺上石板，将这圣穴藏在巨大的土丘之下。然后，仿佛他们的目的已经达到，他们在这地基上预备了一个真正可怖的灵魂坟墓，为他们称为\OriginalTerm{维纳斯}{Venus}的不洁之神建造了一个阴森的死偶像神龛，并在亵渎和受诅咒的祭坛上献上可憎的祭品。因为他们以为，若不将这圣穴埋葬在这些污秽之下，他们的目的就无法完全实现。不幸的人！他们无法理解，他们的企图如何能不被那已战胜死亡的得胜者所知，正如炽热的太阳，当他升于地上，照常行于天空中央，不为全人类所见一样。事实上，他的救恩之力以更明亮的光辉照耀，不仅照亮人的身体，更照亮人的灵魂，早已以自身的光芒充满世界。尽管如此，不敬和邪恶之人的这些计谋长久以来占了上风，没有一个总督或军事指挥官，甚至皇帝本人，曾出现有能力废除这些狂妄的不敬，只有那位蒙万王之王恩宠的人除外。如今，他既在\OriginalTerm{圣神}{divine Spirit}的引导下行事，就不能任由我们所说的圣地因敌人的计谋而被埋在各种污秽之下，并被遗忘和忽视；他也不屈服于那些犯此罪责者的恶意，而是呼求天主的助佑，命令彻底净化这地方，认为被敌人污染最严重的部分应当藉着他获得天主恩宠的特殊标记。于是，他的命令一发出，这些欺骗的器械就从它们高傲的高处被推倒在地，错误的居所连同雕像及其代表的邪魔都被推翻并彻底毁灭。\par

% structure: p0058 source=h2 target=subsection reason=relative-heading-rank; base=h2

\subsection{第二十七章 君士坦丁如何命令将偶像庙宇的材料和土壤本身运到远处。}

皇帝的热诚并不止于此；他进一步命令，将所摧毁的材料，无论石料还是木料，都从原地移走并抛到尽可能远的地方；这道命令也迅速执行了。然而，皇帝并不满足于进展到此；他再次以神圣的热情，命令将地面挖至相当深度，并将被恶魔崇拜的污秽所污染的土壤运到很远的地方。\par

% structure: p0060 source=h2 target=subsection reason=relative-heading-rank; base=h2

\subsection{第二十八章 发现至圣墓穴。}

这也毫无延迟地完成了。但当覆盖泥土之下的原始地面一出现，立刻，完全出乎意料，我们救主复活的可敬而神圣的纪念碑被发现了。于是，这至圣的洞穴忠实地相似他的复活，即：在埋葬于黑暗之后，它再次重现于光明，并向所有前来观看的人提供了那地方曾经发生过奇迹的清晰可见的证据，一个比任何声音都更清晰地见证救主复活的证据。\par

% structure: p0062 source=h2 target=subsection reason=relative-heading-rank; base=h2

\subsection{第二十九章 他如何写信给各省总督和玛加略主教关于建立一座圣堂的事。}

\Emph{立即}，在我所记述的这些事之后，皇帝发出了真正虔诚精神的训令，同时拨出充足的款项，命令在救主墓旁建造一座与天主敬拜相称的祈祷之所，规模宏大而皇家气派。他确实早已有此意向，并如同藉着一种更高的智慧预见了后来将要发生的事。因此，他命令东方各省总督，通过慷慨而不吝啬的开支，确保以高贵而宏伟的规模完成这项工程。他还将下列信件寄给当时耶路撒冷教会的主教，在信中他清楚申明了信仰的救恩教义，言词如下。\par

% structure: p0064 source=h2 target=subsection reason=relative-heading-rank; base=h2

\subsection{第三十章 君士坦丁致玛加略关于建造救主圣堂的信函。}

\OriginalTerm{胜利者君士坦提乌斯·马克西穆斯·奥古斯都}{Victor Constantius, Maximus Augustus} 致 \Person{玛加略}。\par

如此，我们救主的恩宠，没有任何语言的能力足以描述我所要提及的这一奇妙情景。因为，他那最圣洁的受难纪念物，如此长久地深埋地下，竟在漫长岁月中不为人知，直至因那众人之公敌被铲除而得以自由释放的仆人们面前重新显现，这一事实确实超越了一切赞叹。纵使举世公认的智者齐聚一处，竭尽全力要为此事说些相称的话，他们也无法丝毫达到目的。确实，这奇迹的本质超越人类理性之能力，犹如天上的事超越人间之事。因此，我始终首要且唯一的目标便是：随着真理的权威日日以新的奇迹彰显自身，我们的灵魂也能以全部的清醒与热忱的同心合意，更加热切地致力于神圣法律的尊荣。因此，我特别希望你能相信——我想这对其他所有人也是显然的——即我最大的关切莫过于如何以辉煌的建筑来装饰那神圣之地；这地在神圣指引下，我已如同卸下了污秽偶像崇拜的重负；这地从起初在天主眼中便被视为圣洁，如今却因昭示了我们救主受难的明确确证而显得更加圣洁。\par

% structure: p0067 source=h2 target=subsection reason=relative-heading-rank; base=h2

\subsection{第三十一章　圣殿应在墙壁、圆柱和大理石的美丽上超越世上所有教堂}

因此，你的睿智宜为工程所需的一切作出安排与供应，不仅使整座教堂在美丽上超越其他一切，而且建筑物的细节也应如此，以致于帝国任何城市中最美丽的建筑都能被此超越。关于墙壁的建造与装饰，兹通知你：我们的朋友德拉奇利安努斯，代理\OriginalTerm{近卫军长官}{Prætorian Præfects}，以及该省总督，已从我们处领受命令。因为我们敬虔的指示是：工匠与劳工，以及他们从你的睿智处所了解的对工程进展所需的一切，都应由他们立即提供。至于圆柱和大理石，凡你实地视察计划后认为特别珍贵且有用的，请尽力书面告知我们，以便根据你的信中所认为需要的材料数量或种类，从各处按需采购，因为世上最奇妙之地理应得到相称的装饰。\par

% structure: p0069 source=h2 target=subsection reason=relative-heading-rank; base=h2

\subsection{第三十二章　他就屋顶的美化、工人及材料指示总督们}

\Quote{关于教堂的天花板，我想从你处得知，依你判断，它应该采用格子天花板，还是以其他方式完成。如果采用格子天花板，也可以饰以黄金。至于其余事项，请你的圣洁尽早告知上述长官，需要多少工人和工匠，以及需要多少经费。你也要立即向我们报告，不仅关于大理石和圆柱，还有格子天花板，如果你认为这是最美丽的样式。愿天主保佑你，亲爱的弟兄！}\par

% structure: p0071 source=h2 target=subsection reason=relative-heading-rank; base=h2

\subsection{第三十三章　如何建造了我们的救主之教堂，即圣经中所预言的新耶路撒冷}

这便是皇帝的信函；他的指示立即得到执行。因此，就在救主受难的那一地点，建造了一座新耶路撒冷，正对着那古老而著名的耶路撒冷；后者因杀害主的罪污而遭受了彻底荒凉的极端惩罚，这是天主对其不虔敬子民的神圣审判。皇帝正是对着这座城，开始以丰富而奢华的辉煌，为救主战胜死亡建立纪念碑。或许，这就是先知预言中所说的第二座新耶路撒冷，关于这城，在神圣启示的圣经中有如此丰富的见证。\par

首先，他装饰了那神圣的洞穴本身，作为整个工程的主要部分，以及那神圣的纪念碑——曾有一位光芒四射的天使在那里向众人宣告那首先在救主身上彰显的再生。\par

% structure: p0074 source=h2 target=subsection reason=relative-heading-rank; base=h2

\subsection{第三十四章　描述圣墓的结构}

因此，这纪念碑，作为整个工程的首要部分，皇帝热忱的慷慨用稀有的圆柱美化了它，并大量丰富了各种最华丽的装饰。\par

% structure: p0076 source=h2 target=subsection reason=relative-heading-rank; base=h2

\subsection{第三十五章　描述前院与柱廊}

他接下来关注的一片广阔空地，敞向纯净的天空。他以精细打磨的石板铺砌地面，并在三面环绕以极长的柱廊。\par

% structure: p0078 source=h2 target=subsection reason=relative-heading-rank; base=h2

\subsection{第三十六章　描述教堂主体的墙壁、屋顶、装饰与镀金}

因为，在与洞穴相对的一侧，即东侧，教堂本身被建立起来：一座崇高而宏伟的工程，高耸至巨大高度，长度和宽度也极为广阔。这座建筑的内部铺以各色大理石石板；而墙壁的外表面，以精确拼合的光滑石块闪耀，呈现出丝毫不逊于大理石的灿烂。至于屋顶，外部覆以铅皮，以防冬季雨水。而屋顶内部，以雕刻的格子板完成，如同一片广阔的海洋，一系列相连的格子覆盖整座教堂；并且通体覆以最纯净的黄金，使整座建筑仿佛光芒四射。\par

% structure: p0080 source=h2 target=subsection reason=relative-heading-rank; base=h2

\subsection{第三十七章　描述两侧的双柱廊与东面的三门}

此外，教堂两边各有一座柱廊，上下两层都有柱子，长度与教堂本身相当；这些柱廊的屋顶也用黄金装饰。在这些柱廊中，教堂外侧的那些由巨大的柱子支撑，而内侧的那些则立在表面装饰精美的石墩上。三个大门正对东方，用于接纳进入教堂的人群。\par

% structure: p0082 source=h2 target=subsection reason=relative-heading-rank; base=h2

\subsection{第三十八章　半圆殿、十二根柱子和他们的碗的描述}

正对这些大门，整个建筑的冠冕部分是半圆殿，它高耸至教堂的最顶端。这个半圆殿由十二根柱子环绕（按照我们救主宗徒的数目），柱头饰有巨大的银碗，这些银碗是皇帝本人作为华丽的奉献献给他的天主的。\par

% structure: p0084 source=h2 target=subsection reason=relative-heading-rank; base=h2

\subsection{第三十九章　内院、拱廊和门廊的描述}

接着，他围起了\OriginalTerm{中庭}{atrium}，它占据了教堂前面通往入口的空间。这包括：首先是庭院，然后是两侧的柱廊，最后是庭院的大门。之后，在露天市场的中央，工艺精湛的通用入口大门，让外面的路人看到内部景象，这景象不能不令人惊叹。\par

% structure: p0086 source=h2 target=subsection reason=relative-heading-rank; base=h2

\subsection{第四十章　关于他奉献的数量}

因此，皇帝建造了这座圣殿，作为救主复活的显着纪念，并以帝王般的宏伟规模将其装饰得金碧辉煌。他还用无数华丽非凡、材质各异的奉献品来丰富它——金、银、宝石，这些奉献品在大小、数量和种类上的精巧而细致的布置，我们目前没有闲暇来详细描述。\par

% structure: p0088 source=h2 target=subsection reason=relative-heading-rank; base=h2

\subsection{第四十一章　关于在\Place{白冷}和\Place{橄榄山}建造教堂}

在同一地区，他还发现了其他地方，因是两个神圣洞穴的地点而可敬；他也用奢华的宏伟来装饰它们。在一个洞穴，他给予恰当的荣耀，那是我们救主神圣临在首次彰显的地方，当时他屈尊成为血肉之身而生；而在第二个洞穴，他神圣化了从山顶升天至天堂的纪念。他如此高贵地表达了对这些地方的敬畏，同时使他的母亲永垂不朽，她曾是赐予人类如此宝贵恩惠的工具。\par

% structure: p0090 source=h2 target=subsection reason=relative-heading-rank; base=h2

\subsection{第四十二章　皇帝\Person{君士坦丁}的母亲\Person{赫莲娜}皇后出于敬拜目的访问该地并建造了这些教堂}

因为她已决心向天主、万王之王履行虔诚敬拜的义务，并感到自己有责任为她自己的儿子（如今如此伟大的皇帝）和他的儿子们（她自己的孙辈，蒙天主眷顾的\OriginalTerm{凯撒们}{Cæsars}）献上感恩祈祷，尽管年事已高，却拥有非凡的智慧，以年轻人的活力匆忙前去视察这片可敬的土地；同时以真正的帝王般的关怀，访问东方各省、城市和人民。因此，一旦她向救主足迹所踏之地给予了应有的敬畏，按照先知的话说：\Quote{让我们在他脚所站的地方敬拜}，她立即将她的虔诚果实遗赠给后代。\par

% structure: p0092 source=h2 target=subsection reason=relative-heading-rank; base=h2

\subsection{第四十三章　关于\Place{白冷}教堂的进一步说明}

因为她毫不迟延地建造了两座教堂奉献给她所敬拜的天主，一座在救主诞生的洞穴处，另一座在他升天的山上。因为那被称为\Quote{天主与我们同在}的一位，甚至屈尊出生在地上的一个洞穴中，他的诞生地希伯来人称为\Place{白冷}。因此，虔诚的皇后用罕见的纪念物尊崇了那位生下天上圣童的妇人生产的场所，并用一切可能的辉煌美化了圣洞。皇帝本人不久后通过君王的奉献表达了对该地的敬意，并用银、金和刺绣挂毯等贵重礼物增加了母亲的宏伟。此外，皇帝的母亲还在\Place{橄榄山}上建造了一座宏伟的建筑，纪念人类的救主升天，在山顶建造了一座圣教堂和圣殿。的确，可靠的历史告诉我们，救主正是将这个秘密启示传授给他的门徒的。在这里，皇帝也通过多种昂贵的奉献表达了对万王之王的敬畏。因此，虔诚的\Person{赫莲娜·奥古斯塔}，一位虔诚皇帝的母亲，在这两个神秘的洞穴上建造了这两座高贵而美丽的奉献纪念物，永远值得纪念，以荣耀她的救主天主，并作为她神圣热忱的证明，从她儿子那里获得了其皇权的帮助。不久之后，这位年迈的妇人就收获了她劳苦的应得奖赏。在她一生中，直至晚年，都处于极大的繁盛中，并在言语和行为上结出服从神圣诫命的丰硕果实，因而享有轻松安宁的生活，身心精力未衰，最终她从天主那里得到了与她虔诚生活相称的结局，甚至在此生就得到了她善行的报答。\par

% structure: p0094 source=h2 target=subsection reason=relative-heading-rank; base=h2

\subsection{第四十四章　关于\Person{赫莲娜}的慷慨和慈善行为}

因为在她以帝王权威的辉煌巡视东方各省期间，她向各城市的居民集体以及接近她的个人，都充分证明了她的慷慨，同时以慷慨之手向士兵分发赏赐。但她给予赤身露体和无人保护的穷人的礼物尤其丰富。她给一些人金钱，给另一些人充足的衣物；她释放了一些人出狱，或脱离矿场的痛苦奴役；她将另一些人从不公正的压迫中解救出来，还有一些人，她从流放中召回。\par

% structure: p0096 source=h2 target=subsection reason=relative-heading-rank; base=h2

\subsection{第四十五章　\Person{赫莲娜}在教堂中的虔诚行为}

然而，尽管她的品格因我所描述的这些事迹而熠熠生辉，她却远未忽略对天主的个人虔诚。人们时常看见她出入于祂的教会，同时以丰盛的奉献装饰祈祷之所，连最小城市的教堂也不忽略。简言之，这位令人钦佩的妇人穿着朴素谦逊的服装，与崇拜的人群相交，并以一贯虔诚的行为来证明她对天主的敬爱。\par

% structure: p0098 source=h2 target=subsection reason=relative-heading-rank; base=h2

\subsection{第四十六章 她如何立遗嘱，并在八十岁时去世。}

当她在漫长的一生结束时，终于被召去承受更幸福的命运，时年八十岁，且临近去世之时，她拟定并执行了最后遗嘱，将财产遗赠给她唯一的儿子，即皇帝，世界唯一的君主，以及她的孙辈，即他的儿子们，即凯撒们，她分别将她在世界各地拥有的所有财产遗赠给他们。立完遗嘱后，这位三倍有福的妇人在她杰出的儿子面前去世，当时他随侍在侧，照料她，握着她的手：因此，对于正确辨识真理的人而言，这位三倍有福者似乎并未死亡，而是经历了一种真正的改变和过渡，从尘世生命转入天上生命，因为她的灵魂仿佛被重塑为不朽的、天使般的本质，被接升到她的救主面前。\par

% structure: p0100 source=h2 target=subsection reason=relative-heading-rank; base=h2

\subsection{第四十七章 君士坦丁如何埋葬他的母亲，以及他生前如何尊敬她。}

她的遗体也受到了特别的尊敬，由一队庞大的卫兵护送着前往皇城，并安放在一座皇陵中。这就是我们皇帝母亲最后的时日，她是一位值得永远纪念的人，不仅因为她自身的实际虔诚，也因为她生育了如此非凡和令人钦佩的后代。他的品格完全可以称为有福，既因他的孝爱，也因其他原因。他通过自己的影响使她成为如此虔诚的天主崇拜者（尽管她此前并非如此），以至于她似乎从一开始就由人类的救主所教导；除此之外，他以皇家的尊荣充分尊敬她，以致在每个省份，甚至在军队行列中，人们都以\OriginalTerm{奥古斯塔}{Augusta}和\OriginalTerm{皇后}{empress}的称号称呼她，她的肖像也被铸在金币上。他甚至授予她使用和支配皇家财宝的权力，在任何情况下都按她自己的意愿和判断使用：因为她在儿子手中也得到了这一令人羡慕的特权。因此，在他记忆中增光添彩的诸多品质中，我们完全可以包括那种超越一切的孝爱，借此他完全顺从了神圣的诫命，即要求子女对父母给予应有的尊敬。如此，皇帝在巴勒斯坦执行了我以上描述的高贵工程：实际上，在每个省份，他都以比以往时代更为宏伟的规模兴建了新的教堂。\par

% structure: p0102 source=h2 target=subsection reason=relative-heading-rank; base=h2

\subsection{第四十八章 他如何兴建教堂以纪念殉道者，并在君士坦丁堡废除偶像崇拜。}

由于他决意要以其名字命名的城市以特殊荣誉，他用众多神圣建筑装饰了它，既有规模宏大的殉道者纪念堂，也有其他最为辉煌的建筑，不仅在城市内，也在其近郊：因此，他同时向殉道者的记忆致敬，并将他的城市奉献给殉道者的天主。他充满神圣智慧，决心要使其名字所命名的城市清除一切偶像崇拜，从此不得在那被妄称为神的庙宇中崇拜任何雕像，也没有任何祭坛被血的污染玷污：不得有火焚的祭献，不得有邪魔的节庆，也不得有通常由迷信者遵守的任何其他仪式。\par

% structure: p0104 source=h2 target=subsection reason=relative-heading-rank; base=h2

\subsection{第四十九章 十字架在宫殿中的图像，以及达尼尔在公共喷泉的图像。}

另一方面，人们可以看到市场中央的喷泉装饰着代表\OriginalTerm{善牧}{good Shepherd}的雕像，这为研读圣言的人所熟知，还有\OriginalTerm{达尼尔}{Daniel}与狮子的雕像，用黄铜铸造，并饰以金箔闪闪发光。确实，如此大量的神圣之爱充满皇帝的灵魂，以至于在皇宫正殿本身，在一块巨大的嵌板上，展示在镀金镶板天花板的中央，他让人固定了救主受难的象征，由各种宝石与黄金细工镶嵌而成。这个象征似乎被他意图作为帝国本身的守护。\par

% structure: p0106 source=h2 target=subsection reason=relative-heading-rank; base=h2

\subsection{第五十章 他在尼科美底亚及其他城市兴建教堂。}

装饰了以他名字命名的城市之后，他接下来通过建造一座庄严宏伟的教堂来装饰比提尼亚的首府，渴望在这座城市也兴建一座纪念他战胜自己的敌人和天主的敌人的纪念物，以尊崇他的救主，并由他自己出资。他还用美丽的神圣建筑装饰了其他省份的主要城市；例如，在东方的大都市中，那座因安提约古而得名的城市，作为帝国那部分的首府，他将一座规模无与伦比、美丽非凡的教堂奉献给天主的服务。整个建筑被一个广阔的围墙环绕，教堂本身高耸入云，呈八角形，四周环绕着许多厅室、庭院以及上下层房间；整体以大量的金、铜和其他最昂贵的材料华美装饰。\par

% structure: p0108 source=h2 target=subsection reason=relative-heading-rank; base=h2

\subsection{第五十一章 他下令在玛默勒兴建一座教堂。}

如此便是奉皇帝命令建造的主要圣殿。但他听说那同一位救主，昔日曾显现于世上，在远古的时代里，就已向巴勒斯坦的圣人们，在\Place{玛默勒}橡树附近，显示过祂神圣的临在，就下令在那里也建造一座祈祷之所，以敬礼那位如此显现的天主。于是，皇帝的敕令通过分别致函各地方总督而传达，命令迅速完成指定的工程。他还给本史的笔者寄来了一封雄辩的训谕，我认为将其收入本书是合适的，以便使人对其虔诚的勤勉与热忱有一个正确的认识。那么，为了表达他对所听闻的在那地点惯行的恶行的不悦，他用以下的话语向我致函。\par

% structure: p0110 source=h2 target=subsection reason=relative-heading-rank; base=h2

\subsection{君士坦丁关于\Place{玛默勒}致尤西比乌斯的信函}

\Signature{胜利者君士坦丁，最伟大的奥古斯都，致玛加略及巴勒斯坦其余诸位主教。}\par

一件并非寻常重要的恩惠，由我那位真正虔诚的岳母施予我们，因为她通过信函向我们揭露了那些不敬神的人的愚妄行径，这些行径迄今未被你们察觉；因此，这被忽视的罪行现在可以通过我们获得适当的纠正与补救，尽管为时已晚，但却是必要的。因为的确，圣洁之地被不洁的污点玷污，实在是一种严重的亵渎。那么，最亲爱的弟兄们，这究竟是什么事，虽然逃过了你们的明察，却由我所提及的那位因着虔诚的责任感而迫于披露呢？\par

% structure: p0113 source=h2 target=subsection reason=relative-heading-rank; base=h2

\subsection{救主在此地向\Person{亚巴郎}显现}

\Quote{她向我保证，那以\Place{玛默勒}橡树命名的地点（我们发现有\Person{亚巴郎}曾住在此处），正被某些迷信的奴隶用各种方式玷污。她宣称，在那棵树的原址上竖起了本应彻底摧毁的偶像；附近设有一座祭台，不停地举行不洁的献祭。如今，既然这些行为显然既与我们的时代特征不符，也与该地点本身的圣洁不相称，我愿告知各位尊驾，我们尊贵的朋友\Person{阿卡丘斯}伯爵已收到我的书函指示：凡在上述地点发现的偶像，应立即投入火焰；祭台应彻底拆除；若有人在此命令之后，仍在该地点犯下任何不敬神的行为，他将受到应得的惩罚。该地点本身，我们已下令用一座洁净的建筑，即一座教堂，来装饰；以便它成为圣人们合适的集会场所。同时，若有任何违反这些命令的行为发生，你们应毫不迟延地通过信件向我们禀报，以便我们指示对违法者予以最严厉的惩处。因为你们并非不知道，至高的天主最初是在那个地点向\Person{亚巴郎}显现并与他交谈的。在那里，神圣法律的遵守首次开始；在那里，救主亲自协同两位天使，首次恩赐\Person{亚巴郎}祂临在的显现；在那里，天主首次向人显现；在那里，祂向\Person{亚巴郎}应许了未来的后裔，并立即实现了那应许；在那里，祂预言他将成为多民族之父。基于这些原因，在我看来，这个地点不仅应当通过你们的勤勉而保持纯净、免受一切玷污，而且应当恢复其原有的圣洁；使今后在那里除了向全能的天主、我们的救主和万有的主举行合宜的敬礼之外，别无他事。而这一敬礼，你们应当以应有的关注予以照料，如果各位尊驾愿意（我对此确信）满足我的愿望，而我尤其关心对天主的崇拜。愿祂保佑你们，亲爱的弟兄们！}\par

% structure: p0115 source=h2 target=subsection reason=relative-heading-rank; base=h2

\subsection{各处偶像庙宇与雕像的毁灭}

上述诸事，皇帝尽心竭力，以赞颂基督拯救之权能，并以此为他恒常的目标，即荣耀他的救主天主。另一方面，他千方百计斥责异教徒的迷信谬误。因此，各城神庙的入口，因他下令拆去门扇，而暴露于风雨之中；有的庙宇被揭去瓦片，屋顶被毁。还有一些庙宇中，那古老迷信曾夸耀多年的一系列可敬的铜像，如今都展现在皇城的所有公共场所：此处是\OriginalTerm{皮提亚（的阿波罗）}{Pythian}，彼处是\OriginalTerm{斯明提安的阿波罗}{Sminthian Apollo}，均令观者嗤之以鼻；而\OriginalTerm{德尔斐三脚祭台}{Delphic tripods}被放置在赛马场中，\OriginalTerm{赫利孔山的缪斯}{Muses of Helicon}则被安置在宫殿本身。简言之，以他命名的这座城市，处处充满了各地供奉的精美铜像，这些铜像原是受迷信欺骗的受害者曾用无数牺牲和燔祭徒然尊崇为神明的，而今终于因皇帝将这些崇拜对象展示为众目睽睽之下的笑柄和嘲弄而认识其错误。至于那些金像，他则另作处理。因为他一了解到无知的民众对这些金银制成的错误恐吓物怀有虚妄而幼稚的恐惧，便认为应当将这些东西也除去，如同在黑暗中行走的人们面前的绊脚石，从此为众人开辟一条平坦无阻的御道。他既作出这一决定，便觉得无须任何士兵或军队来铲除这邪恶：只需他几位友人便可完成此事，他便一言传旨，打发他们巡视各省。他们仗着皇帝虔诚的意向和自己对天主的虔诚，穿过无数部落和民族，在各城各乡废除这一古代谬误。他们在众人嘲笑与轻蔑中，命令司祭本人将他们的神祇从黑暗的隐处带到光天化日之下；然后剥去其装饰，将那隐藏在彩绘外表之下的丑陋实情展示给众人观看。最后，他们将材料中任何有价值的部分刮下，投入火中熔化以检验其价值，然后将其收藏并保留他们认为必要之物，而将完全无用的部分留给迷信崇拜者，作为他们羞耻的纪念。与此同时，我们可钦羡的皇帝本人也从事于类似的工作。正如我们所说，在这些昂贵的死者雕像被剥去珍贵材料的同时，他也攻击了那些铜像：他下令用绳索将这些神话妄认为神灵的雕像从原处拖出，仿佛将它们掳走一般。\par

% structure: p0117 source=h2 target=subsection reason=relative-heading-rank; base=h2

\subsection{第五十五章 推翻腓尼基阿法卡的偶像神庙并废除其中淫秽习俗}

皇帝其次关心的是，仿佛点燃一支明亮的火炬，借其光芒以帝王的目光环视四周，看看是否仍有隐藏的谬误痕迹留存。正如眼光敏锐的鹰隼在其向天的飞翔中，能从其高远之处看清地上最遥远之物，同样，当他驻跸于自己美丽城市的皇宫之中时，便如同从瞭望台上，发现腓尼基省有一个隐藏而致命的灵魂罗网。这是一座树林和庙宇，不位于任何城市之中，也不在任何公共场所（通常为追求辉煌效果而如此），而是远离寻常大道，在\Place{阿法卡}的\Place{黎巴嫩山}一处山顶上，供奉着名为\Person{维纳斯}的污秽魔鬼。这是一所不洁者的邪恶学校，那些以柔弱毁坏自身肉体的人在此献身。此处，忘其性尊严的男子以柔弱行为讨好魔鬼；此处，妇女的非法交易和奸淫以及其他可怕的可耻行径，在这座庙宇中如同在法纪之外的地点进行。同时，这些恶行因无观察者在场而未被制止，因为品行端正的人无人敢涉足此地。然而，这些行径未能逃过我们威严皇帝的警觉；他以他特有的先见之明亲自察看了这些情况，断定这样的庙宇不配享受天光，于是下令将其连同供物一并彻底摧毁。因此，遵奉皇帝旨意，这些不洁迷信的器具当即被废除，并以武力之手洁净了这地方。如今，那些先前生活毫无约束的人，因皇帝以惩罚相威胁而学会了自制；同样，那些自以为是的外邦迷信者，如今也亲身体验了他们的愚昧。\par

% structure: p0119 source=h2 target=subsection reason=relative-heading-rank; base=h2

\subsection{第五十六章 摧毁埃加城的埃斯科拉庇俄斯神庙}

因为那些自命不凡者的一种广泛传播的错误关乎在\Place{基里基雅}受人崇拜的邪魔，成千上万的人视其为拥有救恩与医治能力的拥有者而敬畏他，他有时向那些在殿中守夜的人显现，有时使患病者恢复健康——但他其实是灵魂的毁灭者，引诱其轻易受骗的信徒脱离真正的救主，使他们陷入不虔的错误。皇帝出于其一贯作风和推进崇拜那既是忌邪的天主又是真正救主者的愿望，下令这座殿也应被夷为平地。一支兵士队伍立即服从此命令，将这座令高尚哲人赞叹的建筑连同其无形之居客一同化为尘土。这居客既非邪魔也非神，毋宁是灵魂的欺骗者，他通过不同世代欺骗人类如此之久。因此，那位曾许诺拯救他人免于厄运与苦难者，竟无法保全自身，如同神话中所传，他被闪电击中而焚毁。然而，我们皇帝的虔敬行为毫无神话或虚饰；而是凭借其救主彰显的大能，这座殿与其他殿一同被彻底推翻，连先前愚妄的痕迹也没有留下。\par

% structure: p0121 source=h2 target=subsection reason=relative-heading-rank; base=h2

\subsection{第57章 外邦人如何离弃偶像崇拜，转向认识天主}

因此，那些曾是迷信奴隶的人，当他们亲眼看见自己的迷惑被揭露，并看到各地庙宇和偶像的实际毁坏时，有些人便归向了基督的救恩之道；另一些人虽未采取这一步，却谴责他们从祖先继承的愚行，并嘲笑他们长久以来当作神的东西。事实上，当他们在自己崇拜之物的华美外表下发现深藏的彻底污秽时，心中还能有其他感受吗？在这外表之下，要么是死人的骨头或干枯的头骨，被魔术师的技艺欺诈地装饰；要么是充满可憎污秽的肮脏破布；要么是一捆干草或秕糠。当他们看见这一切堆集在无生命的偶像中时，便谴责他们父亲和他们自己的极端愚妄，尤其当他们在庙宇的隐秘处所和雕像本身之中发现没有任何居客时：既无邪魔，也无神谕宣示者，既不神，也不先知，正如他们先前所设想的那样；甚至看不到一丝暗淡模糊的幻影。因此，每一个阴暗的洞穴，每一个隐蔽的处所，都向皇帝的使者敞开：庙中不可进入的秘密房间、最内层的圣所，都被士兵的脚践踏；由此，外邦世界如此多世纪以来的心灵盲目，在众人眼前变得清晰可见。\par

% structure: p0123 source=h2 target=subsection reason=relative-heading-rank; base=h2

\subsection{第58章 他如何摧毁\Place{赫略波利}的维纳斯神庙，并在该城建造第一座圣堂}

我所描述的这种行为可被视作皇帝最高贵的成就之一，还有他对每个省份所作的明智安排。例如\Place{腓尼基}城\Place{赫略波利}，在那里，那些以特殊尊称赞美放荡欢乐的人，纵容他们的妻女无耻行淫。但现在，一项充满节制精神的新法规出自皇帝，它断然禁止先前恶习的延续。除此之外，他还送给他们书面的劝勉，仿佛他由天主特别任命为此目的，以便教导所有人贞洁的原则。因此，他并不鄙弃与这些人通信，敦促他们殷勤寻求关于天主的知识。同时，他用相应的行动跟随他的言语，甚至在这座城里建造了一座规模宏大的圣堂；因而，在以往任何时代都未曾听闻的事件，如今首次发生：一座迄今完全沉迷于迷信的城市，如今获得了一座天主的圣堂，并配有司铎和执事，其民众受一位祝圣服务于至高天主的主教照顾。此外，皇帝渴望在这里也有尽可能多的人被赢到真理之中，他为穷人的需要提供了丰厚的供应，甚至想以此邀请他们寻求救恩之道，仿佛他几乎采用那曾说\BibleQuote{无论出于假意，或是出于诚心，基督总被宣讲}之人的话。\par

% structure: p0125 source=h2 target=subsection reason=relative-heading-rank; base=h2

\subsection{第59章 关于\Place{安提约基雅}由\Person{尤斯塔提乌斯}引起的骚动}

然而，在这些事件所带来的普遍喜乐中，正当天主的教会到处并借着一切方式在帝国各处欣欣向荣时，那总是窥伺善人毁灭的嫉妒之灵再次准备对抗我们繁荣的伟大，或许期望皇帝本人会被我们的骚乱与失序所激怒，最终远离我们。因此，他在\Place{安提约基雅}挑起了一场激烈的争论，从而使那里的教会陷入一系列悲惨的灾祸，几乎导致该城完全覆没。教会的成员分裂为两个对立的派别；而民众，甚至连官员和士兵都被激怒到如此程度，以致若非天主的眷顾以及皇帝震怒的恐惧控制了群众的狂怒，这场争斗就会以刀剑解决。在这件事上，皇帝也扮演了灵魂的守护者和医者的角色，以极大的忍耐对需要者施以劝说的疗药。他仿佛以使者般温和地恳求他的子民，派遣一位最受认可、最忠诚的、享有伯爵尊荣的人士到他们中间；同时，他通过多次来信劝勉他们持守和平的精神，并教导他们实践真正的虔敬。在这些规劝取得成效后，他在随后的信件中宽恕了他们的行为，声称他已从引发骚乱的那一方亲自听取了事情的是非曲直。这些信件充满了非凡的学识和教诲，我本应在当前这部作品中收录它们，但恐怕它们会给被指控的人品性留下污点。因此，我将省略这些信件，不愿唤起对过去苦楚的记忆，只将那些他为见证其余人之间和平与和谐的重建而写的信件附在我的叙述中。在这些信件中，他告诫他们不要有任何欲望将另一位通过其干预而恢复和平的教区首领据为己有，并劝勉他们按照教会的惯例，选举万有的共同救主所指示为合适职务的人作为他们的主教。他的信分别致予民众和主教们，措辞如下。\par

% structure: p0127 source=h2 target=subsection reason=relative-heading-rank; base=h2

\subsection{第六十章 \Person{君士坦丁}致\Place{安提约基雅}人的信，指示他们不要将\Person{欧瑟伯}从\Place{凯撒勒雅}调走，而要另寻他人。}

胜利者\Person{君士坦丁}，最大奥古斯都，致\Place{安提约基雅}民众。\par

\Quote{对于明智而有智慧的人类而言，你们之间的和谐是何等令人愉快！而我本人，弟兄们，出于宗教之情，也出于你们自己的生活方式和对我所表现的熱忱，我倾向于以持久的爱心爱你们。正是通过运用正确的理解和明辨，我们才能真正享受我们的福分。有什么能比这明辨更适合你们呢？因此，我强调你们对真理的维持与其说是招致他人的仇恨，不如说是促进了你们的安全，这并不奇怪。确实，在弟兄们中间，那同样行走在真理和正义道路上的意向，通过天主的恩宠，应许他们被登记在他纯洁神圣的家庭中，有什么比欣然同意众人的兴旺更光荣的呢？尤其是因为神圣律法的规诫为你们所提议的意向提供了更好的方向，而我们自己也希望你们的判断得到适当认可的确认。也许你们感到惊讶，并且不明白我这番开场白的意思。我将毫不保留地解释其原因。我承认，在阅读你们的记录时，我通过它们对\Place{凯撒勒雅}的主教\Person{欧瑟伯}的高度赞扬见证，看出你们对他有着强烈的依恋，并希望将他据为己有。我本人也因他的学识和节制而长期了解并尊重他。那么，你们认为我对这件事持有什么想法呢？我渴望寻求并按照严格的原则行事。你们这个愿望给我带来了怎样的焦虑呢？哦，圣洁的信德，你在救主的言语和规诫中为我们的人生提供了一个典范，好像是什么，你自己多么难以抗拒人的罪恶，如果不是你拒绝为利益服务的话！依我之见，那以维持和平为首要目标的人，似乎胜过胜利本身；而当一个正确和光荣的途径可供选择时，当然没有人会犹豫采纳它。那么我问，弟兄们，为什么我们这样决定以至于因我们的选择而加害于他人？为什么我们贪求那些会毁坏我们自身声誉的对象？我本人高度尊重你们认为值得尊重和爱戴的那个人；尽管如此，那些应具有权威并约束所有人的原则完全被忽视，以至于每个人都不满足于自己的处境，所有人都享受适当的权利，这是不正确的；在考虑竞争对手的要求时，也不能认为只有一个人，而是可能有许多人显得配得上与这个人相比。因为只要教会的尊严不受暴力和粗暴的干扰，它们就处于平等地位，到处都值得同样的尊重。而且，对这一个合格性的调查损害他人是不合理的；因为所有教会，无论本身被认为大小，都同样有能力接受和维护神圣的规条，以至于如果一个教会在实践标准上并不比另一个教会差，只要我们敢于大胆宣称真理。如果是这样，我们必须说你们将负有责任，不是保留这位主教，而是错误地移走他；你们的行为将更以暴力而非正义为特征；无论别人怎么想，我敢于清楚大胆地断言，这一措施将提供对你们指控的理由，并引发最具破坏性的派系骚乱：因为即使是胆怯的羊群，当牧人的警惕关怀减弱，它们发现自己失去惯常的引导时，也能显示其牙齿的用途和力量。如果真是这样，如果我的判断没有错，那么弟兄们，请首先考虑这一点——因为许多重要考量将立即呈现——如果你们坚持你们的意图，你们之间的相互善意和感情是否会减少？其次，请记住，那位来到你们中间提供无私忠告的人，现在在天主的审判中得到了应得的赏报；因为你们对他公正行为的崇高见证，他得到了不寻常的回报。最后，按照你们一贯的明智判断，你们要表现出应有的勤勉，选择你们所需的人，谨慎避免所有派系和喧嚣的吵闹；因为这样的吵闹总是错误的，从不相谐的因素的碰撞中，火花和火焰都会升起。我声明，正如我希望取悦天主和你们，并享受与你们善意相称的幸福那样，我爱你们，爱你们温柔的宁静港湾，现在你们已经抛弃了污秽，代之以健全的道德和和谐，将神圣的标准牢牢地植根于船上，并可以说，在向着天堂之光前进的航程中，由铁质的舵引导。那么，请上船吧，那些不朽的货物，因为船上的所有污水已经排干；从此小心确保享受你们所有当前的祝福，以免将来任何时刻你们显得是在冲动或方向错误的熱忱下决定任何措施，或者最初鲁莽地走上了一条不明智的道路。愿天主保护你们，亲爱的弟兄们！}\par

% structure: p0130 source=h2 target=subsection reason=relative-heading-rank; base=h2

\subsection{第六十一章 皇帝致信欧瑟伯，赞扬他拒绝安提约基雅主教职}

\Emph{皇帝就我拒绝安提约基雅主教职致我的信}\par

\Person{胜利者君士坦丁，最大奥古斯都}致\Person{欧瑟伯}\par

\Quote{我已极其仔细地细读了你的信函，并察觉你已严格遵守了教会纪律所规定的准则。坚守既悦乐天主又与传承相符之事，乃是真虔敬的明证。你在此事上有理由自视为有福，因为你在全世界（我可说是这样判断）的评判中被认为堪当照管任何教会。因为所有人都渴望将你据为己有的心愿，无疑提升了你在这一方面的幸运。尽管如此，你那决心遵守天主的诫命和教会宗徒规条的贤明，却在谢绝安提约基雅教会的主教职、并愿意继续留在那间你最初按天主意愿接受监督的教会一事上，做得极其出色。我已就此事写信给安提约基雅的百姓，也写给你在牧职中的同僚们——他们曾为此问题咨询过我；你的圣德在阅读这些信函后，将轻易看出，由于正义本身反对他们的主张，我是按神圣指引写信给他们的。有必要请你的贤明出席他们的会议，以便这项决定能在安提约基雅教会中得到批准。愿天主保佑你，亲爱的弟兄！}\par

% structure: p0134 source=h2 target=subsection reason=relative-heading-rank; base=h2

\subsection{第六十二章 君士坦丁致大公会议的信，反对将\Person{优西比乌}从\Place{凯撒勒雅}调离。}

\Signature{胜利者君士坦丁，至尊奥古斯都，致\Person{狄奥多图斯}、\Person{狄奥多鲁斯}、\Person{纳尔奇苏斯}、\Person{阿埃提乌斯}、\Person{阿尔斐乌斯}，以及其余在\Place{安提约基雅}的主教们。}\par

\Quote{我已细读了你们各位贤明的信函，并极为赞同你们在牧职中的同僚\Person{优西比乌}的明智决定。此外，在得知部分来自你们的信函、部分来自我们杰出的伯爵\Person{阿卡西乌斯}和\Person{斯特拉提乌斯}的信中所阐明的情况后，经充分调查，我已写信给\Place{安提约基雅}的百姓，暗示一条既悦乐天主又有利于教会的途径。我已下令将此信的副本附于本函之后，以使你们自己知道，我作为正义事业的代言人，认为适合写给那百姓的内容：因为我在你们的信函中发现这样的提议，即，根据百姓的选择并经你们自己的意愿认可，\Place{凯撒勒雅}的圣主教\Person{优西比乌}应当主持并照管\Place{安提约基雅}的教会。其实\Person{优西比乌}本人就此事的信函似乎与教会规定的秩序严格一致。然而，使你们的贤明也得知我的意见也是适宜的。因为我获悉，\Place{卡帕多细亚}\Place{凯撒勒雅}的公民、司铎\Person{欧弗洛尼乌斯}，以及同样身为司铎、由\Person{亚历山大}在\Place{亚历山大里亚}任命担任此职的\Place{阿雷图萨}的\Person{乔治}，是经过考验的信德之人。因此，理应向你们的贤明表明，在提出这些人以及你们认为堪当主教尊严的任何其他人选时，你们应以符合宗徒传统的方式决定此问题。因为那样，你们的贤明将能按照教会规条和宗徒传统，以真正的教会纪律所规定的方式指导这次选举。愿天主保佑你们，亲爱的弟兄们！}\par

% structure: p0137 source=h2 target=subsection reason=relative-heading-rank; base=h2

\subsection{第六十三章 皇帝如何展示他根除异端的热情。}

皇帝向教会领袖们发出这样的劝勉，要他们凡事都为了荣耀神圣宗教而做。他以此方式平息了纷争，使天主的教会达至一致和谐。接着，他转向另一项职责，认为有责任根除另一种不敬之人，他们是人类的有害仇敌。这些人是社会的害虫，在宗教礼仪的虚伪外衣下毁灭整个城市；我们的救主在某处以警告之声称他们为假先知和贪婪的豺狼：\BibleQuote{你们要提防假先知！他们来到你们跟前，外披羊毛，内里却是凶残的豺狼。你们可凭他们的果实辨别他们。}于是，通过一道下达给各省总督的命令，他有效地驱逐了所有这类犯法者。除这道法令外，他还亲自向他们发出一封严厉唤醒的训诫，劝告他们认真悔改，以便他们仍能在天主的真教会中找到安全的港湾。那么，请听听他在此信中以何种方式对他们说话。\par

% structure: p0139 source=h2 target=subsection reason=relative-heading-rank; base=h2

\subsection{第六十四章 君士坦丁反对异端的敕令。}

\Signature{胜利者君士坦丁，至尊奥古斯都，致异端者。}\par

如今，藉此诏令，你们 \OriginalTerm{诺瓦替安人}{Novatians}、\OriginalTerm{瓦伦廷派}{Valentinians}、\OriginalTerm{马西昂派}{Marcionites}、\OriginalTerm{保禄派}{Paulians}，以及被称为 \OriginalTerm{弗吕家派}{Cataphrygians} 的人，以及所有以你们私密集会发明并支持异端的人，请明白：你们的教义与何种谎言的织体和虚妄，与何种毁灭性和有毒的谬误不可分割地交织在一起；以致通过你们，健康的灵魂罹患疾病，活人成为永死的猎物。你们这些真理与生命的仇恨者和敌人，与毁灭结盟！你们的一切计谋都与真理对立，却熟悉卑劣的行径；充满荒谬和虚构：藉此你们制造谎言，压迫无辜者，向那相信的人扣留光明。永远在虔诚的面具下越界，你们以污秽充满万物：你们以致命的创伤刺透纯洁无邪的良心，同时——可以这样说——从人眼前抽走白日的光明。但为何我要详述？因为要恰如其分地谈论你们的罪行，所需的时间和闲暇超过我所拥有的。你们的罪过目录如此漫长无尽，如此可憎且全然可怖，以至于一日不足以全部数算。而且，确实，最好耳目避开这样的主题，以免描述每一项邪恶之后，自己信德的纯真和新鲜受损。那么，为何我仍容忍如此泛滥的邪恶？尤其是这长期的宽容导致一些原本健全的人染上这瘟疫般的疾病？为何不立即，仿佛以公开表达不满的方式，斩断如此巨大祸患的根部？\par

% structure: p0142 source=h2 target=subsection reason=relative-heading-rank; base=h2

\subsection{第六十五章 异端者被剥夺聚会场所}

\Quote{\Emph{既然}，那么已无法再容忍你们的毒害性的谬误，我们藉此诏令警告：你们中任何人今后不得擅自聚集。我们已相应地命令，夺取你们所有习惯举行集会的房屋；我们的关心扩展至此，甚至禁止你们举行迷信和愚昧的聚会，不仅公开场所，而且在任何私宅或地点皆不可。因此，你们中那些渴望接纳真实和纯洁宗教的人，请采取更佳途径，进入公教会，与之在圣洁的共融中结合，如此你们将能够达到真理的认识。无论如何，你们扭曲理解的妄想必须完全停止，不再混杂并破坏我们现今时代的幸福：我是指异端者和分裂者不敬且可怜的二元心志。因为，借着天主恩宠所享有的繁荣，努力带领那些过去生活在未来福佑希望中的人，从一切不规则和错误回归正途，从黑暗到光明，从虚妄到真理，从死亡到救恩，是相称的。为了此药方得以有效施行，我们已命令，如前所述，你们被绝对剥夺每一个举行你们迷信聚会的聚集点，我指所有属于异端的祈祷屋——如果这些配得上这名称——并且这些房屋应被立即移交给公教会；任何其他地点没收归公服务，不留任何未来聚会的便利；以便从今以后，你们任何非法集会不得在任何公共场所或私宅出现。让此诏令公之于众。}\par

% structure: p0144 source=h2 target=subsection reason=relative-heading-rank; base=h2

\subsection{第六十六章 异端者中被发现禁书后，多人回归天主教会}

这样，异端者的藏身之处因皇帝的谕令被捣毁，他们所藏匿的凶恶野兽——我指他们不敬教义的主要作者——被驱散。在他们所欺骗的人中，有些被皇帝的威胁所恐吓，隐瞒真实情感，偷偷潜入教会。因为法律指示搜查他们的书籍，那些行邪恶和禁术的人被查出，这些人准备用各种伪装来确保自身安全。然而，也有另一些人，自愿且真正真诚地拥抱了更美好的希望。与此同时，各教会的主教们继续进行严格的审查，彻底拒绝那些以虚假伪装的体面外表试图进入的人，而真诚前来的人则被考验一段时间，经过充分试炼后列入会众。对待那些被指控为公开异端的人就是如此；而那些没有持守不敬教义、只因分裂派顾问的影响而与一个身体分离的人，则毫无困难或拖延地被接纳。因此，众多人仿佛从异乡归来，重访自己的祖国，承认教会为母亲，他们已长久流浪，如今欢欣喜乐地回归。于是，整个身体的成员得以联合，凝聚为一个和谐的整体；唯一的公教会与自身合一，发出完全的光辉，而异端或分裂团体无处存留。完成这伟大工作的荣誉，唯独我们这位受上天保护的皇帝，在他之前的所有人中，能够归于自己。\par

\SourceRecord{Translated by Ernest Cushing Richardson.；Nicene and Post-Nicene Fathers, Second Series；Vol. 1.；Edited by Philip Schaff and Henry Wace.；Buffalo, NY: Christian Literature Publishing Co.,；1890.；<http://www.newadvent.org/fathers/25023.htm>.}

% structure: p0001 source=h1 target=section reason=page-title

\section{君士坦丁传（卷四）}

% structure: p0002 source=h2 target=subsection reason=relative-heading-rank; base=h2

\subsection{第一章　他如何以赏赐和擢升尊荣多人}

皇帝一面如此多方致力于促进天主教会的发展和光荣，并竭力以各种方式传扬救主的道理，却远未忽视世俗事务；相反，在这一方面，他也以不懈的努力，接连不断地向各省人民施予各种恩惠。一方面，他对臣民的普遍福祉表现出慈父般的关怀；另一方面，他也会以各种荣誉款待自己熟识的个人；每一次都以真正高贵的精神施予恩惠。无人能在向皇帝请求恩惠时得不到所求；无人期望从他那里得到恩赐而落空。有人得到金钱赏赐，有人得到土地；有人获得\OriginalTerm{禁军长官}{Praetorian praefecture}之职，有人获得\OriginalTerm{元老身份}{senatorial}，还有人获得\OriginalTerm{执政官级}{consular rank}；许多人被任命为行省总督；另一些人被任命为\OriginalTerm{伯爵}{counts}，分为一级、二级或三级；无数人被授予\OriginalTerm{至显贵}{Most Illustrious}称号及其他许多荣衔；因为皇帝设立了新的爵位，以便使更多人蒙受他的恩宠之标记。\par

% structure: p0004 source=h2 target=subsection reason=relative-heading-rank; base=h2

\subsection{第二章　减免四分之一赋税}

他如何谋求普遍的幸福和繁荣，可从一件最有益且最具普遍性的措施看出，此措施至今仍被感恩地铭记。他免除了每年应缴地税的四分之一，将其赠予土地所有者；因此，若计算这每年的减免，便会发现耕种者每四年可享有一年免除赋税的收入。这项特权以法律确立，并为未来而保障，使得皇帝的仁惠不仅为当时的一代人所铭记，而且也为他们的子女和后代永远纪念。\par

% structure: p0006 source=h2 target=subsection reason=relative-heading-rank; base=h2

\subsection{第三章　平衡更沉重的赋税}

此外，有些人曾对先前皇帝统治时期所做的土地测量提出异议，抱怨其财产负担过重；在此事上，他也本着公正的原则，派遣专员去平衡赋税，并确保那些提出申诉的人获得豁免。\par

% structure: p0008 source=h2 target=subsection reason=relative-heading-rank; base=h2

\subsection{第四章　他用私人财产对钱财诉讼的败诉者施行慷慨}

在司法仲裁案件中，为了使因他的判决而败诉的人离开时不像获胜的诉讼人那样不满意，他亲自从自己的私人财产中，在有些情况下赐予土地，在另一些情况下赐予金钱，给予败诉的一方。他以这种方式确保败诉者，因曾出现在他面前，能像诉讼的获胜者一样满意；因为他认为，在任何情况下，都不应有人从与如此君王的会面中沮丧悲伤地退去。因此，诉讼双方都带着愉快和欣喜的面容离开审判现场，而皇帝高尚的慷慨赢得了普遍的赞叹。\par

% structure: p0010 source=h2 target=subsection reason=relative-heading-rank; base=h2

\subsection{第五章　藉我们救主的标记战胜斯基泰人}

我为何要即使简短和附带地叙述，他如何使野蛮民族臣服于罗马的权势；他如何是第一个征服从未学会服从的\OriginalTerm{斯基泰}{Scythian}和\OriginalTerm{萨尔马提亚}{Sarmatian}部落的人，并迫使他们，无论多么不愿意，承认罗马的统治？因为在他之前的皇帝实际上曾向斯基泰人纳贡；罗马人每年缴纳贡赋，承认自己为野蛮人的奴仆；这一耻辱，我们的皇帝再也无法忍受，也不能认为与他胜利的历程相称，继续支付他前任所缴纳的贡赋。因此，他完全信赖他的救主的助佑，也向这些敌人举起了他得胜的旌旗，很快便使他们全部归顺；对顽强抵抗他权威的人，他以武力镇压；另一方面，他通过明智派遣的使团安抚其余的人，使他们从无法无天和野蛮的生活中被改造到有秩序和文明的状态。这样，斯基泰人终于学会了承认服从罗马的势力。\par

% structure: p0012 source=h2 target=subsection reason=relative-heading-rank; base=h2

\subsection{第六章　因奴隶叛乱而征服萨尔马提亚人}

至于萨尔马提亚人，天主亲自将他们置于君士坦丁的统治之下，并以如下方式制服了一个充满蛮族傲慢的民族。他们受到斯基泰人的攻击，便将武器交给奴隶以击退敌人。这些奴隶首先战胜了入侵者，然后调转武器攻击他们的主人，将主人全部驱逐出家园。被驱逐的萨尔马提亚人发现他们唯一的得救希望在于君士坦丁的保护；而他，惯于拯救人的性命，将他们全部接纳到罗马帝国的疆界内。凡能服役的人，他编入自己的军队；其余的人，他分配土地供他们耕种以自养；因此他们自己承认，过去的灾难产生了幸福的结果，因为他们现在享受罗马的自由，而非野蛮的蛮族生活。天主以这种方式将许多不同的蛮族部落加添到他的疆域。\par

% structure: p0014 source=h2 target=subsection reason=relative-heading-rank; base=h2

\subsection{第七章　来自不同野蛮民族的使者从皇帝接受礼物}

事实上，使节们不断从各国来到，献上他们最珍贵的礼物作为君王的接纳。因此，我曾亲自站在皇宫入口附近，观察到一群引人注目的蛮族随从，他们的服饰和装饰各不相同，发式和胡须样式也迥异。他们面目狰狞可怖，身材高大异常：有些肤色发红，有些白如雪，还有一些介于两者之间。因为在我提及的那些人中，可以看到\Place{布莱米}部落、\Place{印度}人和\Place{埃塞俄比亚}人的标本，\Quote{那个广泛分散的民族，人类中最遥远的族群。}所有这些人都像一幅绘制的行列，依次向皇帝献上各自民族最珍视的礼物：有的献上金冠，有的献上镶有宝石的冠冕；有的带来金发男孩，有的带来绣金织花的蛮族服装；有的带来马匹，有的带来盾牌和长矛，以及弓箭，以此表示他们的效劳和联盟，供皇帝接纳。他分别接受这些礼物，并仔细收藏，以如此慷慨的方式回应，使赠礼者顿时富有。他还以罗马的高贵官职尊荣他们中最显贵者；因此他们中许多人从此更愿意留居我们中间，不再有重返故乡的愿望。\par

% structure: p0016 source=h2 target=subsection reason=relative-heading-rank; base=h2

\subsection{第八章 他也写信给曾派使节来访的波斯王，为他境内的基督徒说情。}

波斯王也通过派遣使节和赠礼作为和平与友谊的保证，表示愿意与君士坦丁结盟。在处理这项条约时，皇帝在答谢礼物的豪华程度上远远超过了首先向他表示敬意的这位君主。他还听说波斯有许多天主的教会，并且那里有大量人群聚集在基督的羊栈内，为此消息满心欢喜，决定将他对普世福祉的关怀也延伸到那个国家，因为他的目标是同等关怀每一个国家中的所有人。\par

% structure: p0018 source=h2 target=subsection reason=relative-heading-rank; base=h2

\subsection{第九章 君士坦丁奥古斯都致波斯王沙普尔的信函，内中含有一个真正虔诚的对天主和基督的宣认。}

\WorkTitle{致波斯王信函抄本}\par

作为对神圣信德的持守，我得以分享真理之光；在真理之光的引导下，我增进对神圣信德的认识。因此，正如我的行为本身所表明的，我信奉最圣洁的宗教；我宣称这一敬礼即是教导我更深入认识至圣天主的敬礼；凭借他的神圣大能，从大洋的边界开始，我逐一唤醒了世界每个民族，使他们获得稳固的希望；因此那些曾在最残酷的暴君奴役下呻吟、因日常苦难的压迫而几乎完全毁灭的民族，通过我的努力被恢复到了远为幸福的状态。我承认自己不断尊崇和纪念这位天主；我喜欢以纯洁无邪的思想在高天之上默观他的荣耀。\par

% structure: p0021 source=h2 target=subsection reason=relative-heading-rank; base=h2

\subsection{第十章 作者谴责偶像，并光荣天主。}

我屈膝呼求这位天主，对祭献的血腥、它们污秽可憎的气味以及一切妖术的烈火深感恐惧；因为那些被这些仪式玷污的亵渎不敬的迷信，已经使外邦世界的许多人，甚至整个民族堕落并陷入毁灭。因为万有的主不能容忍那些他出于自己的仁爱和对人类需求的考虑而启示给众人使用的恩赐，被歪曲成满足任何人私欲的工具。他对人的唯一要求是心灵的纯洁和灵魂的无玷；他以此标准衡量德行和虔诚的行为。因为他喜悦于节制与温和的行为；他爱慕谦和的人，憎恶骚动的心；他喜悦信德，惩罚不信；一切傲慢的权力都被他摧毁，他报复骄傲者的狂妄自大。虽然傲慢自大的人被彻底推翻，他却以应得的赏报酬报谦卑和宽恕的人；他甚至如此尊崇公正统治的王国，并以特殊的帮助加强它，使明智的君王安享和平的宁静。\par

% structure: p0023 source=h2 target=subsection reason=relative-heading-rank; base=h2

\subsection{第十一章 反对暴君和迫害者；论瓦莱里安的被掳。}

因此，我的兄弟，我相信我承认这独一天主——万物的创造者和父——并没有错；我许多在权力上的前任，被错误的疯狂引入歧途，竟敢否认他，但他们全都遭到了如此可怕和毁灭性的报应，以致于后世所有世代都将其灾祸视为对任何想要步其后尘者的最有效警告。在这些人中，我相信那位被神圣惩罚的闪电从我们这里驱赶出去、放逐到你的领土上的人，他的耻辱为你那著名的胜利增添了光彩。\par

% structure: p0025 source=h2 target=subsection reason=relative-heading-rank; base=h2

\subsection{第十二章 他声明，目睹迫害者的覆灭后，如今他为基督徒享有的和平而喜乐。}

确实，如今在我们这个时代，我所描述的那类人的惩罚公之于众，这是一件幸事。因为我自己亲眼目睹了那些最近以不敬神的谕令骚扰敬拜天主者的结局。为此，应感谢天主，因为他卓越的眷顾，所有遵守他神圣律法的人都因重新享受和平而欢欣。因此我完全相信一切事物都处于最好、最安全的状态，因为天主正借着他们纯洁而忠信的宗教服务和在其神圣品格上意见一致的判断，施恩将所有的人聚集到自己身边。\par

% structure: p0027 source=h2 target=subsection reason=relative-heading-rank; base=h2

\subsection{第十三章 他恳请对其国内的基督徒施以慈爱关怀。}

\Quote{请你想象一下，我听见这样符合我心愿的消息时有多么喜悦：波斯最美丽的地区充满了那些人，我目前只为他们说话，我是指基督徒。因此，我祈求您和他们都享有丰盛的繁荣，并且您的福祉和他们的福祉能均等。因为这样，您将体验到那位是并主和众父之父的天主的慈悲和恩宠。如今，因为您的权力很大，我将这些人委托给您保护；因为您的虔诚显著，我将他们交托给您照料。请用您一贯的仁慈和善良爱护他们；因为通过这个信德的证据，您将为自己和我们两方得到无法衡量的利益。}\par

% structure: p0029 source=h2 target=subsection reason=relative-heading-rank; base=h2

\subsection{第十四章　君士坦丁的热切祈祷如何为基督徒带来平安}

因此，世界各国如同受单一舵手的技艺引导，都自愿服从作为天主仆人的统治者的管理，罗马帝国的和平一直未受打扰，他的所有属下均享有安宁和休息的生活。与此同时，皇帝相信虔诚之人的祈祷有助于维持公共福利，因而感到必须热切寻求这样的祈祷；不但自己祈求天主的帮助和恩宠，而且命令教会的主教们为他献上祈祷。\par

% structure: p0031 source=h2 target=subsection reason=relative-heading-rank; base=h2

\subsection{第十五章　他命令在自己的硬币和画像上以祈祷的姿态表现自己}

从他命令将自己的肖像印在帝国的金币上，眼睛仰望，如同祈祷天主的姿态；这种钱币在整个罗马世界流通。此外，在某些城市的宫殿入口处悬挂着他的全身画像，眼中仰望天堂，手伸开如祈祷状，可以看出他的灵魂被神圣信德的能力深深打动了。\par

% structure: p0033 source=h2 target=subsection reason=relative-heading-rank; base=h2

\subsection{第十六章　他通过法律禁止在偶像神庙中放置自己的画像}

他以这种方式，甚至通过绘画的媒介，表现自己惯常祈祷天主。与此同时，他通过明确的法令禁止在任何偶像神庙中放置自己的任何肖像，为的是连他个人的轮廓也不会受到禁止的迷信错误的污染。\par

% structure: p0035 source=h2 target=subsection reason=relative-heading-rank; base=h2

\subsection{第十七章　论他在宫殿中的祈祷和阅读圣经}

那些注意到他如何将自己的宫殿塑造成天主的教堂，并自己为其中聚集的人们提供热心的榜样的人，可以看见他虔诚的更高尚的证据：他如何将圣经拿在手中，专心研读这些受到天主默感的神圣话语；然后他会与皇室所有成员一起定期祈祷。\par

% structure: p0037 source=h2 target=subsection reason=relative-heading-rank; base=h2

\subsection{第十八章　他命令普遍守主日和预备日}

他也规定一天应被视为祈祷的特殊日子：我是指真正为所有日子的首先和最重要的那一天，我们的主和救主的日子。他的全家都委托给执事和其他奉献于天主的侍者，他们因严肃的生活和各种美德而显得卓越；而他忠诚的护卫队，对他个人充满爱情和忠诚，发现他们的皇帝是虔诚实践的教导者，并像他一样崇敬主的救赎日，在那一天行他喜爱的祈祷。这位有福的君主也向他的所有属下推荐同样的守节；他的热切希望是逐步引导全人类崇拜天主。因此他命令罗马帝国的所有属下守主日为休息日，并崇敬安息日前的那一天；我想，是为了纪念人类的救主在那一天所完成的事。由于他希望教导他的整个军队热心崇敬救主的日子（该日名称来自光和太阳），他自由给予那些参与神圣信德的人有闲暇参加天主的教会的礼仪，以便他们能够毫无阻碍地完成宗教崇拜。\par

% structure: p0039 source=h2 target=subsection reason=relative-heading-rank; base=h2

\subsection{第十九章　他命令连他的异教士兵也要在主日祈祷}

至于那些仍不知天主真理的人，他通过第二条法令规定，他们应在每个主日出现在城市附近的开阔平原上，在信号给出时齐声向天主祈祷他们先前学习过的祈祷。他劝告他们，他们的信心不应当建立在他们的长矛、盔甲或肉体强壮之上，而应当承认至高的天主为所有善事和胜利本身的赐予者；他们应当定时向天主举手祈祷，仰望天堂，抬起心灵更高地望向天上之王，祈求他作为胜利的主宰、他们的保存者、护卫者和帮助者。皇帝本人规定了所有军队要用的祈祷，命令他们用拉丁文宣读以下词语：\par

% structure: p0041 source=h2 target=subsection reason=relative-heading-rank; base=h2

\subsection{第二十章　君士坦丁给他的士兵规定的祈祷格式}

\Quote{我们承认你为唯一的天主；我们称你为我们的王，恳求你的援助。凭你的恩宠，我们得胜；通过你，我们比我们的敌人更强大。我们感谢你过去的恩赐，并信任你将来的福祉。我们齐声向你祈祷，恳求你长久保守我们的皇帝君士坦丁和他虔诚的儿子们安全胜利。}这就是他的军队在主日要完成的义务，也是他们被教导要向天主祈祷的内容。\par

% structure: p0043 source=h2 target=subsection reason=relative-heading-rank; base=h2

\subsection{第二十一章　他命令将救主的十字架标志铭刻在他的士兵的盾牌上}

不但如此，他还命令将救赎的胜利标志铭刻在他的士兵的盾牌上；并命令他的战斗部队在行进时，前导的不是像以往那样的金像，而只有十字架旗帜。\par

% structure: p0045 source=h2 target=subsection reason=relative-heading-rank; base=h2

\subsection{第二十二章　论他在祈祷中的热心和对复活节的崇敬}

皇帝本人，作为我们宗教神圣奥迹的分享者，每日在规定时辰退入皇宫最内室，在那里独与天主交谈，屈膝恳求，祈求所需的恩惠。尤其在救恩的复活节，他加倍虔诚，如同履行祭司的职责，尽心竭力，并在这场庆节的纪念中超越众人。他还将神圣的守夜变为白昼般明亮，命人在全城点燃极长的蜡炬；此外，火炬处处散布光明，使这神秘的守夜璀璨辉煌，胜过白昼。及至白昼复临，他效法救主的恩慈行为，向各个民族、行省和人民的臣民慷慨施予，普施厚惠。\par

% structure: p0047 source=h2 target=subsection reason=relative-heading-rank; base=h2

\subsection{第二十三章 他如何禁止偶像崇拜，却尊敬殉道者与教会庆节。}

他在事奉天主方面，便是这样神圣地服务。同时，帝国境内所有臣民，无论文职或军职，都发现处处设防抵御偶像崇拜，各种祭献均被禁止。又通过一项法令，规定必须遵守主日，并传达给各省总督；他们遵皇帝之命，尊重纪念殉道者的日子，并适当尊敬教会中的庆节时期。所有这些意旨均得以完全实现，令皇帝十分满意。\par

% structure: p0049 source=h2 target=subsection reason=relative-heading-rank; base=h2

\subsection{第二十四章 他称自己为主教，负责教会外部事务。}

因此，有一次当他宴请一群主教时，脱口说出\Quote{他自己也是一位主教}，在我听来，他以下面的话对他们说：\Quote{你们是管辖教会内部的主教；而我也是主教，由天主派立，照管教会外部的一切。}他的措施确实与他的话相符：因为他以主教般的关怀照管他的臣民，并尽力劝勉他们度虔敬的生活。\par

% structure: p0051 source=h2 target=subsection reason=relative-heading-rank; base=h2

\subsection{第二十五章 禁止祭献、秘仪、角斗士搏斗，以及尼罗河的淫秽崇拜。}

与这种热忱一致，他接连颁布法律和敕令，禁止任何人向偶像祭献、咨询占卜者、设立神像，或以角斗士的血腥搏斗污染城市。此外，由于埃及人，尤其是亚历山大里亚人，惯常通过由柔弱的男子组成的祭司团来尊崇他们的河流，又通过一项法律，命令将整个这类邪恶之人除灭，使此后无人再受类似不洁之毒的污染。尽管迷信的居民担心河流会因此停止其惯常的泛滥，但天主亲自显示了他对皇帝律法的认可，安排一切事物完全出乎他们意料。因为那些以其邪恶行为污染城市的人果然不再出现；但河流，仿佛流经之地已得净化以迎接它，涨得比以往更高，以其肥沃的水流完全漫溢大地。这有效地告诫受骗的民众，要远离不洁之人，并将他们的昌盛唯独归于那赐予万善者。\par

% structure: p0053 source=h2 target=subsection reason=relative-heading-rank; base=h2

\subsection{第二十六章 关于无子者及遗嘱法的现行法律修订。}

皇帝在各行省所施的这类恩惠确是如此众多，足以为任何想要记载它们的人提供丰富的素材。其中可以列举那些他完全重新制定并建立在更公平基础上的法律：其改革性质可以简要而容易地说明。无子者在旧法下被剥夺继承祖产的权利，这是一条残酷的法令，将他们视为实际上的罪犯。皇帝废除了这条法律，并规定处于这种情形的人可以继承。他基于公平正义的原则处理这个问题，认为故意犯罪者应受到与其罪行相应的惩罚。但自然本身拒绝给予许多人子女，他们或许渴望众多后代，却因身体虚弱而失望。另一些人无子，并非因为厌恶后代，而是因为他们对\OriginalTerm{哲学}{philosophy}的热爱使他们不情愿婚姻结合。还有妇女献身于天主的事业，保持了纯洁无瑕的童贞，将自己的灵魂和身体完全奉献给贞洁和圣善的生活。那么，这种行为应被视为有罪，还是更应被钦佩和赞美呢？因为渴望这种状态本身就是可敬的，而维持它则超越了自然本身的力量。确实，那些因身体虚弱而失去生育希望的人值得同情，而非惩罚；而献身于更高目标的人不应受到惩罚，而应得到特别的钦佩。基于这些合理的理性原则，皇帝纠正了这条法律的缺陷。再者，关于临终者的遗嘱，旧法规定即便在最后一息，也必须以某些特定的言辞表达，并规定了确切的格式和用语。这种做法导致了许多欺诈行为，阻碍了死者意愿的完全实现。我们的皇帝一旦意识到这些弊端，也改革了这条法律，宣布临终者应被允许以尽可能少的言语、用他喜欢的任何词语表明他最后的愿望；并可以任何书面形式，甚至口头形式立下遗嘱，只要在适当的见证人面前进行，以便他们能够忠实地履行其职责。\par

% structure: p0055 source=h2 target=subsection reason=relative-heading-rank; base=h2

\subsection{第二十七章 他在其他法令中，规定基督徒不得为犹太人奴仆，并确认大公会议决议的有效性。}

他还通过了一项法律，规定任何基督徒不得因犹太主人而处于奴役之下，理由是：那些被救主所赎回的人，不应被那些杀害先知和主本身的民族置于奴隶制的轭下，这是不合理的。如果此后发现任何人处于这种情况，奴隶应获得自由，主人应被处以罚款。\par

他还将自己的权柄赋予主教们在教务会议上通过的决议，并禁止各省总督废止他们的任何法令：因为他视天主的司铎高于任何审判官。他为臣民的利益制定了一千条类似的规定；但现在没有时间详细描述它们，以传达他在这方面的帝王智慧；现在我也不必详述，作为至高天主的忠实仆人，他如何从早到晚致力于寻找施恩的对象，以及他对所有人如何平等而普遍地仁慈。\par

% structure: p0058 source=h2 target=subsection reason=relative-heading-rank; base=h2

\subsection{第二十八章 他给教会的恩赐，以及对贞女和穷人的慷慨}

然而，他的慷慨尤其体现在对天主的教会上。在某些情况下，他赐予土地；在其他情况下，他发放食物供应以支持穷人、孤儿和寡妇；此外，他还表现出极大的关心和远见，为赤身露体和贫困的人提供衣服。然而，他特别尊敬那些献身于实践神圣哲学的人。因此，他对天主最圣洁、永贞的团体表示出近乎崇敬的尊重：因为他确信，这些人所奉献的天主自己就住在他们的灵魂里。\par

% structure: p0060 source=h2 target=subsection reason=relative-heading-rank; base=h2

\subsection{第二十九章 关于\Person{君士坦丁}的演讲和演说}

对于他自己，他有时彻夜不眠，用神圣知识充实自己的心灵；他花了很多时间撰写演讲，其中许多他公开宣讲；因为他认为，他有责任通过诉诸理性来治理臣民，并在所有方面确保对他的权威的理性服从。因此，他有时会亲自召集集会，届时会有大量人群参加，希望听到一位皇帝扮演哲学家的角色。如果在演讲过程中有任何机会触及神圣话题，他立即站直身体，以严肃的表情和低沉的声音，仿佛虔诚地引导听众进入神圣教理的奥秘；当他们以欢呼声向他致意时，他会用手势指示他们举目望天，将钦佩留给至高君王独自，并以崇拜和赞美来尊崇他。他通常将演讲的主题分为几部分：首先彻底驳斥多神教的错误，证明外邦人的迷信不过是欺诈和掩饰不敬虔；然后他断言天主的独一主权；由此转向他的眷顾，既有普遍的也有特殊的；接着进入救恩的安排，他展示其必要性及其与情况的适应性；然后依次进入天主审判的教理。在这里他特别有力地打动听众的良心，同时谴责贪婪和暴力的人，以及那些沉溺于无度贪欲的人。他甚至使一些在场的熟人感到他话语的严厉鞭笞，使他们因内疚而低头站立，因为他以最清晰和最令人印象深刻的方式作证反对他们，说他们要为他们的行为向天主交账。他提醒他们，天主亲自将世界的帝国赐给他，而他本人基于同样的神圣原则，将其中一些部分委托给他们治理；但所有人在适当的时候都将同样被召唤向万物的至高主宰汇报他们的行为。这就是他持续的见证；这就是他的劝诫和教导。他本人以真诚的信德感受并表达这些情感；但他的听众却不愿学习，对忠告充耳不闻；他们虽然高声喝彩接受他的话，但被无法满足的贪欲驱使实际上无视了它们。\par

% structure: p0062 source=h2 target=subsection reason=relative-heading-rank; base=h2

\subsection{第三十章 他在一个贪婪的人面前标出一个坟墓的量度，从而使他羞愧}

有一次，他亲自这样对他的一个朝臣说：\Quote{我的朋友，我们要把无度的贪欲带到多远？} 然后他用偶然握在手中的长矛画出一个人的尺寸，继续说道：\Quote{即使你能获得这整个世界的一切财富，是的，整个世界本身，你最终带走的也不会超过我标出的这个小点，如果那真的属于你的话。}\par

% structure: p0064 source=h2 target=subsection reason=relative-heading-rank; base=h2

\subsection{第三十一章 他因过度仁慈而被嘲笑}

与此同时，既然没有对死刑的恐惧来阻止犯罪，因为皇帝本人一贯倾向于仁慈，而且没有一个省总督用适当的刑罚来惩治罪行，这种情况给帝国的整体管理带来了不小的责备；是否公正，让每个人自己判断：就我自己而言，我只请求允许记录事实。\par

% structure: p0066 source=h2 target=subsection reason=relative-heading-rank; base=h2

\subsection{第三十二章 关于\Person{君士坦丁}写给圣人之集会的演讲稿}

皇帝习惯于用拉丁语撰写他的演讲稿，这些演讲稿由为此特殊任务指定的翻译者翻译成希腊语。其中一篇如此翻译的演讲稿，我打算作为样本附在本书中，即他题写为\Quote{对圣人的集会}并献给天主的教会的那一篇，以便没有人有理由认为我在这方面的见证只是空洞的赞美。\par

% structure: p0068 source=h2 target=subsection reason=relative-heading-rank; base=h2

\subsection{第三十三章 他如何站立聆听\Person{欧瑟伯}为纪念我们救主之墓所作的演讲}

然而，有一件事我却不可略而不记，就是这位令人钦佩的君王在我面前所做的一件事。有一次，我因对他的虔诚怀有信心，便大胆请求准允在他面前宣讲一篇关于救主坟墓的道理。他欣然允诺，并在皇宫内院中，当着许多听众的面，与其他听众一同站立聆听。我恳请他坐在旁边的御座上，但他不肯；他凝神静听我讲论的各个主题，并亲自为其中包含的神学真理作证。过了些时候，因演说篇幅甚长，我本欲结束，但他不许，并劝我继续讲完。我再次恳请他坐下，他却不悦，说，粗心大意地聆听有关天主之道的讨论是不对的；又说，这种姿势对他自己有益，因为恭敬地站立聆听神圣真理是合适的。于是，我讲完了道理，便回家，继续日常的工作。\par

% structure: p0070 source=h2 target=subsection reason=relative-heading-rank; base=h2

\subsection{第三十四章 他致信尤西比乌斯论复活节及圣经抄本}

君士坦丁常关心天主教会之福祉，他亲自致信于我，论及提供圣经抄本的事宜，并论及至圣的复活节庆。因我曾将自己对那节日奥秘含义的阐释献给他；他如何赐予我回信的殊荣，凡读下面信函者自可明了。\par

% structure: p0072 source=h2 target=subsection reason=relative-heading-rank; base=h2

\subsection{第三十五章 君士坦丁致尤西比乌斯书，称赞其论复活节演讲}

\Signature{胜利者君士坦丁，至大奥古斯都，致尤西比乌斯。}\par

\Quote{诚然，要适当地论述基督的奥秘，并以相称的方式解释关于复活节庆的争议——其起源及其珍贵而艰苦的实现——这实在是一桩艰巨的工作，甚至非语言本身所能及。因为即便是那些能够理解天主事理的人，也不足以充分描述天主的事理。尽管如此，我仍满心钦佩你的学识与热忱，我不仅自己欣然阅读了你的作品，并且按照你自己的意愿，指示将它传递给许多我们神圣宗教的真诚信徒。既然我们如此喜悦地接受来自你的睿智的这份恩惠，请更频繁地以那些著作使我们喜乐，你承认自己自幼便受其训练，以致我像是在催促一个自愿的人——如人们所说——劝你从事你习惯的工作。确实，我们对你所持的高瞻远瞩的评判，证明那将你的著作译成拉丁文的人绝非不能胜任此任务，尽管这样的译本不可能完全等同原著本身的卓越。愿天主保佑你，亲爱的弟兄。} 他关于此事的信函如此。而关于提供圣经抄本以供在教会中诵读的信，内容如下。\par

% structure: p0075 source=h2 target=subsection reason=relative-heading-rank; base=h2

\subsection{第三十六章 君士坦丁致尤西比乌斯书论预备圣经抄本}

\Signature{胜利者君士坦丁，至大奥古斯都，致尤西比乌斯。}\par

\Quote{因天主我们的救主的恩眷眷顾，许多人已联合到那以我命名的城市中最神圣的教会。看来，既然那城在其余各方面都迅速繁荣，教会的数目也应当增加。因此，请你欣然接受我在这方面的决定。我想应指示你的睿智，命人用预备好的羊皮纸，以清晰的字迹，以方便携带的形式，由精于技艺的专业抄写员，抄写五十部圣经，其供应和使用你知道对教会的训导最为需要。本教区的\OriginalTerm{财务官}{catholicus}也已接获我们仁慈的函令，要小心供应一切所需以预备这些抄本；你则需特别留意，务使它们尽快完成。你凭此信也有权使用两辆公共马车运送，这样，抄本誊写清楚后，便极易送达为朕亲阅；你可委托你教会的一位执事担任此差，他抵达此处后，将得享朕的慷慨。愿天主保佑你，亲爱的弟兄。}\par

% structure: p0078 source=h2 target=subsection reason=relative-heading-rank; base=h2

\subsection{第三十七章 抄本如何备办}

皇帝的御令如此，随即执行了这工作本身；我们以精美装帧的三叠四叠本送至他处。此事实有皇帝另一封致谢信为证；他在信中听说我乡的君士坦提亚城——其居民素来沉迷迷信——因宗教感而放弃了昔日的偶像崇拜，便表示喜悦，并赞许他们的行为。\par

% structure: p0080 source=h2 target=subsection reason=relative-heading-rank; base=h2

\subsection{第三十八章 迦萨市镇如何因信基督教而升格为城，并得名君士坦提亚}

事实上，如今在巴勒斯坦省称为君士坦提亚的地方，因接受了救恩的宗教，既蒙天主眷顾，又得皇帝特别尊荣，首次升格为城，并以他虔敬的姊妹的更为尊贵的名字取代了旧名。\par

% structure: p0082 source=h2 target=subsection reason=relative-heading-rank; base=h2

\subsection{第三十九章 腓尼基某地亦升格为城，其他城市偶像崇拜被废止，教堂建成}

\Emph{类似的}变化也在其他几个城市中发生；例如，在腓尼基的那座以皇帝名字命名的城市，其居民将无数偶像投入火焰，而采纳了救恩的信德原则。其他省份的许多人，无论是在城市还是乡村，都成为寻求天主救恩知识的自愿追求者；他们视过去视为最神圣的各种偶像为无价值之物而加以摧毁；自愿拆毁其中包含的高大庙宇和神龛；并抛弃他们先前的观念，或者说是错误，开始并完全建立了新的教堂。但由于详细叙述这位虔诚君主的行为并非我分内之事，而是那些有幸始终与他相伴者的职责，我将满足于简要记述我个人所知的事实，然后继续叙述他生命的最后日子。\par

% structure: p0084 source=h2 target=subsection reason=relative-heading-rank; base=h2

\subsection{他在统治的三个十年期将凯撒尊位赐予三个儿子，并奉献耶路撒冷的教堂}

至此，他已完成在位第三十年。在此期间，他的三个儿子在不同时间被接纳为帝国的同僚。第一个儿子，与他同名的君士坦丁，大约在他统治第十年获得此殊荣。次子君士坦提乌斯，以其祖父命名，约在第二十年被宣布为凯撒；第三子君士坦斯，其名字表达了他品格的坚定与稳固，在他父亲统治第三十年时晋升到同一尊位。这样，他养育了三个后代，仿佛天主圣三一般的虔敬儿子们，并依次在每个十年期使他们分享他的帝国权威；他认为这正是三十年庆典，适合向万有的至高主宰感恩，同时相信，他在耶路撒冷以其热忱的慷慨所建造的教堂，其奉献礼也可有利地举行。\par

% structure: p0086 source=h2 target=subsection reason=relative-heading-rank; base=h2

\subsection{同时，他因埃及所起的争论命令在提洛召开主教会议}

\Emph{同时}，那敌视一切良善的嫉妒之灵，如同遮蔽最明亮阳光的乌云，试图再次挑起纷争扰乱埃及教会的安宁，破坏这一喜庆的欢乐。然而，我们蒙天主眷顾的皇帝又一次召集了许多主教的主教会议，让他们仿佛天主的军队一样武装起来，对抗这恶意的灵；命令他们从整个埃及和利比亚、从亚细亚和从欧罗巴前来，以便先解决争议，然后执行上述神圣建筑的奉献礼。他吩咐他们，在途中，要在腓尼基的首府城市调解分歧；提醒他们，在怀有相互敌意时，无权从事天主的服务，因为天主的律法明确禁止那些不和的人献礼物，除非他们先和好并彼此愿意和平（\BibleRef{玛 5:23-24}）。这些有益的训诫，皇帝不断生动地呈现在自己眼前，并据此劝诫他们以完全一致与和谐的精神承担当前的责任，在一封内容如下的信中。\par

% structure: p0088 source=h2 target=subsection reason=relative-heading-rank; base=h2

\subsection{君士坦丁致提洛主教会议的信}

\Signature{凯旋者君士坦丁，至尊奥古斯都，致提洛神圣主教会议。}\par

\Quote{诚然，使天主教会保持合一，使基督的仆人们在此时此刻免受一切指责，乃是最切合且最有利于我们这时代的昌盛。然而，有些人被一种恶毒而疯狂的争竞精神所驱使（因为我不愿指责他们故意过着与身份不符的生活），正试图制造那种在我看来是一切祸患中最有害的普遍混乱；因此，我劝告你们，你们既已积极，就当毫不拖延地聚集起来，召开一次会议：保护那些需要保护的人，为处于危险中的弟兄们提供救治，将分裂的成员召回判断的一致，趁仍有机会时纠正错误；这样，你们就能使如此众多的省份恢复那应有的和谐——而那和谐，奇怪而可悲的反常！竟被少数人的傲慢所破坏。我相信，所有人都一致认为，这条路同时取悦于全能的天主（也是我本人最高的愿望），而且，如果你们能成功地恢复和平，也将为你们自己带来不小的荣誉。那么，不要迟延，要以加倍的勤勉赶快结束当前的纷争，使之符合时机，以我们事奉的救主特别要求于我们的那种真正真诚和信德的精神聚集在一起——我几乎可以说，他是在所有场合以可听见的声音这样要求的。至于我，敬虔的热忱绝不会有所欠缺。我已经按照你们信中所指示的，尽了一切努力。我派遣了那些你们希望在场的诸位主教，让他们与你们共同商议。我派遣了一位执政官级的人物\Person{狄约尼削}，他将提醒那些应当与你们一同出席会议的主教们他们的职责，并亲临会议监督进程，尤其维持良好秩序。同时，若有人（虽然我认为极不可能）胆敢在此情况下违抗我的命令，拒绝出席，将会立即派遣使者，凭御令将那人流放，并教训他：在捍卫真理时，违抗皇帝的谕令是不适宜的。至于其余的事，将由各位圣座，不受敌意或偏袒的影响，而是依照教会和宗徒的秩序，为正犯的过错或不经意的失误，想出适当的补救办法；这样，你们就能使教会立即免受一切指责，减轻我的忧虑，并通过恢复和平的祝福给如今分裂的人，为你们自己赢得最高的荣誉。愿天主保佑你们，亲爱的弟兄们！}\par

% structure: p0091 source=h2 target=subsection reason=relative-heading-rank; base=h2

\subsection{第四十三章 来自各省的主教参加了\Place{耶路撒冷}圣堂的奉献礼。}

这些命令刚被付诸实施，又有一位使者带着皇帝的文牍到达，催促会议立即启程前往\Place{耶路撒冷}，不得迟延。于是，他们全都离开\Place{腓尼基}省，前往目的地，并利用公共交通工具。这样，\Place{耶路撒冷}成了来自各省的杰出主教们的聚集点，全城挤满了天主的仆人们的庞大集会。\Place{马其顿}人派来了他们都会的主教；\Place{潘诺尼亚}人和\Place{默西亚}人派来了他们当中最美善的天主青年羊群中的一员。还有一位来自\Place{波斯}的圣洁主教也在那里，他精通圣经；\Place{比提尼亚}和\Place{色雷斯}的主教们也光临了会议；来自\Place{基里基雅}的最有名望的主教也未曾缺席，\Place{卡帕多细亚}的领袖们也来了，他们以学识和口才尤为出众。简言之，整个\Place{叙利亚}和\Place{美索不达米亚}、\Place{腓尼基}和\Place{阿拉伯}、\Place{巴勒斯坦}、\Place{埃及}和\Place{利比亚}，连同\Place{特拜}的居民，都共同壮大了天主的仆人的庞大集会，他们身后还有来自各省的无数群众。他们由皇家护卫陪同，宫中还派来了值得信赖的官员，奉命以皇帝的费用来增添节日的辉煌。\par

% structure: p0093 source=h2 target=subsection reason=relative-heading-rank; base=h2

\subsection{第四十四章 关于公证员\Person{马里亚努斯}对他们的接待；对穷人的金钱分发；以及对圣堂的献礼。}

这些官员的指导和首领是皇帝一位极有用的仆人，一位在信德和虔诚方面卓越的人，并且精通圣经；在暴君当权时期，他曾因对信仰的宣认而荣耀地受人注目，因此理所当然地被委任安排当前的这些事务。因此，他忠实地服从皇帝的命令，以谦恭的款待接待了与会者，并以极其盛大的规模用宴席和宴会招待他们。他还向赤身露体和贫穷的人，以及因缺乏食物和日常生活必需品而受苦的男女群众大量分发金钱和衣物。最后，他用皇帝慷慨的献礼装饰并美化了圣堂本身，从而完全完成了他所受托的服务。\par

% structure: p0095 source=h2 target=subsection reason=relative-heading-rank; base=h2

\subsection{第四十五章 与会主教的各种演讲；以及本史作者\Person{欧瑟比乌斯}的演讲。}

同时，因天主仆人的祈祷和讲道，这庆节更加光彩。其中一些人颂扬这位虔诚皇帝的慷慨奉献，他将自己完全献于人类的救主，并颂扬他为纪念救主所建造的宏伟建筑；另一些人则以关于我们宗教神圣教义的公开讨论，仿佛向在场所有人的耳朵提供了一场属灵的盛宴；还有些人诠释圣经段落，揭示其隐藏的意义；而那些无力做到这些的人，则以他们为普遍和平、为天主的教会、为作为如此多恩宠工具上的皇帝本人、以及为他虔诚的儿子们奉献的祈祷，向天主献上无血的祭献和奥妙的礼仪。我自己虽然不配享有这样的特权，也发表了各种公开演说以尊崇这一盛典，其中一部分通过书面描述解释了皇帝建筑的细节，另一部分则试图从先知异象中为建筑所展示的象征收集恰当的例证。就这样，在我们皇帝在位第三十年，基督圣殿奉献的庆典在喜乐中得以举行。\par

% structure: p0097 source=h2 target=subsection reason=relative-heading-rank; base=h2

\subsection{第四十六章 \Person{优西比乌} 随后将其关于救主圣殿的描述和三秩庆辰的演说当面呈献给 \Person{君士坦丁}}

关于我们救主圣堂的结构、圣岩的形状、工程本身的辉煌以及金银宝石无数的奉献，我已尽我所能地进行了描述，并以单篇论著形式奉献给皇帝。我将在适当的机会将其附录于本作品。同时，我还要附上我关于他登基三十周年庆辰的演说辞，这篇演说辞在不久之后，当我旅行至以他命名的城市时，在皇帝本人面前发表。这是我第二次有机会在皇宫本身荣耀至高的天主；当时我虔诚的听众表现出最大的喜乐，他后来证实了这一点，当时他款待了在场的主教们，并以各种尊荣厚待他们。\par

% structure: p0099 source=h2 target=subsection reason=relative-heading-rank; base=h2

\subsection{第四十七章 \Place{尼西亚} 公会议在 \Person{君士坦丁} 在位第二十年举行，耶路撒冷圣殿奉献则在第三十年}

这第二次公会议由皇帝在耶路撒冷召集，就我们所知，它是仅次于他曾在著名的 \Place{比提尼亚} 城市召集的第一次公会议的最大会议。那一次确实是胜利的集会，在他即位第二十年举行，是一次在胜利之城为战胜敌人而感恩的庆典。本次集会则为第三十周年增添了光彩，期间皇帝在我们救主坟墓处的圣殿举行奉献礼，作为向施与一切美善的天主所献的和平祭。\par

% structure: p0101 source=h2 target=subsection reason=relative-heading-rank; base=h2

\subsection{第四十八章 \Person{君士坦丁} 对过分赞美他的人不悦}

当所有这些仪式都完成后，皇帝品德的神圣特质继续成为普遍赞美的主题。其中一位天主的仆人竟然大胆到在他面前称他有福，认为他配得在此世拥有绝对而普世的权柄，并在来世注定要与天主圣子共享王权。然而，\Person{君士坦丁} 听到这些话时感到愤怒，禁止那位说话者使用这样的言辞，反而劝勉他恳切地为他祈祷，希望无论是今生还是来世，他都能配得成为天主的仆人。\par

% structure: p0103 source=h2 target=subsection reason=relative-heading-rank; base=h2

\subsection{第四十九章 其子 \Person{君士坦提乌斯·凯撒} 的婚姻}

在他即位第三十年结束时，他为他的次子主持了婚礼，而长子的婚礼早已完成。这是一个充满极大喜乐和欢庆的场合，皇帝亲自在仪式上陪同他的儿子，并以丰盛的款待宴请男女宾客，男女分席。同样，大量的礼物丰裕地分发给各城市和人民。\par

% structure: p0105 source=h2 target=subsection reason=relative-heading-rank; base=h2

\subsection{第五十章 来自印度人的使节和礼物}

大约此时，来自居住在遥远东方地区的印度人的使节抵达，带来了各种明亮宝石和不同于我们已知物种的动物等礼物。他们将这些奉献献给皇帝，从而承认他的主权甚至延伸至印度洋，并承认他们国家的君主们通过画像和雕像向他表示敬意，承认他的帝国和最高权威。因此，东方的印度人现在归顺于他的统治，正如西方大洋的不列颠人在他统治初期所做的那样。\par

% structure: p0107 source=h2 target=subsection reason=relative-heading-rank; base=h2

\subsection{第五十一章 \Person{君士坦丁} 将帝国分给三个儿子，并教导他们治国与宗教}

这样，他既已在世界的两个极端建立了自己的权力，便将整个治下的全部领土如同分配遗产一样分给了他的三个儿子，这些儿子是他最亲爱的关切对象。他将祖父的份额分给长子，将帝国的东部地区分给次子，将位于这两部分之间的国家分给第三子。由于渴望赐予他的孩子们一种对灵魂真正宝贵和有益的遗产，他谨慎地用真正的宗教原则来熏陶他们，亲自做他们认识神圣事物的导师，并任命有可靠虔敬的人为他们的教师。同时，他为他们指派了最有成就的世俗学问教师，其中一些人教导他们战争技艺，另一些人则训练他们政治学，还有一些人教授法律科学。此外，每个人都得到真正君王般的随从，包括步兵、持矛兵和卫兵，以及各种其他军事力量；分别由皇帝曾亲自验证其军事技能和对儿子们忠诚的将领、护民官和将军指挥。\par

% structure: p0109 source=h2 target=subsection reason=relative-heading-rank; base=h2

\subsection{第五十二章 他们成年后，他仍是他们在虔敬方面的导师}

当凯撒们尚处幼年时，便有合适的顾问辅佐他们处理国务；及至成年，他们父亲的训诲便已足够。父亲在场时，以自身榜样激励他们，劝勉他们追随其虔诚的道路；不在时，则通过书信为他们提供适合其帝王身份的处事准则，其中首要且最大的叮嘱是，要珍视对万有之主的知识和敬拜，胜过财富，甚至胜过帝国本身。最后，他允许他们不受约束地掌管帝国的公共行政，但他的首要请求是，他们要关心天主的教会的事奉，并勇敢地自称为基督的门徒。如此受训，与其说是因教诲而顺从，不如说是出于他们对德性的自发渴望，他的儿子们不仅超额完成了父亲的叮嘱，而且全心致力于事奉天主，甚至在皇宫内也与全家成员一同遵守教会的规条。因为父亲的远见已确保他儿子们的所有侍从都是基督徒。不仅如此，最高级别的军官和掌管公共事务的人也都是同一信德的宣认者：因为皇帝信赖那些献身于天主事奉之人的忠诚，如同坚固可靠的保障。当这位三倍蒙福的君主完成这些安排，从而为整个帝国确保了秩序与安宁时，万恩的施予者天主认为将他迁至更美产业的时机已到，便召他偿还自然之债。\par

% structure: p0111 source=h2 target=subsection reason=relative-heading-rank; base=h2

\subsection{在位约三十二年，享寿六十余岁，身体依然健康}

他完成了三十二年统治的年限，只差几个月零几天，而他整个寿命大约是其两倍。在这个年纪，他仍拥有健康强健的身体，毫无瑕疵，且比青年更有活力；仪表高贵，体力足以胜任任何辛劳；因此他能参加军事训练、骑马、忍受旅途劳累、投入战斗，并为被他征服的敌人竖立战利品，除此之外，他还赢得了那些无血的胜利，他惯常以此战胜那些反对他的人。\par

% structure: p0113 source=h2 target=subsection reason=relative-heading-rank; base=h2

\subsection{论那些因贪婪和伪善而滥用其极度仁爱之人}

同样，他的精神品质也达到了人类完美的最高点。事实上，他因每一种品德而卓越，但尤其以仁爱著称；然而，这一德行却使许多人指责他，因为邪恶之人的卑劣，他们将自身的罪行归咎于皇帝的宽容。说实话，我自己可以为在此期间盛行的严重罪恶作证；我是指贪婪和无原则之人的暴力，他们同样侵害社会各阶层，以及那些潜入教会、冒用基督徒名称和品格的可耻伪善者。他自身的仁爱和心地善良，他信德的真诚和品格的诚实，使皇帝相信这些 \Quote{自称基督徒之人} 的宣认，他们狡猾地保持着对他个人真心爱戴的表象。他对这种人的信任有时迫使他做出与自己身份不相称的行为，嫉妒趁机在这件事上玷污了他品格的荣光。\par

% structure: p0115 source=h2 target=subsection reason=relative-heading-rank; base=h2

\subsection{君士坦丁直至生命终结仍从事各种著作}

然而，这些犯罪者很快便受到了天主的惩罚。回到我们的皇帝。他如此彻底地训练了自己的推理心智，以至于直到最后他仍在撰写各种主题的论述，经常当众发表演说，并教导听众宗教的神圣教义。他还习惯于从事政治和军事问题的立法；总之，为人类整体的福祉谋划一切。值得特别注意的是，就在他离世前不久，他在惯常的听众面前发表了一篇葬礼演说，在其中详细论述了灵魂的不朽、坚持敬虔生活之人的状态，以及天主为爱祂之人所预备的福分。另一方面，他通过丰富而确凿的论证，揭示了那些走相反道路之人的结局，生动地描述了不敬虔之人的最终毁灭。他在这些主题上的有力见证似乎深深触动了周围人的良心，以至于一位自诩的哲学家，当他被问及对所听内容的看法时，证明了他的话语真实，并虽不情愿却真正地赞扬了他驳斥多神崇拜的论证。皇帝在去世前与朋友们进行这样的交谈，仿佛为自己过渡到更幸福的生活铺平了道路。\par

% structure: p0117 source=h2 target=subsection reason=relative-heading-rank; base=h2

\subsection{如何率领主教们远征波斯，并携带一座教堂形式的帐篷}

同样值得记录的是，大约在我写作此刻之时，皇帝听说东方有一些野蛮人叛乱，便认为战胜这一敌人仍在前方等待着他，决定远征波斯人。于是他立即调动军队，同时将他计划中的行军告知当时在他宫廷中的主教们，他认为应该带其中一些人作为同伴和一些必要的事奉天主的助手。他们则欣然表示愿意跟随他，声称不愿离开他，并承诺藉着为他向天主祈祷来与他一同并为他作战。皇帝听到这个回答，满心喜悦，向他们透露了他计划中的行军路线；随后他令人准备了一座极其华丽的帐篷，其形状宛如一座圣堂，供他自己在即将到来的战争中使用。他打算在此与主教们一同向那位赐予一切胜利的天主祈祷。\par

% structure: p0119 source=h2 target=subsection reason=relative-heading-rank; base=h2

\subsection{第五十七章 他如何接见波斯使节，并在复活节与其他人在夜间守夜。}

同时，波斯人听闻皇帝的战备，对与他的军队交战的前景感到十分惊恐，于是派遣使节前来请求和平条件。皇帝本人是真诚的和平爱好者，立即接受了这些提议，并欣然与该民族建立了友好关系。此时，隆重的复活节即将到来；他在这个节日向天主献上祈祷的贡品，并与其他人一起在夜间守夜。\par

% structure: p0121 source=h2 target=subsection reason=relative-heading-rank; base=h2

\subsection{第五十八章 论\Place{君士坦丁堡}为纪念宗徒而建圣堂之事。}

此后，他开始在他的同名城市中建造一座纪念宗徒的圣堂。他将这座建筑建得极高，并用各色大理石石板从地基到屋顶加以镶嵌，使之光彩夺目。他还建造了内部精细的格子天花板，并全部镀金。外部覆盖物用黄铜代替瓦片，保护建筑免受雨水侵蚀；同样用黄金华丽地大量装饰，反射出太阳的光芒，令远处观看者目眩。穹顶完全由黄铜和黄金制成的精美雕刻装饰环绕。\par

% structure: p0123 source=h2 target=subsection reason=relative-heading-rank; base=h2

\subsection{第五十九章 同上圣堂的进一步描述。}

皇帝以如此壮丽的方式美化了这座圣堂。建筑被一个广阔的露天区域所环绕，四侧以柱廊为界，柱廊将区域和圣堂本身围合。与这些柱廊相邻的是一排排庄严的厅堂，带有浴室和散步道，此外还有许多适合管理该处人员使用的房间。\par

% structure: p0125 source=h2 target=subsection reason=relative-heading-rank; base=h2

\subsection{第六十章 他也在该圣堂中修建了自己的墓穴纪念碑。}

皇帝奉献所有这些建筑，是渴望永久纪念我们救主的宗徒。然而，他修建这座建筑还有另一个目的，一个起初不为人知，后来却众人皆知的目的。事实上，他选择这个地方是预见自己的死亡，以非凡信德的热情预期他的身体将与宗徒们共享他们的名号，从而在死后与他们一起成为在此地为他们举行敬礼的对象。于是他让人在圣堂内设置十二口石棺，如同神圣的柱子，以纪念和缅怀宗徒的数目，他自己的石棺放在中间，两侧各有六口宗徒的石棺。这样，正如我所说，他以谨慎的远见为自己死后提供了一个光荣的安息之所，并早已秘密作出这一决定，现在他将此圣堂奉献给宗徒，相信这种对宗徒纪念的尊崇将对他的灵魂大有裨益。天主并没有使他失望，没有辜负他如此热切期待和渴望的事。因为在他完成了复活节的首轮礼仪，并以一种令他自己和所有人都喜乐欢欣的方式度过了我们主的这个神圣日子之后，那位协助他完成这一切行动、且他终身热诚事奉的天主，乐意在一个幸福的时刻将他迁往更美好的生命。\par

% structure: p0127 source=h2 target=subsection reason=relative-heading-rank; base=h2

\subsection{第六十一章 他在\Place{海伦波利斯}患病，并为领受圣洗而祈祷。}

起初，他感到身体有些轻微不适，很快便发展为确切的疾病。因此，他前往自己城中的温泉浴场；随后前往以他母亲命名的城市\Place{海伦波利斯}。他在殉道圣堂里度过了一段时间，向天主献上祈祷和恳求。最终，他确信自己的生命即将结束，感到是时候寻求洗去过去一生的罪污了，他坚信，无论他作为凡人犯过什么过错，他的灵魂都将通过圣洗的奥秘之言和救赎之水得到净化。怀着这些想法，他俯伏在圣堂内，向天主献上祈祷和忏悔，在这圣堂里，他生平第一次接受了覆手和祈祷。此后他前往\Place{尼科美底亚}的郊区，在那里召集主教们会面，向他们说了以下的话。\par

% structure: p0129 source=h2 target=subsection reason=relative-heading-rank; base=h2

\subsection{第六十二章 \Person{君士坦丁}恳请主教们为他施行\OriginalTerm{圣洗圣事}{Baptism}。}

\Quote{我所长久盼望的时刻已经来到，我热切渴望并祈祷能获得天主的救恩。这一刻已经来临，我也能领受那赋予不朽的印记的福分；这便是我能领受救恩印记的时刻。我原想在约但河中领受圣洗，因为我们的救主为了我们的榜样，曾在那里受圣洗。然而，深知何事为我们有益的天主，乐意我在此地领受这福分。那么，就如此行吧，不要再延迟了。因为，如果那位掌管生死的主愿意延长我在世的生命，并命我此后与天主子民为伍，作为祂教会的一员与他们一同祈祷，我将从此为自己规定一种合乎祂事奉的生活方式。}\par

% structure: p0131 source=h2 target=subsection reason=relative-heading-rank; base=h2

\subsection{第六十三章 他受圣洗后如何感谢天主}

于是他扬声向天主献上感恩的祈祷，随后又加上了这些话：\Quote{现在我知道我真是有福的；现在我确信我被视为配得永生，并成为天主之光的分享者。}他又对那些无缘享受他所领受的福分的人的不幸处境表示怜悯。当他的护民官与将军们出现在他面前，为即将失去他而哀恸流泪，祈求他的寿数得以延长时，他回答他们说，他现在已拥有真生命；只有他自己才能知道他所得的福分是何等宝贵；因此，他切望早日离世归主，而非延迟。接着，他开始完成必要的事务安排：将其帝国城市的罗马居民遗赠年度捐赠；将帝国遗产如祖业般分给他的子女；总之，按照自己的意愿作了一切处置。\par

% structure: p0133 source=h2 target=subsection reason=relative-heading-rank; base=h2

\subsection{第六十四章 君士坦丁于圣神降临节正午去世}

这一切都发生在一个极其重要的庆节期间，我指的是那庄严神圣的圣神降临节，它以七周为期限，并以那一日——圣经证明我们共同的救主升天及圣神降临人间之日——为印封。在此庆节期间，皇帝领受了我所描述的特恩；而在最后的那一日，人们可恰如其分地称之为庆节中的庆节，他在正午时刻被接到天主的面前，将其必朽的躯体留给他的同胞，而将那能理解并爱慕天主的部分带到与天主的共融中。这就是君士坦丁现世生命的终结。现在让我们来看随后发生的事。\par

% structure: p0135 source=h2 target=subsection reason=relative-heading-rank; base=h2

\subsection{第六十五章 士兵与军官的哀悼}

\Emph{立刻}，聚集的持枪兵士与亲卫队撕裂衣服，俯伏在地，以头撞地，发出哀号与悲痛的哭喊，呼唤他们的皇帝与主上——更确切地说，如同真正的儿女，呼唤他们的父亲；而护民官与百夫长则称他为他们的保护者、护卫者与恩主。其余的士兵也以恭敬的秩序前来哀悼，如同羊群失去他们善良的牧人。与此同时，百姓在城中狂奔，有的以高声呼喊表达内心的悲伤，有的则因悲痛而茫然失措。每个人都视此事件为己身的灾祸，哀悼他的去世，仿佛感到自己被剥夺了人人共享的福分。\par

% structure: p0137 source=h2 target=subsection reason=relative-heading-rank; base=h2

\subsection{第六十六章 遗体从尼科美底亚移往君士坦丁堡皇宫}

此后，士兵们将遗体从卧榻抬起，放入金棺，用紫色覆盖物包裹，然后移往以皇帝名字命名的城市。在那里，遗体被安放在皇宫正殿的高处，周围是用金烛台点燃的蜡烛，呈现出一幅奇妙的景象，自世界初创以来，从未有人见过在日光之下有这样的景象。因为在皇宫的中央大殿中，皇帝的遗体穿着君权的标志——冠冕与紫袍——安卧在高处，四周有众多侍从环绕，日夜不停地看守。\par

% structure: p0139 source=h2 target=subsection reason=relative-heading-rank; base=h2

\subsection{第六十七章 他生前受到的尊荣，伯爵及其他官员在他去世后仍予以奉行}

最高的军事将领、伯爵们以及全体官吏，他们素来习惯于向皇帝敬礼，现在仍然毫不改变地履行这一义务，甚至在他死后，他们按指定时间进入寝殿，屈膝向放在棺材中的皇帝致敬，仿佛他仍活着。之后，元老们和所有曾担任过任何尊贵职务的人出现，行了同样的敬礼。随后是各阶层的群众，带着他们的妻子儿女前来观礼。这些敬意持续了相当长的时间，因为士兵们决定这样守护遗体，直到他的儿子们抵达并亲自办理父亲的葬礼。没有哪个凡人像这位有福的皇帝一样，死后仍然继续统治，并接受与生前相同的敬礼：他是所有活着的人中唯一从天主获得这种赏报的。这是一份合适的赏报，因为他在所有行动中尊崇至高的天主和他的基督，因此天主本人也乐意使他的遗骸在人间仍然保持帝国权威，以此向凡并非全然缺乏理解的人指明，他的灵魂注定要享有的不朽而永无止境的帝国。这就是此处所发生的事。\par

% structure: p0141 source=h2 target=subsection reason=relative-heading-rank; base=h2

\subsection{第六十八章 军队决定此后将\OriginalTerm{奥古斯都}{Augustus}称号授予其子嗣的决议}

\Emph{与此同时}，护民官从他们指挥的部队中选出那些忠诚热忱、长久以来为皇帝所知的军官，派遣他们向两位\OriginalTerm{凯撒}{Cæsars}报告这一近期事件。他们遂执行了这项任务。然而，当各省的士兵得知皇帝去世的消息后，全体士兵仿佛受到超自然冲动的驱使，一致决议，似乎伟大的皇帝仍然活着，不承认任何人除了他的儿子们作为罗马世界的君主；并且他们很快决定这些儿子不应再保留凯撒之名，而应每人获封\OriginalTerm{奥古斯都}{Augustus}称号，这一称号表明皇权的最高至上性。这就是军队所采取的措施；这些决议通过书信相互传达，因此各军团的共同愿望在同一时刻传遍整个帝国。\par

% structure: p0143 source=h2 target=subsection reason=relative-heading-rank; base=h2

\subsection{第六十九章 君士坦丁在罗马的哀悼，以及死后通过绘画向他所表示的敬意}

当皇帝去世的消息抵达帝都时，罗马元老院和人民感到这是最沉重、最痛苦的不幸，沉溺于过度的悲伤。浴场和市场关闭，公共表演以及所有其他有闲之人惯常享受的娱乐都中止了。那些先前过着奢侈安逸生活的人，现在在街上愁容满面地行走，而所有人一致赞美已故者的名字，视他为天主所爱、真正配得皇帝尊严的人。他们的悲伤并不仅限于言语；他们还通过奉献绘画以纪念他，像他生前一样表示尊敬。这些图画的设计描绘了天堂本身，将皇帝描绘为在天穹之上的一座空灵府邸中安息。他们也宣布唯有他的儿子们为皇帝和\OriginalTerm{奥古斯都}{Augusti}，并恳切请求准许他们在帝都接收皇帝遗体并为他举行葬礼。\par

% structure: p0145 source=h2 target=subsection reason=relative-heading-rank; base=h2

\subsection{第七十章 由他的儿子君士坦提乌斯在君士坦丁堡举行的葬礼}

然而，他的次子此时已经抵达，便开始在他所命名的城市中为他的父亲举行葬礼，他亲自领队，由排列成军阵的士兵分队前导，后面跟着大批群众，遗体本身由持矛兵和重装步兵连队环绕。当队伍抵达奉献于我们救主宗徒的教堂时，棺材被安葬在那里。年轻的皇帝君士坦提乌斯通过他的亲临和适当举行这一神圣仪式，这样向他的亡父表示敬意。\par

% structure: p0147 source=h2 target=subsection reason=relative-heading-rank; base=h2

\subsection{第七十一章 在宗徒教堂为君士坦丁葬礼举行的神圣礼仪}

当君士坦提乌斯随同军旅退去后，天主的仆人们上前，与群众和全体信众一起，以祈祷举行神圣敬礼。同时，颂赞的贡品献给这位有福君王的品德，他的遗体安放在一个高大醒目的纪念碑上，全体群众与天主的司祭一同为他的灵魂献上祈祷，不是没有眼泪——不，而是多有哭泣；这样履行了一项与虔诚亡者的愿望相宜的礼仪。在这方面，天主的恩惠也向他的仆人彰显，因为他不仅将帝国的继承权留给他自己亲爱的儿子们，而且他有福的灵魂的尘世居所，按照他自己恳切的愿望，获准分享宗徒的纪念碑，与他们的名号和天主子民的名号相关联，因举行神圣礼仪和奥迹礼仪而受尊荣，并参与圣人们的祈祷。这样，他在死后也继续拥有帝国权力，仿佛以重获新生的生命控制着宇宙的统治权，以\OriginalTerm{胜利者}{Victor}、\OriginalTerm{至大者}{Maximus}、\OriginalTerm{奥古斯都}{Augustus}的称号，在他自己的名下保持罗马世界的统治权。\par

% structure: p0149 source=h2 target=subsection reason=relative-heading-rank; base=h2

\subsection{第七十二章 论凤凰}

我们不能把他比作那埃及鸟，据说它独一无二，在芬芳香料中自献而死，又从自己灰烬中复活，以先前的形貌高飞。他倒更像他的救主，后者如同那播种的种子，从一粒生出许多，因天主的祝福而丰收，用其果实覆盖了整个世界。同样，我们这位三倍有福的君王也借着他诸子的传承而倍增。他的雕像在各省与他们的并列而竖立；君士坦丁的名字甚至在他有生之年结束后仍被承认和尊崇。\par

% structure: p0151 source=h2 target=subsection reason=relative-heading-rank; base=h2

\subsection{第73章 钱币上如何描绘君士坦丁升天}

\Emph{一种钱币}也被铸造，印有下述图案。一面出现我们有福君王的形象，头部紧密蒙着面纱：另一面显示他坐着如同御车者，由四匹马拉着，一只手从天伸下接他升天。\par

% structure: p0153 source=h2 target=subsection reason=relative-heading-rank; base=h2

\subsection{第74章 他所尊崇的天主也当得起地尊崇了他作为回报}

至高天主藉着君士坦丁——在所有统治者中唯独他公开承认基督信仰——向我们显明的这些证据，表明他如何看待那些有特权敬拜他及其基督的人，与那些选择了相反派别、胆敢攻击他的教会而激起他敌意的人之间的巨大差异。后者的悲惨结局在每一事例上都提供了他不悦的标记，正如君士坦丁的死向所有人传递了他神圣之爱的明确保证。\par

% structure: p0155 source=h2 target=subsection reason=relative-heading-rank; base=h2

\subsection{第75章 他在对天主的虔诚上超越了所有先前的皇帝}

他在罗马皇帝中独自、卓越地站立，作为天主的崇拜者；独自作为基督教义向所有人的大胆宣告者；独自向他的教会给予尊荣，前人从未如此；独自彻底消除了多神论的错误，反对一切形式的偶像崇拜：因此，在生前死后，他独独被认为配得上无人能说他人已达成的尊荣；以至于没有人，无论希腊人或蛮族，甚至古罗马人本身，曾被呈现给我们堪与他相比。\par

\SourceRecord{Translated by Ernest Cushing Richardson.；Nicene and Post-Nicene Fathers, Second Series；Vol. 1.；Edited by Philip Schaff and Henry Wace.；Buffalo, NY: Christian Literature Publishing Co.,；1890.；<http://www.newadvent.org/fathers/25024.htm>.}
